% Genetics — Biology Upper Sixth — Cours signé OnBuch+
% Official MINESEC syllabus. Compiler avec : tectonic BIO03-genetics.tex
\newcommand{\DOCMATIERE}{Biology}
\newcommand{\DOCNIVEAU}{Upper Sixth}
\newcommand{\DOCMODULE}{Upper Sixth — Biology}
\newcommand{\DOCLECON}{3}
\newcommand{\DOCTITRE}{Genetics}
% OnBuch+ — préambule « cours signés » ENRICHI (compilé avec tectonic / XeLaTeX)
% Système visuel « Pop » d'OnBuch+ : fond crème (jamais blanc), bordures encre
% épaisses, ombres « dures » décalées, titres Archivo Black en capitales.
% Version MATHÉMATIQUES : graphes (pgfplots/tikz), boîtes pédagogiques
% (définition, propriété, méthode, attention, le savais-tu, exercice résolu,
% exercice, TP, bilan), page de garde et signature OnBuch+ ancrée en bas.
%
% Macros attendues dans main.tex AVANT \input{preamble} :
%   \DOCMATIERE  \DOCNIVEAU  \DOCMODULE  \DOCLECON (numéro)  \DOCTITRE
\documentclass[11pt]{article}

\usepackage{fontspec}
\usepackage[english]{babel}
\usepackage[a4paper,margin=2cm,top=3.4cm,bottom=2.8cm,headheight=1.4cm,footskip=1.3cm]{geometry}
\usepackage{amsmath,amssymb,amsfonts}
\usepackage{mathtools} % \xmapsto, \coloneqq... souvent utilisés par les modèles
\usepackage{cancel} % simplifications barrées (bilans de forces)
% \abs{z} (module, valeur absolue) et \norm{u} (norme), utilisés par les modèles
\providecommand{\abs}{}\let\abs\relax\DeclarePairedDelimiter{\abs}{\lvert}{\rvert}
\providecommand{\norm}{}\let\norm\relax\DeclarePairedDelimiter{\norm}{\lVert}{\rVert}
\usepackage[table]{xcolor} % [table] : fournit aussi \rowcolors (alternance de lignes)
\usepackage{pagecolor}
\usepackage{tikz}
\usetikzlibrary{arrows.meta,positioning,calc,shapes.geometric,shapes.misc,shapes.arrows,decorations.pathmorphing,decorations.pathreplacing,decorations.markings,patterns,shadows,mindmap,trees,fit,backgrounds,matrix,intersections,angles,quotes}
\usepackage{pgfplots}
\usepackage[version=4]{mhchem} % équations nucléaires \ce{^{235}_{92}U}
\usepackage[european,siunitx]{circuitikz} % schémas électriques
\pgfplotsset{compat=1.18}
\usepgfplotslibrary{fillbetween,groupplots} % aires entre courbes (calcul intégral)
\usepackage{enumitem}
\usepackage{fancyhdr}
% [most] charge toutes les bibliothèques tcolorbox (listings, minted, raster,
% documentation...) et gonfle inutilement la mémoire TeX sur un document long
% (~300 boîtes) : on ne charge que ce qui est réellement utilisé.
\usepackage[skins,breakable]{tcolorbox}
\usepackage{graphicx}
\usepackage{colortbl}
\usepackage{booktabs}
\usepackage{tabularx}
\usepackage{diagbox} % case d'en-tête diagonale des tableaux à double entrée
% Colonnes extensibles centrée / gauche / droite (C, L, R), souvent utilisées par les modèles
\newcolumntype{C}{>{\centering\arraybackslash}X}
\newcolumntype{L}{>{\raggedright\arraybackslash}X}
\newcolumntype{R}{>{\raggedleft\arraybackslash}X}
\usepackage{array}
\usepackage{multicol}
\usepackage{siunitx}
% parse-numbers=false : accepte aussi \SI{2,5 \times 10^{-3}}{...}
\sisetup{locale=UK,output-decimal-marker={.},group-minimum-digits=4,parse-numbers=false}
\DeclareSIUnit\ohm{\ensuremath{\symup{\Omega}}} % U+2126 absent de la police
\DeclareSIUnit\division{div} % réglages d'oscilloscope (ms/div, V/div)
\DeclareSIUnit\tr{tr} % tours (vitesse de rotation, ex. \SI{1500}{\tr\per\minute})
\usepackage{qrcode}
% \degree : commande fréquemment utilisée par les modèles pour un angle ou une
% température en dehors de \SI{}{} (non fournie par les paquets chargés).
\providecommand{\degree}{^{\circ}}
% \qed (fin de démonstration), utilisé par les modèles sans amsthm
\providecommand{\qed}{\hfill$\square$}
% \barre : double barre des valeurs interdites dans un tableau de variations (tkz-tab)
\providecommand{\barre}{\ensuremath{\|}}
% environnement proof (amsthm non chargé) utilisé par les modèles
\ifdefined\proof\else\newenvironment{proof}[1][Proof]{\par\noindent\textit{#1.}\ }{\qed\par}\fi
\usepackage{lastpage}
\usepackage{hyperref}
\hypersetup{hidelinks}

% ============================================================
% Palette « Pop » OnBuch+ — mêmes hex que la classe P (plus_ui.dart)
% ============================================================
\definecolor{poporange}{HTML}{F59321}
\definecolor{popdark}{HTML}{E07A0C}
\definecolor{popcream}{HTML}{FFF6E8}
\definecolor{popink}{HTML}{1C1714}
\definecolor{poppurple}{HTML}{7A5AE0}
\definecolor{popgreen}{HTML}{2E9E6A}
\definecolor{poppink}{HTML}{E0457B}
\definecolor{popblue}{HTML}{2F7DE1}
\definecolor{popgold}{HTML}{C98A12}
\definecolor{popmuted}{HTML}{8A7F76}
\definecolor{popline}{HTML}{EADFD2}
\definecolor{popgray}{HTML}{8A7F76}
\colorlet{popgrey}{popgray}
% Alias tolérants (le modèle nomme parfois les couleurs différemment)
\colorlet{popwhite}{white}
\colorlet{popblack}{popink}
\colorlet{popred}{poppink}
\colorlet{popyellow}{popgold}
% Teintes claires pour fonds de tableaux / zones
\colorlet{popblueL}{popblue!10!white}
\colorlet{poporangeL}{poporange!14!white}
\colorlet{popgreenL}{popgreen!12!white}
\colorlet{poppurpleL}{poppurple!12!white}
\colorlet{poppinkL}{poppink!12!white}
\colorlet{popgoldL}{popgold!14!white}

\pagecolor{popcream}

% ============================================================
% Typographie
% ============================================================
\defaultfontfeatures{Ligatures=TeX}
\setmainfont{PlusJakartaSans-Medium.ttf}[Path=../../../fonts/,BoldFont=PlusJakartaSans-Bold.ttf]
% Maths en sans-serif (Fira Math) avec chiffres et lettres droites de Plus
% Jakarta, pour que les formules se fondent dans le texte.
\usepackage[math-style=ISO,bold-style=ISO]{unicode-math}
\setmathfont{FiraMath-Regular.otf}[Path=../../../fonts/]
\setmathfont{PlusJakartaSans-Medium.ttf}[Path=../../../fonts/,range={up/num,up/{Latin,latin},\mathup}]
% (icomma retiré : décimales avec point en anglais)
% Caractères mathématiques Unicode absents des polices de texte (Plus Jakarta,
% Archivo Black) : rendus par leur équivalent mathématique, en texte comme en
% formule — sinon ils s'affichent en carré vide (ex. « Arithmétique dans ℤ »).
\usepackage{newunicodechar}
\newunicodechar{ℝ}{\ensuremath{\mathbb{R}}}
\newunicodechar{ℤ}{\ensuremath{\mathbb{Z}}}
\newunicodechar{ℂ}{\ensuremath{\mathbb{C}}}
\newunicodechar{ℕ}{\ensuremath{\mathbb{N}}}
\newunicodechar{ℚ}{\ensuremath{\mathbb{Q}}}
\newunicodechar{𝔻}{\ensuremath{\mathbb{D}}}
\newunicodechar{α}{\ensuremath{\alpha}}
\newunicodechar{β}{\ensuremath{\beta}}
\newunicodechar{γ}{\ensuremath{\gamma}}
\newunicodechar{δ}{\ensuremath{\delta}}
\newunicodechar{ε}{\ensuremath{\varepsilon}}
\newunicodechar{θ}{\ensuremath{\theta}}
\newunicodechar{λ}{\ensuremath{\lambda}}
\newunicodechar{μ}{\ensuremath{\mu}}
\newunicodechar{σ}{\ensuremath{\sigma}}
\newunicodechar{τ}{\ensuremath{\tau}}
\newunicodechar{φ}{\ensuremath{\varphi}}
\newunicodechar{ψ}{\ensuremath{\psi}}
\newunicodechar{ω}{\ensuremath{\omega}}
\newunicodechar{Σ}{\ensuremath{\Sigma}}
% majuscules produites par \MakeUppercase dans les titres (α-aminés -> Α)
\newunicodechar{Α}{\ensuremath{\alpha}}
\newunicodechar{Β}{\ensuremath{\beta}}
\newunicodechar{Γ}{\ensuremath{\Gamma}}
\newunicodechar{Δ}{\ensuremath{\Delta}}
\newunicodechar{Ε}{\ensuremath{\varepsilon}}
\newunicodechar{Θ}{\ensuremath{\Theta}}
\newunicodechar{Λ}{\ensuremath{\Lambda}}
\newunicodechar{Μ}{\ensuremath{\mu}}
\newunicodechar{Τ}{\ensuremath{\tau}}
\newunicodechar{Φ}{\ensuremath{\Phi}}
\newunicodechar{Ψ}{\ensuremath{\Psi}}
\newunicodechar{Ω}{\ensuremath{\Omega}}
\newunicodechar{↦}{\ensuremath{\mapsto}}
\newunicodechar{∈}{\ensuremath{\in}}
\newunicodechar{∉}{\ensuremath{\notin}}
\newunicodechar{∧}{\ensuremath{\wedge}}
\newunicodechar{⇔}{\ensuremath{\Leftrightarrow}}
\newunicodechar{⟺}{\ensuremath{\Longleftrightarrow}}
\newunicodechar{⇒}{\ensuremath{\Rightarrow}}
\newunicodechar{⟹}{\ensuremath{\Longrightarrow}}
\newunicodechar{∘}{\ensuremath{\circ}}
\newunicodechar{∪}{\ensuremath{\cup}}
\newunicodechar{∩}{\ensuremath{\cap}}
\newunicodechar{≡}{\ensuremath{\equiv}}
\newunicodechar{⊂}{\ensuremath{\subset}}
\newunicodechar{⊥}{\ensuremath{\perp}}
\newunicodechar{∃}{\ensuremath{\exists}}
\newunicodechar{∀}{\ensuremath{\forall}}
\newunicodechar{∛}{\ensuremath{\sqrt[3]{\,}}}
\newunicodechar{∜}{\ensuremath{\sqrt[4]{\,}}}
\newunicodechar{⃗}{\ensuremath{{}^{\to}}}
\newunicodechar{ⁿ}{\ifmmode^{n}\else\textsuperscript{\ensuremath{n}}\fi}
\newunicodechar{ˣ}{\ifmmode^{x}\else\textsuperscript{\ensuremath{x}}\fi}
\newunicodechar{ᵇ}{\ifmmode^{b}\else\textsuperscript{\ensuremath{b}}\fi}
\newunicodechar{ʳ}{\ifmmode^{r}\else\textsuperscript{\ensuremath{r}}\fi}
\newunicodechar{ᵘ}{\ifmmode^{u}\else\textsuperscript{\ensuremath{u}}\fi}
\newunicodechar{ʸ}{\ifmmode^{y}\else\textsuperscript{\ensuremath{y}}\fi}
\newunicodechar{ᵃ}{\ifmmode^{a}\else\textsuperscript{\ensuremath{a}}\fi}
\newunicodechar{ᵉ}{\ifmmode^{e}\else\textsuperscript{\ensuremath{e}}\fi}
\newunicodechar{ᵏ}{\ifmmode^{k}\else\textsuperscript{\ensuremath{k}}\fi}
\newunicodechar{ᵖ}{\ifmmode^{p}\else\textsuperscript{\ensuremath{p}}\fi}
\newunicodechar{ᴺ}{\ifmmode^{N}\else\textsuperscript{\ensuremath{N}}\fi}
\newunicodechar{ᶜ}{\ifmmode^{c}\else\textsuperscript{\ensuremath{c}}\fi}
\newunicodechar{ʰ}{\ifmmode^{h}\else\textsuperscript{\ensuremath{h}}\fi}
\newunicodechar{ᵠ}{\ifmmode^{q}\else\textsuperscript{\ensuremath{q}}\fi}
\newunicodechar{ₙ}{\ifmmode_{n}\else\textsubscript{\ensuremath{n}}\fi}
\newunicodechar{ₐ}{\ifmmode_{a}\else\textsubscript{\ensuremath{a}}\fi}
\newunicodechar{ᵢ}{\ifmmode_{i}\else\textsubscript{\ensuremath{i}}\fi}
\newunicodechar{ᵣ}{\ifmmode_{r}\else\textsubscript{\ensuremath{r}}\fi}
\newunicodechar{ₖ}{\ifmmode_{k}\else\textsubscript{\ensuremath{k}}\fi}
\newunicodechar{ᵦ}{\ifmmode_{\beta}\else\textsubscript{\ensuremath{\beta}}\fi}
\newunicodechar{ₓ}{\ifmmode_{x}\else\textsubscript{\ensuremath{x}}\fi}
\newfontfamily\popdisplay{ArchivoBlack-Regular.ttf}[Path=../../../fonts/]
\newfontfamily\poplogo{SpaceGrotesk-Bold.ttf}[Path=../../../fonts/]
\newfontfamily\popsemi{PlusJakartaSans-SemiBold.ttf}[Path=../../../fonts/]
\newfontfamily\popxbold{PlusJakartaSans-ExtraBold.ttf}[Path=../../../fonts/]
\newfontfamily\popxboldls{PlusJakartaSans-ExtraBold.ttf}[Path=../../../fonts/,LetterSpace=3.0]

\setlist[enumerate]{leftmargin=1.7em,itemsep=5pt,topsep=4pt}
\setlist[itemize]{leftmargin=1.7em,itemsep=4pt,topsep=4pt,label={\color{poporange}\textbullet}}
\setlength{\parindent}{0pt}
\setlength{\parskip}{5pt}
\renewcommand{\arraystretch}{1.25}

% pgfplots : style Pop par défaut
\pgfplotsset{
  every axis/.append style={axis line style={popink,line width=1pt},tick style={popink},
    grid style={popline,line width=0.5pt},label style={font=\small},tick label style={font=\footnotesize}},
  popcurve/.style={poporange,line width=1.8pt,no markers,smooth},
}
% Même style disponible dans un \draw TikZ simple (hors pgfplots)
\tikzset{popcurve/.style={poporange,line width=1.8pt}}
\tikzset{popgrid/.style={popline,very thin}}
\colorlet{popcurve}{poporange} % le nom du style est parfois utilisé comme couleur

% ============================================================
% Pill / squiggle
% ============================================================
\newcommand{\poppill}[3]{%
  \tikz[baseline=-0.55ex]\node[draw=popink,line width=1.2pt,rounded corners=2.6pt,%
    fill=#2,inner xsep=6.5pt,inner ysep=3pt]{%
    \popxboldls\fontsize{7}{7}\selectfont\color{#3}\MakeUppercase{#1}};%
}
\newcommand{\popsquiggle}[1]{%
  \tikz[baseline=-0.4ex]{
    \draw[#1,line width=2.6pt,line cap=round]
      plot [smooth] coordinates {(0,0.09) (0.34,0.01) (0.68,0.21) (1.02,0.01) (1.36,0.10)};
  }%
}
% Mot-clé surligné (marqueur orange sous le texte)
\newcommand{\cle}[1]{{\popxbold\color{popdark}#1}}

% ============================================================
% En-tête / pied
% ============================================================
\pagestyle{fancy}
\fancyhf{}
\renewcommand{\headrulewidth}{0pt}
\renewcommand{\footrulewidth}{0pt}
\lhead{\raisebox{-0.15ex}{\poplogo\fontsize{13}{13}\selectfont\color{popink}OnBuch\color{poporange}+}}
\rhead{\poppill{\DOCMATIERE\ \textbullet\ \DOCNIVEAU\ \textbullet\ Chapter \DOCLECON}{white}{popink}}
\lfoot{\popxbold\fontsize{7.6}{7.6}\selectfont\color{popmuted}\MakeUppercase{Signed course OnBuch+}\ \textbullet\ {\popsemi\DOCTITRE}}
\rfoot{\tikz[baseline=-0.6ex]\node[rounded corners=8.5pt,fill=popink,minimum height=17pt,minimum width=17pt,inner xsep=5pt,inner sep=0]{\popxbold\fontsize{8}{8}\selectfont\color{popcream}\thepage\,/\,\pageref*{LastPage}};}

% ============================================================
% Titres
% ============================================================
\newcommand{\chaptitre}[2]{%
  \begin{tcolorbox}[enhanced,colback=poporange,colframe=popink,boxrule=2.6pt,arc=11pt,
      drop shadow=popink,left=15pt,right=15pt,top=12pt,bottom=11pt,before skip=3pt,after skip=18pt]
    \poppill{#2}{popink}{popcream}\par\vspace{9pt}
    {\raggedright\hyphenpenalty=10000\exhyphenpenalty=10000\popdisplay\fontsize{22}{24}\selectfont\color{popcream}\MakeUppercase{#1}\par}
  \end{tcolorbox}
}
% Section numérotée : pastille orange avec le numéro + titre Archivo
\newcounter{popsec}
\newcommand{\coursec}[1]{%
  \stepcounter{popsec}\setcounter{popsub}{0}%
  \par\vspace{16pt}\noindent
  \begin{minipage}{\linewidth}
  \tikz[baseline=-0.7ex]\node[draw=popink,line width=1.6pt,fill=poporange,rounded corners=4pt,minimum size=22pt,inner sep=0]{\popdisplay\fontsize{12}{12}\selectfont\color{popcream}\thepopsec};%
  \hspace{8pt}{\popdisplay\fontsize{15}{17}\selectfont\color{popink}\MakeUppercase{#1}}\par
  \vspace{4pt}\noindent{\color{poporange}\rule{36pt}{3.6pt}}
  \end{minipage}\par\nopagebreak\vspace{8pt}
}
\newcounter{popsub}
\newcommand{\courssub}[1]{%
  \stepcounter{popsub}%
  \par\vspace{9pt}\noindent{\popxbold\fontsize{11.5}{13}\selectfont\color{popink}{\color{poporange}\thepopsec.\thepopsub}\hspace{6pt}#1}\par\nopagebreak\vspace{4pt}
}

% ============================================================
% Boîtes pédagogiques — même recette : fond blanc, bordure épaisse colorée,
% ombre dure encre, badge plein. Titre optionnel [#1].
% ============================================================
\tcbset{popbase/.style={breakable,enhanced,colback=white,boxrule=2.3pt,arc=9pt,
  shadow={3pt}{-3pt}{0pt}{popink},left=14pt,right=14pt,top=11pt,bottom=11pt,before skip=11pt,after skip=15pt,
  fontupper=\popsemi\fontsize{10.6}{14.5}\selectfont\color{popink},
  fontlower=\popsemi\fontsize{10.6}{14.5}\selectfont\color{popink}}}
% Coche dessinée (le glyphe ✓ n'existe pas dans les polices du document)
\newcommand{\popcheck}[1]{\tikz[baseline=-0.5ex]\draw[#1,line width=1.6pt,line cap=round,line join=round](-0.13,0.0)--(-0.04,-0.09)--(0.13,0.1);}
\providecommand{\coche}{\popcheck{popgreen}}
\providecommand{\resultat}[1]{\par\smallskip\noindent\fcolorbox{popgreen}{popgreenL}{\parbox{\dimexpr\linewidth-2\fboxsep-2\fboxrule\relax}{\textbf{Result.} #1}}\par\smallskip}
\newcommand{\popboxhead}[3]{\poppill{#1}{#2}{white}\ifx\relax#3\relax\else\hspace{6pt}{\popxbold\fontsize{10.6}{12}\selectfont\color{popink}#3}\fi\par\vspace{7pt}}

\newtcolorbox{aretenir}[1][]{popbase,colframe=popgreen,before upper={\popboxhead{Key points}{popgreen}{#1}}}
\newtcolorbox{exemplebox}[1][]{popbase,colframe=poporange,before upper={\popboxhead{Exemple}{poporange}{#1}}}
\newtcolorbox{definition}[1][]{popbase,colframe=popblue,before upper={\popboxhead{Definition}{popblue}{#1}}}
\newtcolorbox{propriete}[1][]{popbase,colframe=poppurple,before upper={\popboxhead{Property}{poppurple}{#1}}}
\newtcolorbox{methode}[1][]{popbase,colframe=popgold,before upper={\popboxhead{Method}{popgold}{#1}}}
\newtcolorbox{attention}[1][]{popbase,colframe=poppink,before upper={\popboxhead{Warning}{poppink}{#1}}}
\newtcolorbox{experience}[1][]{popbase,colframe=popblue,colback=popblueL,before upper={\popboxhead{Experiment}{popblue}{#1}}}
% « Le savais-tu ? » : carte violette pleine (contexte, culture, Cameroun)
\newtcolorbox{savaistu}[1][]{popbase,colback=poppurple,colframe=popink,
  fontupper=\popsemi\fontsize{10.4}{14.2}\selectfont\color{white},
  before upper={\poppill{Did you know?}{white}{poppurple}\ifx\relax#1\relax\else\hspace{6pt}{\popxbold\color{white}#1}\fi\par\vspace{7pt}}}
% Exercice résolu : énoncé en haut, \tcblower puis la solution sur fond crème
\newcounter{popexo}\newcounter{popexr}
\newtcolorbox{exoresolu}[1][]{popbase,colframe=poporange,colbacklower=popcream,
  segmentation style={popink,line width=1.2pt,dashed},
  before upper={\stepcounter{popexr}\popboxhead{Worked example \thepopexr}{poporange}{#1}},
  before lower={\poppill{Solution}{popink}{popcream}\par\vspace{6pt}}}
% Exercice (non résolu en place) — la correction est dans la partie Corrigés
\newtcolorbox{exercice}[1][]{popbase,colframe=popink,
  before upper={\stepcounter{popexo}\popboxhead{Exercise \thepopexo}{popink}{#1}}}
% Corrigé
\newtcolorbox{corrige}[1][]{popbase,colframe=popgreen,colback=popgreenL,
  before upper={\popboxhead{Solution}{popgreen}{#1}}}
% Objectifs / prérequis / bilan
\newtcolorbox{objectifs}{popbase,colframe=popink,colback=poporangeL,
  before upper={\popboxhead{Chapter objectives}{popink}{}}}
\newtcolorbox{prerequis}{popbase,colframe=popink,colback=popblueL,
  before upper={\popboxhead{Prerequisites}{popblue}{}}}
\newtcolorbox{bilan}{popbase,colframe=popink,colback=white,boxrule=2.8pt,
  before upper={\popboxhead{Fiche bilan}{poporange}{L'essentiel à connaître pour l'examen}}}
% Figure encadrée avec légende
\newtcolorbox{popfigure}[1][]{enhanced,breakable=false,colback=white,colframe=popline,boxrule=1.4pt,arc=8pt,
  left=10pt,right=10pt,top=10pt,bottom=8pt,before skip=10pt,after skip=12pt,halign=center,
  lower separated=false,halign lower=center,
  fontlower=\popsemi\fontsize{9}{11.5}\selectfont\color{popmuted},
  #1}
\newcommand{\legende}[1]{\par\vspace{6pt}{\popsemi\fontsize{9}{11.5}\selectfont\color{popmuted}#1}}

% ============================================================
% Signature OnBuch+
% ============================================================
\newcommand{\poplogoinline}[1]{{\poplogo\fontsize{#1}{#1}\selectfont\color{popink}OnBuch\color{poporange}+}}
\newcommand{\signatureonbuch}{%
  \begin{tcolorbox}[enhanced,colback=popink,colframe=popink,boxrule=2.2pt,arc=10pt,
      drop shadow=poporange,left=14pt,right=14pt,top=10pt,bottom=10pt,before skip=0pt,after skip=0pt]
    \tikz[baseline=-0.7ex]\node[circle,fill=popgreen,draw=popcream,line width=1.3pt,minimum size=22pt,inner sep=0]{\popcheck{white}};%
    \hspace{9pt}\begin{minipage}[c]{0.62\linewidth}
      {\popxbold\fontsize{12}{13}\selectfont\color{popcream}Signed\ \textbullet\ {\poplogo OnBuch\color{poporange}+}}\par\vspace{2pt}
      {\raggedright\popsemi\fontsize{8}{10.5}\selectfont\color{popline}\MakeUppercase{OnBuch+ teaching team}\\\MakeUppercase{Follows the official MINESEC syllabus}\par}
    \end{minipage}\hfill
    {\poplogo\fontsize{22}{22}\selectfont\color{popcream}OnBuch\color{poporange}+}
  \end{tcolorbox}%
}

% ============================================================
% Page de garde
% #1 titre · #2 matière · #3 niveau · #4 module · #5 n° leçon · #6 durée ·
% #7 description / objectifs (texte ou itemize)
% ============================================================
\newcommand{\pagegarde}[7]{%
  \thispagestyle{empty}
  \newgeometry{a4paper,margin=1.7cm,top=1.6cm,bottom=1.4cm}%
  \begin{tikzpicture}[remember picture,overlay]
    \fill[poporange!18] ([xshift=-3.2cm,yshift=-3.6cm]current page.north east) circle (4.6cm);
    \fill[poppurple!14] ([xshift=2.4cm,yshift=7.6cm]current page.south west) circle (3.8cm);
  \end{tikzpicture}%
  {\poplogo\fontsize{30}{30}\selectfont\color{popink}OnBuch\color{poporange}+}\hfill
  \raisebox{0.6ex}{\poppill{Signed course \textbullet\ Premium}{popink}{popcream}}\par
  \vspace{1pt}\popsquiggle{poppurple}\par
  \vspace{8pt}
  \begin{tcolorbox}[enhanced,colback=poporange,colframe=popink,boxrule=2.8pt,arc=13pt,
      drop shadow=popink,left=18pt,right=18pt,top=12pt,bottom=13pt,after skip=10pt]
    \poppill{#2 \textbullet\ #3}{popink}{popcream}\hspace{5pt}\poppill{#4}{white}{popink}\par\vspace{8pt}
    {\popdisplay\fontsize{14}{14}\selectfont\color{popink}CHAPTER #5}\par\vspace{4pt}
    {\raggedright\hyphenpenalty=10000\exhyphenpenalty=10000\popdisplay\fontsize{25}{27}\selectfont\color{popcream}\MakeUppercase{#1}\par}
    \vspace{8pt}
    {\popxbold\fontsize{9.5}{9.5}\selectfont\color{popink}\MakeUppercase{Suggested time: #6}}
  \end{tcolorbox}
  \begin{tcolorbox}[enhanced,colback=white,colframe=popink,boxrule=2.2pt,arc=10pt,
      drop shadow=popink,left=16pt,right=16pt,top=10pt,bottom=10pt,after skip=8pt]
    \poppill{In this chapter}{poppurple}{white}\par\vspace{5pt}
    \popsemi\fontsize{9.3}{12.6}\selectfont\color{popink}#7
  \end{tcolorbox}
  \noindent\begin{minipage}[t]{0.487\linewidth}
    \begin{tcolorbox}[enhanced,colback=popgreen,colframe=popink,boxrule=2.2pt,arc=10pt,
        drop shadow=popink,left=11pt,right=11pt,top=10pt,bottom=10pt,height=3.1cm,valign=center]
      \tcbox[colback=white,colframe=popink,boxrule=1.3pt,arc=5pt,boxsep=0pt,nobeforeafter,
        left=3pt,right=3pt,top=3pt,bottom=3pt,valign=center]{\qrcode[height=1.9cm]{https://play.google.com/store/apps/details?id=com.luvvix.onbuch&hl=en}}%
      \hspace{8pt}\begin{minipage}[c]{0.5\linewidth}
        {\popdisplay\fontsize{11}{12.5}\selectfont\color{white}\MakeUppercase{Download OnBuch}}\par\vspace{4pt}
        {\raggedright\popsemi\fontsize{8.4}{11}\selectfont\color{white}Free \textbullet\ Google Play\\AI tutor, past papers, results\par}
      \end{minipage}
    \end{tcolorbox}
  \end{minipage}\hfill
  \begin{minipage}[t]{0.487\linewidth}
    \begin{tcolorbox}[enhanced,colback=poppurple,colframe=popink,boxrule=2.2pt,arc=10pt,
        drop shadow=popink,left=11pt,right=11pt,top=10pt,bottom=10pt,height=3.1cm,valign=center]
      \tikz[baseline=-0.7ex]\node[circle,fill=white,draw=popink,line width=1.3pt,minimum size=1.6cm,inner sep=0]{\popdisplay\fontsize{24}{24}\selectfont\color{poppurple}+};%
      \hspace{8pt}\begin{minipage}[c]{0.6\linewidth}
        {\popdisplay\fontsize{11}{12.5}\selectfont\color{white}\MakeUppercase{Go OnBuch+}}\par\vspace{4pt}
        {\raggedright\popsemi\fontsize{8.4}{11}\selectfont\color{white}All signed courses, unlimited AI tutor, PDF sheets — OnBuch+ tab in the app\par}
      \end{minipage}
    \end{tcolorbox}
  \end{minipage}
  \vspace{8pt}
  \popcontenu
  \vfill
  \signatureonbuch
  \restoregeometry
  \newpage
}

% Rangée « Ce cours contient » (page de garde)
\newcommand{\poptile}[3]{% #1 couleur #2 gros texte #3 légende
  \begin{tcolorbox}[enhanced,colback=#1,colframe=popink,boxrule=2pt,arc=9pt,shadow={2.5pt}{-2.5pt}{0pt}{popink},
      left=6pt,right=6pt,top=9pt,bottom=9pt,halign=center,valign=center,height=1.75cm,width=\linewidth,nobeforeafter]
    {\popdisplay\fontsize{12.5}{13}\selectfont\color{white}\MakeUppercase{#2}}\par\vspace{3pt}
    {\centering\popsemi\fontsize{7.8}{9.5}\selectfont\color{white}#3\par}
  \end{tcolorbox}}
\newcommand{\popcontenu}{%
  \par\noindent{\popxboldls\fontsize{7.5}{7.5}\selectfont\color{popmuted}THIS SIGNED COURSE CONTAINS}\par\vspace{7pt}
  \noindent\begin{minipage}[t]{0.235\linewidth}\poptile{popblue}{Course}{complete, illustrated, step by step}\end{minipage}\hfill
  \begin{minipage}[t]{0.235\linewidth}\poptile{poporange}{Methods}{and worked examples}\end{minipage}\hfill
  \begin{minipage}[t]{0.235\linewidth}\poptile{poppink}{Exercises}{exam style + solutions}\end{minipage}\hfill
  \begin{minipage}[t]{0.235\linewidth}\poptile{popgreen}{Summary sheet}{the essentials to remember}\end{minipage}\par}

% Sommaire
\newtcolorbox{sommaire}{enhanced,colback=white,colframe=popink,boxrule=2.2pt,arc=10pt,
  drop shadow=popink,left=18pt,right=18pt,top=13pt,bottom=13pt,before skip=6pt,after skip=20pt,
  before upper={\poppill{Contents}{poppurple}{white}\par\vspace{9pt}\popsemi\fontsize{10.6}{14.5}\selectfont\color{popink}}}

% Signature de fin de cours — poussée tout en bas de la dernière page
\newcommand{\findecours}{%
  \par\vfill\nopagebreak
  \begin{center}\popsquiggle{poporange}\end{center}\vspace{4pt}
  \signatureonbuch
}
\AtBeginDocument{\def\llbracket{\mathopen{[\mkern-3.5mu[}}\def\rrbracket{\mathclose{]\mkern-3.5mu]}}} % intervalles d'entiers (glyphe ⟦ absent de la police)
% Glyphes absents de Fira Math : redéfinis à partir de glyphes disponibles
\AtBeginDocument{%
  \def\setminus{\mathbin{\backslash}}\let\smallsetminus\setminus
  \def\vdots{\mathord{\vbox{\baselineskip3.2pt\lineskiplimit0pt\kern4pt\hbox{.}\hbox{.}\hbox{.}}}}%
  \def\ddots{\mathinner{\mkern1mu\raise6.5pt\hbox{.}\mkern2mu\raise3.6pt\hbox{.}\mkern2mu\raise0.7pt\hbox{.}\mkern1mu}}%
  \def\triangle{\mathord{\scalebox{0.9}{$\symup{\Delta}$}}}%
  \def\nabla{\mathord{\raisebox{\depth}{\scalebox{1}[-1]{$\symup{\Delta}$}}}}%
  \def\checkmark{\popcheck{popdark}}%
  \def\star{\mathbin{\ast}}\def\bigstar{\mathord{\ast}}%
  \def\diamond{\mathbin{\scalebox{0.6}{$\Diamond$}}}%
  \def\bigcup{\mathop{\mathchoice{\raisebox{-0.3em}{\scalebox{1.7}{$\cup$}}}{\scalebox{1.25}{$\cup$}}{\cup}{\cup}}\displaylimits}%
  \def\bigcap{\mathop{\mathchoice{\raisebox{-0.3em}{\scalebox{1.7}{$\cap$}}}{\scalebox{1.25}{$\cap$}}{\cap}{\cap}}\displaylimits}%
  \def\lesseqqgtr{\mathrel{\lessgtr}}\def\bigoplus{\mathop{\oplus}\displaylimits}%
}
\providecommand{\ssi}{if and only if} % abréviation fréquente
\AtBeginDocument{\providecommand{\wideparen}{\overparen}} % arc de cercle

%% FIGFIX-BEGIN v5 — correctifs globaux des figures (docs/onbuch-generation/tools/figfix)
\usepackage{adjustbox}\usepackage{etoolbox}\tcbuselibrary{raster}
% toute figure TikZ/pgfplots plus large que la ligne est réduite à la largeur disponible
\BeforeBeginEnvironment{tikzpicture}{\begin{adjustbox}{max width=\linewidth}}
\AfterEndEnvironment{tikzpicture}{\end{adjustbox}}
% légendes pgfplots lisibles par-dessus les courbes
\pgfplotsset{every axis/.append style={legend style={fill=white,fill opacity=0.92,text opacity=1,draw=black!15,inner sep=3pt}}}
% étiquettes d'arêtes (syntaxe edge["..."]) avec fond blanc
\tikzset{every edge quotes/.append style={fill=white,inner sep=1.2pt}}
%% FIGFIX-END
\begin{document}

\pagegarde{Genetics}{Biology}{Upper Sixth}{Upper Sixth — Biology}{3}{15 periods}{This chapter explores the principles of heredity, from Mendel's classical experiments to modern population genetics. Students learn to analyse monohybrid and dihybrid crosses, sex linkage, epistasis, multiple alleles, Hardy-Weinberg equilibrium, variation and pedigree analysis — core skills for the GCE Advanced Level Biology examination.
\vspace{4pt}
\begin{multicols}{2}\small
\begin{itemize}
\item By the end of this chapter I can define and correctly use key genetic terms (gene, allele, genotype, phenotype, homozygous, heterozygous, dominant, recessive).
\item By the end of this chapter I can state and apply Mendel's laws of segregation and independent assortment to solve monohybrid and dihybrid crosses.
\item By the end of this chapter I can distinguish complete dominance, codominance and incomplete dominance and predict their phenotypic ratios.
\item By the end of this chapter I can explain sex determination mechanisms and solve problems involving sex-linked traits such as haemophilia and colour blindness.
\item By the end of this chapter I can analyse gene linkage, crossing over and calculate crossover (recombination) values.
\item By the end of this chapter I can interpret epistatic interactions, lethal genes and multiple allele systems such as the ABO blood groups.
\end{itemize}\end{multicols}}

\begin{sommaire}
\begin{multicols}{2}
\begin{enumerate}
\item Foundations: Genetic Terms and Mendel's Laws
\item Patterns of Dominance: Complete, Codominance and Incomplete Dominance
\item Sex Determination and Sex Linkage
\item Dihybrid Inheritance, Linkage and Crossing Over
\item Gene Interactions: Epistasis, Lethal Genes and Multiple Alleles
\item Population Genetics: Gene Frequencies and the Hardy-Weinberg Equation
\item Variation, Polymorphism and Pedigree Analysis
\item Integration activity
\item Methods and worked examples
\item Exercises
\item Solutions to the exercises
\item Summary sheet
\end{enumerate}
\end{multicols}
\end{sommaire}

\begin{objectifs}
\begin{itemize}
\item By the end of this chapter I can define and correctly use key genetic terms (gene, allele, genotype, phenotype, homozygous, heterozygous, dominant, recessive).
\item By the end of this chapter I can state and apply Mendel's laws of segregation and independent assortment to solve monohybrid and dihybrid crosses.
\item By the end of this chapter I can distinguish complete dominance, codominance and incomplete dominance and predict their phenotypic ratios.
\item By the end of this chapter I can explain sex determination mechanisms and solve problems involving sex-linked traits such as haemophilia and colour blindness.
\item By the end of this chapter I can analyse gene linkage, crossing over and calculate crossover (recombination) values.
\item By the end of this chapter I can interpret epistatic interactions, lethal genes and multiple allele systems such as the ABO blood groups.
\item By the end of this chapter I can use the Hardy-Weinberg equation to calculate gene and genotype frequencies in populations.
\item By the end of this chapter I can describe types and causes of variation and explain polymorphism with examples.
\item By the end of this chapter I can construct and interpret pedigree charts to deduce modes of inheritance.
\end{itemize}
\end{objectifs}

\begin{prerequis}
\begin{itemize}
\item Cell structure and the role of chromosomes, genes and DNA (Lower Sixth: cell biology).
\item Mitosis and meiosis, including gamete formation and chromosome behaviour (Lower Sixth).
\item Sexual reproduction in flowering plants and animals (Lower Sixth / Form Five).
\item Basic probability and use of ratios and fractions (Ordinary Level mathematics).
\item The concept of species and populations (Lower Sixth: ecology/diversity).
\end{itemize}
\end{prerequis}

\newpage

\coursec{Foundations: Genetic Terms and Mendel's Laws}

In a village near Bafoussam, a farmer wonders why his red-flowered hibiscus, crossed with a white variety, produces only pink offspring --- while his neighbour's maize shows both purple and yellow kernels on the same cob. At the Bamenda Regional Hospital, a couple is puzzled: both parents have normal blood clotting, yet their son has haemophilia. Meanwhile, a blood bank in Yaound\'e must explain how a child can be blood group O when both parents are group AB\ldots or can they? These everyday puzzles of inheritance are exactly what genetics explains. This chapter gives you the tools to \cle{predict}, \cle{calculate} and \cle{interpret} how traits pass from one generation to the next.

The central problem is simple to state but took humanity thousands of years to solve: \emph{what are the rules that govern the transmission of characters from parents to offspring?} Before 1866, people believed in ``blending inheritance'' --- that parental characters mix like paints. One Austrian monk, working quietly in a monastery garden, proved this wrong using nothing more than pea plants, patience and careful counting.

\courssub{Genetics, Heredity and Variation}

\begin{definition}[Genetics, heredity and variation]
\cle{Genetics} is the branch of biology that studies \cle{heredity} (the transmission of characters from parents to offspring) and \cle{variation} (the differences observed between individuals of the same species).
\end{definition}

Heredity explains why a mango seed always grows into a mango tree and why children resemble their parents. Variation explains why no two children of the same couple (except identical twins) are exactly alike. Genetics is the science that reconciles these two apparently opposite observations: characters are transmitted faithfully, \emph{yet} offspring are never exact copies of their parents.

\begin{savaistu}[Gregor Mendel and the garden pea]
Gregor Mendel (1822--1884), an Augustinian monk at Brno (now in the Czech Republic), carried out his crossing experiments between 1856 and 1863 and published his results in 1866. His work was ignored for over 30 years and was only rediscovered around 1900, after his death. He is now honoured as the \emph{father of genetics}. Remarkably, he deduced the behaviour of ``factors'' (now called genes) without ever seeing a chromosome.
\end{savaistu}

\textbf{Why did Mendel choose \emph{Pisum sativum}, the garden pea?}\par
Mendel's success was not luck: the pea plant was an almost perfect experimental organism.

\begin{itemize}
\item \textbf{Short life cycle:} several generations can be obtained in a few years, so results come quickly.
\item \textbf{Many offspring per cross:} each pod contains several seeds, giving large numbers that make ratios statistically reliable.
\item \textbf{Clear contrasting characters:} peas show easily distinguished pairs of traits (tall vs dwarf, round vs wrinkled seeds, yellow vs green cotyledons) with no intermediates.
\item \textbf{Controlled reproduction:} the pea flower is normally self-pollinating, but it is easy to remove the stamens (emasculation) and dust on pollen from a chosen plant, so the experimenter controls exactly which parents mate.
\item \textbf{Pure-breeding lines available:} varieties that always produce the same trait when self-pollinated could be bought from seed merchants.
\end{itemize}

Mendel also introduced two methodological innovations that made his work scientific: he studied \emph{one character at a time}, and he \emph{counted} his offspring and expressed the results as numerical ratios.

\courssub{Essential Genetic Vocabulary}

Every genetics problem at GCE Advanced Level is solved with the same small set of words. Learn them precisely --- examiners penalise vague usage.

\begin{definition}[Gene, allele, genotype and phenotype]
A \cle{gene} is a segment of DNA occupying a fixed position, the \cle{locus}, on a chromosome, which determines a particular character. An \cle{allele} is one of the alternative forms of a gene found at the same locus (e.g.\ the allele for red seed coat and the allele for white seed coat). The \cle{genotype} is the genetic constitution of an individual --- the combination of alleles it carries. The \cle{phenotype} is the observable expression of the genotype (what the organism looks like), which may also be influenced by the environment.
\end{definition}

\begin{exemplebox}[A Cameroonian example: seed coat colour in local beans]
Consider seed coat colour in beans (\emph{Phaseolus vulgaris}) grown around Dschang. Suppose the gene for seed coat colour has two alleles: $R$ (red coat) and $r$ (white coat). A bean plant of genotype $RR$ or $Rr$ has red seeds (phenotype: red); only genotype $rr$ gives white seeds. Two plants can therefore share the same \emph{phenotype} (red seeds) while having different \emph{genotypes} ($RR$ or $Rr$) --- a distinction that only a cross can reveal.
\end{exemplebox}

\begin{definition}[Homozygous, heterozygous, dominant and recessive]
An individual is \cle{homozygous} for a gene when it carries two identical alleles at that locus ($RR$ or $rr$); it is \cle{heterozygous} when it carries two different alleles ($Rr$). A \cle{dominant} allele is expressed in the phenotype of a heterozygote; a \cle{recessive} allele is only expressed when homozygous. By convention, the dominant allele is written with a capital letter and the recessive allele with the same letter in lower case.
\end{definition}

\begin{definition}[Generations, pure-breeding lines and crosses]
The \cle{P generation} is the parental generation; the \cle{F1} (first filial generation) is the offspring of the P cross; the \cle{F2} is the offspring obtained by crossing F1 individuals among themselves (or self-pollinating them). A \cle{pure-breeding line} is a population which, when self-fertilised or crossed among itself, always produces offspring identical for the character considered --- its members are homozygous. A \cle{back cross} is any cross of an offspring with one of its parents (or with an individual of parental genotype); a \cle{test cross} is a back cross with the homozygous recessive parent, used to reveal an unknown genotype.
\end{definition}

\courssub{Mendel's Monohybrid Experiments}

A \cle{monohybrid cross} follows the inheritance of a \emph{single} pair of contrasting characters. Mendel crossed pure-breeding \textbf{tall} pea plants with pure-breeding \textbf{dwarf} pea plants.

\textbf{The observations}\par
\begin{itemize}
\item \textbf{P:} tall $\times$ dwarf (pure-breeding).
\item \textbf{F1:} \emph{all} the offspring were tall. The dwarf character had apparently ``disappeared''. Mendel called tall \emph{dominant} and dwarf \emph{recessive}.
\item \textbf{F2:} when F1 plants were self-pollinated, the dwarf character \emph{reappeared}. Mendel counted 787 tall and 277 dwarf plants --- a ratio of approximately \textbf{3 tall : 1 dwarf}.
\end{itemize}

The same 3:1 ratio appeared for every one of the seven characters Mendel studied. This constancy demanded an explanation, and Mendel supplied it: each plant carries two ``factors'' (alleles) for each character; these factors \emph{separate} when gametes are formed, so each gamete carries only one; fertilisation restores the pair.

\begin{popfigure}
\begin{center}
\begin{tikzpicture}[font=\small, box/.style={draw=popblue, thick, rounded corners, fill=popblueL, align=center, text width=3.1cm, minimum height=1cm}, arr/.style={->, thick, popdark}]
\node[box, align=center] (P) at (0,0){\textbf{P}\\ Tall $\times$ Dwarf\\ $TT \times tt$};
\node[box, align=center] (G) at (4.6,0){\textbf{Gametes}\\ $T$ \ and \ $t$};
\node[box, align=center] (F1) at (9.2,0){\textbf{F1}\\ all $Tt$\\ all Tall};
\node[box, align=center] (F1x) at (9.2,-2.6){\textbf{F1 $\times$ F1}\\ $Tt \times Tt$};
\node[box, align=center] (F2) at (4.6,-2.6){\textbf{F2}\\ $1\,TT : 2\,Tt : 1\,tt$\\ $3$ Tall : $1$ Dwarf};
\draw[arr] (P) -- (G);
\draw[arr] (G) -- (F1);
\draw[arr] (F1) -- (F1x);
\draw[arr] (F1x) -- (F2);
\end{tikzpicture}
\end{center}
\legende{Mendel's monohybrid cross: pure-breeding tall $\times$ dwarf peas give a uniform tall F1, and selfing the F1 gives an F2 in the phenotypic ratio 3 tall : 1 dwarf.}
\end{popfigure}

\begin{propriete}[Mendel's First Law --- the Law of Segregation]
Each individual carries two alleles for each character; these two alleles \textbf{segregate} (separate) during gamete formation, so that each gamete carries only \textbf{one} allele of each pair. Fertilisation restores the diploid pair at random. (You are expected to state this law precisely and to use it to explain the 3:1 ratio; no formal proof is required, but you must be able to justify it with a genetic diagram.)
\end{propriete}

The law explains the F2 ratio perfectly. The F1 plant ($Tt$) produces two types of gametes, $T$ and $t$, in equal proportions. Random fusion of these gametes gives $TT$, $Tt$, $Tt$, $tt$ --- that is, a genotypic ratio of $1:2:1$ and, because $T$ is dominant, a phenotypic ratio of $3:1$.

\courssub{Genetic Diagrams and the Punnett Square}

\begin{methode}[How to set out a genetic diagram]
Follow these steps in every genetics question --- they carry method marks:
\begin{enumerate}
\item \textbf{Parental phenotypes:} state them in words (e.g.\ Tall $\times$ Dwarf).
\item \textbf{Parental genotypes:} deduce and write them (e.g.\ $TT \times tt$). Pure-breeding means homozygous.
\item \textbf{Gametes:} apply the Law of Segregation --- each parent contributes one allele per gamete. \emph{Circle the gametes.}
\item \textbf{Fertilisation:} combine gametes using a Punnett square (checkerboard) or lines showing fusion.
\item \textbf{Offspring:} state the genotypes, then the genotypic ratio, then the phenotypes and the phenotypic ratio.
\end{enumerate}
\end{methode}

\begin{popfigure}
\begin{center}
\begin{tikzpicture}[font=\small]
\draw[thick, popdark] (0,0) grid (3,3);
\node at (-0.7,2.5) {\textbf{Gametes}};
\node at (0.5,3.45) {$T$};
\node at (1.5,3.45) {$t$};
\node at (-0.5,2.5) {$T$};
\node at (-0.5,1.5) {$t$};
\node[fill=popblueL, minimum size=0.9cm] at (0.5,2.5) {$TT$};
\node[fill=popblueL, minimum size=0.9cm] at (1.5,2.5) {$Tt$};
\node[fill=popblueL, minimum size=0.9cm] at (0.5,1.5) {$Tt$};
\node[fill=popgreenL, minimum size=0.9cm] at (1.5,1.5) {$tt$};
\node[align=center] at (1.5,-0.7) {Genotypes: $1\,TT : 2\,Tt : 1\,tt$};
\node[align=center] at (1.5,-1.3) {Phenotypes: $3$ Tall : $1$ Dwarf};
\end{tikzpicture}
\end{center}
\legende{Punnett square for the heterozygous cross $Tt \times Tt$. Male gametes head the columns, female gametes head the rows; each box is a possible zygote.}
\end{popfigure}

\begin{attention}[Genotype ratio vs phenotype ratio --- the classic GCE trap]
For the cross $Tt \times Tt$, the \textbf{genotypic} ratio is $1:2:1$ ($TT:Tt:tt$) but the \textbf{phenotypic} ratio is $3:1$ (tall:dwarf), because $TT$ and $Tt$ look identical. Candidates routinely write ``3:1'' when asked for genotypes, or ``1:2:1'' when asked for phenotypes. Read the question twice and label each ratio explicitly.
\end{attention}

\begin{exoresolu}[Monohybrid cross in beans]
A pure-breeding bean plant with red seed coats is crossed with a pure-breeding white-seeded plant. The F1 are all red-seeded. (a) Which allele is dominant? (b) Predict the F2 genotypic and phenotypic ratios.
\tcblower
\textbf{(a)} Since the F1 are all red, red is dominant. Let $R$ = red, $r$ = white.\par
\textbf{(b)} Parents: $RR \times rr$; gametes $R$ and $r$; F1 all $Rr$. Crossing F1 $\times$ F1: $Rr \times Rr$; gametes $R$, $r$ from each parent. The Punnett square gives $RR$, $Rr$, $Rr$, $rr$.
\begin{align*}
\text{Genotypic ratio} &= 1\,RR : 2\,Rr : 1\,rr\\
\text{Phenotypic ratio} &= 3\ \text{red} : 1\ \text{white}
\end{align*}
\end{exoresolu}

\courssub{The Test Cross}

Suppose you buy a red-seeded bean plant at the market. Its phenotype tells you it carries at least one $R$ allele --- but is it $RR$ or $Rr$? The phenotype alone cannot answer. Mendel's solution was the \cle{test cross}: cross the unknown individual with a \textbf{homozygous recessive} ($rr$) individual. Because the recessive parent can only contribute $r$ gametes, the offspring phenotypes directly reveal the gametes --- and therefore the genotype --- of the unknown parent.

\begin{popfigure}
\begin{center}
\begin{tikzpicture}[font=\small, box/.style={draw=popgreen, thick, rounded corners, fill=popgreenL, align=center, text width=4.6cm, minimum height=1.6cm}, arr/.style={->, thick, popdark}]
\node[box, align=center] (A) at (0,0){\textbf{If unknown is $Rr$}\\ gametes: $R$, $r$\\ tester $rr$ gives $r$ only\\ offspring: $1\,Rr : 1\,rr$\\ \textbf{1 red : 1 white}};
\node[box, align=center] (B) at (7.4,0){\textbf{If unknown is $RR$}\\ gametes: $R$ only\\ tester $rr$ gives $r$ only\\ offspring: all $Rr$\\ \textbf{all red}};
\draw[arr] (0,-1.6) -- (0,-2.4) node[midway, right, align=left, text width=4.6cm]{};
\draw[arr] (7.4,-1.6) -- (7.4,-2.4);
\node[align=center] at (3.7,-2.9) {A 1:1 ratio in the offspring proves the unknown parent is \textbf{heterozygous}.};
\end{tikzpicture}
\end{center}
\legende{The test cross: crossing an unknown red-seeded plant with a homozygous recessive ($rr$) tester. A 1:1 offspring ratio reveals a heterozygote; all-red offspring reveal a homozygous dominant.}
\end{popfigure}

\begin{exemplebox}[Interpreting a test cross result]
A tall pea plant of unknown genotype is test-crossed with a dwarf plant ($tt$). Out of 80 offspring, 41 are tall and 39 are dwarf. The ratio is approximately $1:1$, so the unknown plant is heterozygous, $Tt$: it produced $T$ and $t$ gametes in equal numbers, and the tester contributed only $t$. Had the unknown been $TT$, all 80 offspring would have been tall.
\end{exemplebox}

\begin{attention}[Exam technique: how to earn the method marks]
Always (i) write the parental phenotypes \emph{and} genotypes, (ii) state the gametes \emph{in circles}, and (iii) show fertilisation clearly with a Punnett square or fusion lines. A correct final ratio with no working earns few marks; a fully labelled diagram earns method marks even if the final counting slips.
\end{attention}

\courssub{From One Character to Two: The Law of Independent Assortment}

Having explained single characters, Mendel asked a bolder question: what happens when \emph{two} characters are inherited together --- for example, seed shape (round/wrinkled) \emph{and} cotyledon colour (yellow/green)? He crossed pure-breeding round-yellow peas with wrinkled-green peas and followed both characters into the F2. He found that the characters behaved \emph{independently}: the 3:1 ratio for shape was combined with the 3:1 ratio for colour in every possible way, giving the famous F2 phenotypic ratio $9:3:3:1$.

\begin{propriete}[Mendel's Second Law --- the Law of Independent Assortment]
Alleles of \emph{different} genes (located on different chromosome pairs) assort \textbf{independently} of one another during gamete formation: the allele a gamete receives for one character does not influence which allele it receives for another character. (Stated here as a bridge; it will be developed fully with dihybrid crosses in Section 4.)
\end{propriete}

\begin{aretenir}[Key points of this section]
\begin{itemize}
\item Genetics studies heredity and variation; Mendel founded it using \emph{Pisum sativum}, chosen for its short cycle, many offspring, contrasting traits and controllable pollination.
\item Genotype $\neq$ phenotype: a dominant phenotype can hide either a homozygous or a heterozygous genotype.
\item Mendel's First Law: the two alleles of a gene segregate into different gametes; each gamete carries one allele.
\item Monohybrid F2: genotypic ratio $1:2:1$; phenotypic ratio $3:1$ (complete dominance).
\item The test cross (unknown $\times$ homozygous recessive) reveals an unknown genotype: a $1:1$ ratio means heterozygous.
\item Mendel's Second Law: genes on different chromosomes assort independently --- the foundation of dihybrid inheritance.
\end{itemize}
\end{aretenir}

\begin{exercice}[Practice]
\begin{enumerate}
\item In guinea pigs, black coat ($B$) is dominant to white coat ($b$). A heterozygous black female is crossed with a white male. Give the expected genotypic and phenotypic ratios of the offspring.
\item A farmer crosses two heterozygous red-seeded bean plants and harvests 400 seeds. How many would you expect to be white-seeded?
\item A tall pea plant is test-crossed and all 60 offspring are tall. What is the genotype of the tall parent? Explain.
\end{enumerate}
\end{exercice}

\begin{corrige}[Exercise 1]
\textbf{1.} Cross: $Bb \times bb$. Gametes: $B$, $b$ from the female; $b$ only from the male. Offspring: $Bb$ and $bb$ in equal proportions. Genotypic ratio $1\,Bb : 1\,bb$; phenotypic ratio $1$ black : $1$ white.\par
\textbf{2.} Cross: $Rr \times Rr$ gives $\tfrac{1}{4}$ $rr$ (white). Expected white seeds: $\tfrac{1}{4} \times 400 = 100$.\par
\textbf{3.} The parent is homozygous dominant, $TT$. A heterozygote ($Tt$) would have produced roughly half dwarf offspring ($1:1$ ratio); since all 60 offspring are tall, the parent contributed only $T$ gametes.
\end{corrige}

\coursec{Patterns of Dominance: Complete, Codominance and Incomplete Dominance}

In the previous section we established the vocabulary of genetics (gene, allele, genotype, phenotype, homozygote, heterozygote) and Mendel's two laws. Mendel was fortunate: in every one of his seven pea characters, one allele completely hid the other in the heterozygote, giving the famous $3:1$ ratio in the $F_2$. But is the heterozygote \emph{always} identical to one of the parents? Careful observation of nature says no. A red-flowered snapdragon crossed with a white-flowered one gives \emph{pink} offspring; a red cow crossed with a white cow gives a \emph{roan} calf carrying both red and white hairs. In this section we study the three great \cle{patterns of dominance} --- complete dominance, incomplete dominance and codominance --- and learn to recognise each one from experimental data, a skill the GCE Advanced Level examiners test almost every year.

\courssub{Complete Dominance}

\textbf{Observation.} A man who can roll his tongue and a woman who cannot have children; all the children can roll their tongues. Yet a grandchild (from two roller parents) may fail to roll. The ``non-rolling'' character disappeared in the $F_1$ and reappeared in the $F_2$. This is exactly what Mendel saw with round and wrinkled pea seeds.

\begin{definition}[Complete dominance]
\cle{Complete dominance} is the pattern of inheritance in which one allele (the \cle{dominant} allele) fully masks the expression of the other allele (the \cle{recessive} allele) in the heterozygote, so that the heterozygote is phenotypically \textbf{indistinguishable} from the homozygous dominant individual.
\end{definition}

\begin{exemplebox}[Complete dominance in the pea and in humans]
In \emph{Pisum sativum}, seed shape is controlled by a gene with two alleles: $R$ (round, dominant) and $r$ (wrinkled, recessive). Genotypes $RR$ and $Rr$ both give round seeds; only $rr$ gives wrinkled seeds. In humans, tongue rolling is traditionally taught in the same way: rollers may be $TT$ or $Tt$, non-rollers $tt$. (Modern studies show tongue rolling is in fact influenced by several factors, but it remains a convenient classroom model of complete dominance.)
\end{exemplebox}

\textbf{Genetic diagram (monohybrid cross, complete dominance).} Crossing two pure-breeding parents ($RR \times rr$), then selfing the $F_1$ ($Rr \times Rr$):

\begin{center}
\begin{tabular}{|l|c|c|}
\hline
\rowcolor{popblueL} $F_1$ gametes & $R$ & $r$ \\
\hline
$R$ & $RR$ (round) & $Rr$ (round) \\
\hline
$r$ & $Rr$ (round) & $rr$ (wrinkled) \\
\hline
\end{tabular}
\end{center}

The $F_2$ \textbf{genotypic ratio} is $1\,RR : 2\,Rr : 1\,rr$, but because $Rr$ looks like $RR$, the \textbf{phenotypic ratio} is $\cle{3:1}$ (3 round : 1 wrinkled).

\begin{attention}[The phenotype hides the genotype]
In complete dominance you \textbf{cannot} tell $RR$ from $Rr$ by looking. To distinguish them, a \cle{test cross} (crossing with the homozygous recessive $rr$) is required: $RR \times rr$ gives all dominant offspring, while $Rr \times rr$ gives a $1:1$ ratio.
\end{attention}

\courssub{Incomplete Dominance}

\textbf{Observation.} A gardener crosses a pure red-flowered snapdragon (\emph{Antirrhinum majus}) with a pure white-flowered one. If red were completely dominant, all the $F_1$ should be red. In fact, every $F_1$ plant is \textbf{pink} --- a colour seen in \emph{neither} parent. Neither allele can fully mask the other.

\begin{definition}[Incomplete dominance]
\cle{Incomplete dominance} (partial dominance) is the pattern of inheritance in which neither allele is completely dominant over the other, so that the heterozygote shows an \textbf{intermediate, blended phenotype} between the two homozygotes.
\end{definition}

A useful intuition: one functional allele produces only \emph{half} the pigment dose of two functional alleles. With two $C^R$ alleles the flower is red; with one $C^R$ and one $C^W$ allele, half the pigment gives pink; with no functional allele, white.

\textbf{Notation.} Since neither allele is dominant, we avoid capital/small letters and write, for example, $C^R$ (red) and $C^W$ (white): both alleles are written as superscripts on the same letter.

\begin{experience}[The snapdragon cross]
Pure red ($C^R C^R$) $\times$ pure white ($C^W C^W$) $\rightarrow$ $F_1$ all pink ($C^R C^W$). Selfing the $F_1$:

\begin{center}
\begin{tabular}{|l|c|c|}
\hline
\rowcolor{popblueL} $F_1$ gametes & $C^R$ & $C^W$ \\
\hline
$C^R$ & $C^R C^R$ (red) & $C^R C^W$ (pink) \\
\hline
$C^W$ & $C^R C^W$ (pink) & $C^W C^W$ (white) \\
\hline
\end{tabular}
\end{center}

$F_2$: $1$ red : $2$ pink : $1$ white. The phenotypic ratio is $\cle{1:2:1}$, identical to the genotypic ratio.
\end{experience}

\begin{popfigure}
\begin{tikzpicture}[scale=1.0]
% Parental generation
\node[align=center] at (1.2,4.6) {\textbf{P}};
\draw[fill=poppink!70!red, draw=popdark] (0,3.6) circle (0.45);
\node[align=center, text width=2.2cm] at (0,2.7) {Red\\ $C^R C^R$};
\draw[fill=white, draw=popdark] (2.4,3.6) circle (0.45);
\node[align=center, text width=2.2cm] at (2.4,2.7) {White\\ $C^W C^W$};
\node at (1.2,3.6) {$\times$};
\draw[->, thick, popblue] (1.2,2.3) -- (1.2,1.7);
% F1
\node[align=center] at (1.2,1.35) {\textbf{F}$_1$};
\draw[fill=poppink, draw=popdark] (1.2,0.5) circle (0.45);
\node[align=center, text width=2.4cm] at (1.2,-0.4) {All pink\\ $C^R C^W$};
\draw[->, thick, popblue] (2.0,0.5) -- (3.0,0.5);
\node[align=center] at (3.6,0.5) {self};
% F2
\node[align=center] at (6.4,2.1) {\textbf{F}$_2$};
\draw[fill=poppink!70!red, draw=popdark] (5.2,0.5) circle (0.4);
\node[align=center, text width=1.8cm] at (5.2,-0.4) {1 red};
\draw[fill=poppink, draw=popdark] (6.4,0.5) circle (0.4);
\draw[fill=poppink, draw=popdark] (6.4,1.35) circle (0.4);
\node[align=center, text width=1.8cm] at (6.4,-0.4) {2 pink};
\draw[fill=white, draw=popdark] (7.6,0.5) circle (0.4);
\node[align=center, text width=1.8cm] at (7.6,-0.4) {1 white};
\end{tikzpicture}
\legende{Incomplete dominance in the snapdragon: red $\times$ white gives pink $F_1$, and a $1:2:1$ ratio in the $F_2$.}
\end{popfigure}

\courssub{Codominance}

\textbf{Observation.} A cattle breeder crosses a pure red bull with a pure white cow. The calf is neither red, nor white, nor pink: it is \textbf{roan} --- its coat contains \emph{both} red hairs \emph{and} white hairs, side by side. Both alleles are fully and separately expressed.

\begin{definition}[Codominance]
\cle{Codominance} is the pattern of inheritance in which \textbf{both alleles are fully and simultaneously expressed} in the heterozygote, so that the heterozygote shows \textbf{both} parental phenotypes together (not a blend).
\end{definition}

\begin{exemplebox}[Roan cattle and the MN blood group]
\textbf{(a) Roan coat in cattle.} With alleles $C^R$ (red) and $C^W$ (white): $C^R C^R$ = red, $C^W C^W$ = white, $C^R C^W$ = roan (red \emph{and} white hairs intermingled). Crossing two roan animals gives $1$ red : $2$ roan : $1$ white.

\textbf{(b) MN blood group in humans.} The gene $L$ has alleles $L^M$ and $L^N$, each coding for a distinct glycoprotein (antigen) on red blood cells. Genotype $L^M L^M$ = group M; $L^N L^N$ = group N; $L^M L^N$ = group MN, in which \emph{both} antigens M and N are detected on the same red cells. Two MN parents produce children in the ratio $1$ M : $2$ MN : $1$ N.
\end{exemplebox}

\begin{propriete}[Phenotypic ratio equals genotypic ratio]
In \textbf{both} codominance and incomplete dominance, every genotype has its own distinguishable phenotype. Consequently the $F_2$ \textbf{phenotypic ratio equals the genotypic ratio}: $1:2:1$. (In complete dominance the phenotypic ratio is $3:1$ because two genotypes share one phenotype.) You should be able to justify this from a Punnett square --- it is frequently asked.
\end{propriete}

\begin{attention}[Codominance is NOT a blend]
Do not confuse the two patterns. In \textbf{incomplete dominance} the heterozygote is a \emph{mixture/intermediate}: pink flowers, where no red or white petal exists. In \textbf{codominance} both traits appear \emph{separately and together}: a roan cow has distinct red hairs and distinct white hairs; an MN person has both M and N antigens. Pink $\neq$ roan. Examiners penalise candidates who write ``pink cow''.
\end{attention}

\courssub{Comparing the Three Patterns}

\begin{popfigure}
\begin{tikzpicture}
\begin{axis}[
    ybar, bar width=0.55cm,
    width=12cm, height=6.2cm,
    ymin=0, ymax=3.4,
    ylabel={Number of parts in $F_2$},
    symbolic x coords={Dom. phenotype, Heterozygote, Rec. phenotype},
    xtick=data,
    x tick label style={align=center, font=\small},
    legend style={at={(0.5,1.02)}, anchor=south, legend columns=3, font=\small},
    ytick={0,1,2,3},
]
\addplot[fill=popblue] coordinates {(Dom. phenotype,3) (Heterozygote,0) (Rec. phenotype,1)};
\addplot[fill=poporange] coordinates {(Dom. phenotype,1) (Heterozygote,2) (Rec. phenotype,1)};
\addplot[fill=popgreen] coordinates {(Dom. phenotype,1) (Heterozygote,2) (Rec. phenotype,1)};
\legend{Complete (3:1), Incomplete (1:2:1), Codominance (1:2:1)}
\end{axis}
\end{tikzpicture}
\legende{Comparison of $F_2$ phenotypic ratios. In complete dominance the heterozygote is grouped with the dominant phenotype ($3:1$); in the other two patterns it forms its own class ($1:2:1$).}
\end{popfigure}

\begin{center}
\begin{tabularx}{\linewidth}{|l|X|X|X|}
\hline
\rowcolor{popblueL}
\textbf{Feature} & \textbf{Complete dominance} & \textbf{Incomplete dominance} & \textbf{Codominance} \\
\hline
Heterozygote appearance & Like the dominant homozygote & Intermediate blend (e.g. pink) & Both traits shown together (e.g. roan) \\
\hline
$F_2$ genotypic ratio & $1:2:1$ & $1:2:1$ & $1:2:1$ \\
\hline
$F_2$ phenotypic ratio & $3:1$ & $1:2:1$ & $1:2:1$ \\
\hline
Number of $F_2$ phenotypes & 2 & 3 & 3 \\
\hline
Examples & Pea seed shape; tongue rolling & Snapdragon flower colour & Roan cattle; MN blood group \\
\hline
\end{tabularx}
\end{center}

\begin{aretenir}[Key points]
\begin{itemize}
\item The $F_2$ \textbf{genotypic} ratio of a monohybrid cross is \emph{always} $1:2:1$; only the \textbf{phenotypic} ratio changes with the dominance pattern.
\item Complete dominance $\rightarrow 3:1$; incomplete dominance and codominance $\rightarrow 1:2:1$.
\item Incomplete dominance = blended intermediate; codominance = both traits expressed side by side.
\end{itemize}
\end{aretenir}

\courssub{Identifying the Dominance Pattern from Exam Data}

\begin{methode}[How to deduce the dominance pattern]
\begin{enumerate}
\item \textbf{Look at the $F_1$} (or the offspring of two pure-breeding contrasting parents). If all offspring resemble one parent $\rightarrow$ complete dominance. If they show a new, intermediate phenotype $\rightarrow$ suspect incomplete dominance. If they show both parental traits together $\rightarrow$ suspect codominance.
\item \textbf{Compute the offspring ratio} of the $F_2$ (or of a heterozygote $\times$ heterozygote cross). Convert the numbers to the simplest ratio (divide by the smallest number).
\item \textbf{Match the ratio:} $3:1$ $\rightarrow$ complete dominance; $1:2:1$ with an intermediate class $\rightarrow$ incomplete dominance; $1:2:1$ with a ``both-traits'' class $\rightarrow$ codominance; $1:1$ $\rightarrow$ test cross (heterozygote $\times$ homozygous recessive).
\item \textbf{Check with a genetic diagram:} assign symbols, predict the offspring, and confirm your prediction matches the data exactly.
\end{enumerate}
\end{methode}

\begin{attention}[Exam tip]
When the heterozygote differs from \textbf{both} parents, immediately suspect incomplete dominance or codominance --- then decide by asking: is the heterozygote a \emph{blend} (incomplete) or does it show \emph{both traits separately} (codominance)?
\end{attention}

\begin{exoresolu}[Deducing the pattern from data]
A farmer crosses a pure-breeding red-flowered plant with a pure-breeding white-flowered plant. All 40 $F_1$ plants have pink flowers. When two $F_1$ plants are crossed, 120 seeds germinate and give: 29 red, 61 pink, 30 white. (a) Explain the inheritance pattern. (b) Give the genotypes of all plants. (c) State the expected numbers and compare with the observed numbers.
\tcblower
\textbf{(a)} The $F_1$ shows a new intermediate phenotype (pink), different from both parents: this is \cle{incomplete dominance}. The $F_2$ ratio $29:61:30 \approx 1:2:1$ confirms it.

\textbf{(b)} Let $C^R$ = red allele, $C^W$ = white allele. Parents: $C^R C^R$ (red) $\times$ $C^W C^W$ (white); $F_1$: $C^R C^W$ (pink); $F_2$: $C^R C^R$ red, $C^R C^W$ pink, $C^W C^W$ white.

\textbf{(c)} Expected ratio $1:2:1$ on 120 plants:
\[
\text{red} = \tfrac{1}{4}\times 120 = 30, \quad
\text{pink} = \tfrac{2}{4}\times 120 = 60, \quad
\text{white} = \tfrac{1}{4}\times 120 = 30.
\]
Observed ($29, 61, 30$) matches expected ($30, 60, 30$) very closely; small differences are due to chance sampling.
\end{exoresolu}

\begin{exoresolu}[Codominance in cattle]
Two roan cattle are crossed. (a) Give the expected genotypic and phenotypic ratios among the calves. (b) A farmer wants a herd of only roan cattle. Explain, using a genetic diagram, why crossing roan $\times$ roan can never achieve this, and suggest the cross that gives $100\%$ roan.
\tcblower
\textbf{(a)} Let $C^R$ = red, $C^W$ = white. Roan $\times$ roan: $C^R C^W \times C^R C^W$ gives $1\,C^R C^R : 2\,C^R C^W : 1\,C^W C^W$, i.e. $1$ red : $2$ roan : $1$ white. Genotypic and phenotypic ratios are both $1:2:1$.

\textbf{(b)} Because the roan animal is a heterozygote, its gametes segregate ($\tfrac{1}{2}C^R$, $\tfrac{1}{2}C^W$), so pure red and pure white calves will always reappear (each with probability $\tfrac{1}{4}$). To obtain $100\%$ roan, cross a \textbf{red} ($C^R C^R$) with a \textbf{white} ($C^W C^W$): every calf receives $C^R$ from one parent and $C^W$ from the other, so all are $C^R C^W$ (roan).
\end{exoresolu}

\begin{savaistu}[Dominance patterns around us in Cameroon]
Local agriculture offers many revision anchors. In some maize varieties grown in the West and North-West Regions, kernel colour shows incomplete dominance: crossing deep-purple with white kernels can yield intermediate pale-purple cobs. Among local poultry, crosses between black and white chickens sometimes produce ``blue'' or speckled birds --- speckled plumage, where black and white feathers both appear, recalls codominance, while an evenly diluted ``blue'' grey recalls incomplete dominance. When you revise, attach each pattern to one Cameroonian image: $3:1$ (tongue rolling in your family), pink snapdragon ($1:2:1$ blend), roan calf ($1:2:1$ both colours).
\end{savaistu}

\begin{exercice}[Practice]
\begin{enumerate}
\item In a certain plant, $F^R F^R$ gives red flowers, $F^R F^W$ gives pink and $F^W F^W$ gives white. A pink-flowered plant is crossed with a white-flowered plant. Give the expected genotypic and phenotypic ratios of the offspring.
\item A man of blood group MN marries a woman of group M. What blood groups are possible among their children, and in what proportions?
\item Two pink snapdragons produce 80 offspring. How many of each phenotype do you expect?
\item In guinea pigs, black coat ($B$) is completely dominant over white ($b$). A black male whose father was white is crossed with a white female. Predict the offspring.
\end{enumerate}
\end{exercice}

\begin{corrige}[Exercise 1]
\textbf{1.} $F^R F^W \times F^W F^W$: gametes $F^R, F^W$ and $F^W$ only. Offspring: $\tfrac{1}{2}F^R F^W$ (pink) : $\tfrac{1}{2}F^W F^W$ (white). Ratio $1:1$; no red offspring possible.

\textbf{2.} $L^M L^N \times L^M L^M$: offspring $\tfrac{1}{2}L^M L^M$ (group M) and $\tfrac{1}{2}L^M L^N$ (group MN). Group N is impossible.

\textbf{3.} Pink $\times$ pink gives $1:2:1$; on 80 offspring: $20$ red, $40$ pink, $20$ white.

\textbf{4.} The black male must be $Bb$ (he received $b$ from his white father, $bb$). $Bb \times bb$ gives $\tfrac{1}{2}Bb$ black : $\tfrac{1}{2}bb$ white, a $1:1$ test-cross ratio.
\end{corrige}

\coursec{Sex Determination and Sex Linkage}

In the previous section we studied how alleles of a single gene interact at a locus: complete dominance, codominance and incomplete dominance. In all those cases, the two sexes were affected in exactly the same way --- a heterozygous four o'clock plant is pink whether it is male or female. Yet careful observation of human families tells a different story for certain traits: red-green colour blindness and haemophilia appear almost exclusively in boys, and a colour-blind grandfather often transmits the condition to his grandsons through his daughter, who is herself perfectly normal. Traits whose inheritance depends on sex cannot be explained by the simple Mendelian ratios we have met so far. To understand them, we must first ask a more fundamental question: \emph{what makes an individual male or female?} The answer lies in the chromosomes themselves, and it opens the door to one of the most elegant chapters of genetics: \cle{sex determination} and \cle{sex linkage}.

\courssub{The Chromosomal Basis of Sex}

\textbf{Autosomes and sex chromosomes.}\par
When the chromosomes of a human cell are photographed at metaphase of mitosis, then cut out and arranged in homologous pairs in order of decreasing size, the resulting picture is called a \cle{karyotype}. A human karyotype shows 46 chromosomes, that is 23 pairs. Twenty-two of these pairs look identical in males and in females: they are called \cle{autosomes}, and they carry genes for ordinary body characteristics (eye colour, blood group, enzymes, and so on). The 23rd pair is different: it is the pair of \cle{sex chromosomes}, and its composition differs between the sexes.

\begin{definition}[Autosomes and sex chromosomes]
\cle{Autosomes} are chromosomes that are identical in morphology and number in both sexes of a species. \cle{Sex chromosomes} (also called heterosomes or gonosomes) are chromosomes that differ between the sexes and carry genes involved in sex determination. In humans there are 44 autosomes (22 pairs) and 2 sex chromosomes: \textbf{XX in the female}, \textbf{XY in the male}.
\end{definition}

The X chromosome is large and carries hundreds of genes, most of which have nothing to do with sex (genes for blood clotting factors, for colour vision, for muscle proteins). The Y chromosome is much smaller and carries very few genes; its essential role is to carry the \emph{SRY} gene (\emph{Sex-determining Region of Y}), which switches the developing embryonic gonad towards becoming a testis. The testis then secretes hormones that direct the development of male characteristics. In the absence of a Y chromosome, the gonad develops into an ovary. This is why, in humans, \textbf{it is the presence or absence of the Y chromosome that determines maleness}.

\begin{popfigure}
\begin{center}
\begin{tikzpicture}[scale=1]
% X chromosome
\draw[thick, popblue] (0,1.6) -- (0.35,0.9) -- (0,0.2);
\draw[thick, popblue] (0.7,1.6) -- (0.35,0.9) -- (0.7,0.2);
\draw[thick, popblue] (0,1.6) -- (0.7,1.6);
\draw[thick, popblue] (0,0.2) -- (0.7,0.2);
\fill[popdark] (0.35,0.9) circle (0.06);
\node[popblue, font=\bfseries] at (0.35,-0.35) {X};
% Y chromosome
\draw[thick, poporange] (2.6,1.15) -- (2.85,0.75) -- (2.6,0.35);
\draw[thick, poporange] (3.1,1.15) -- (2.85,0.75) -- (3.1,0.35);
\draw[thick, poporange] (2.6,1.15) -- (3.1,1.15);
\draw[thick, poporange] (2.6,0.35) -- (3.1,0.35);
\fill[popdark] (2.85,0.75) circle (0.06);
\node[poporange, font=\bfseries] at (2.85,-0.35) {Y};
% braces and labels
\draw[decorate, decoration={brace, amplitude=4pt}, popmuted] (-0.25,1.7) -- (-0.25,0.1);
\node[popmuted, font=\small, align=center, text width=3.2cm] at (-2.3,0.9) {large, many genes (clotting, vision)};
\draw[decorate, decoration={brace, amplitude=4pt}, popmuted] (3.35,1.25) -- (3.35,0.25);
\node[popmuted, font=\small, align=center, text width=3.2cm] at (5.4,0.75) {small, few genes (SRY)};
% non-homologous regions
\node[font=\small, align=center, text width=4.5cm, popink] at (1.75,2.3) {The X and Y pair only at tiny homologous tips; most of X has NO partner on Y};
\end{tikzpicture}
\end{center}
\legende{Comparison of the human X and Y chromosomes. The X is large and gene-rich; the Y is small. Because most of the X has no homologous region on the Y, genes on the X are present in a single copy in males.}
\end{popfigure}

\courssub{Sex Determination in Humans: the Genetic Diagram}

Because the female is XX, every ovum she produces carries one X chromosome: the female is the \cle{homogametic sex} (she produces only one type of gamete with respect to the sex chromosomes). The male is XY, so meiosis in the testis produces two classes of spermatozoa in equal numbers: half carry an X, half carry a Y. The male is the \cle{heterogametic sex}. Fertilisation is random, so an X-bearing sperm has exactly the same chance of fertilising the ovum as a Y-bearing sperm.

\begin{popfigure}
\begin{center}
\begin{tikzpicture}[scale=1]
% Parents
\node[font=\bfseries, poppink] at (0,3.4) {Mother};
\node[font=\bfseries, popblue] at (6,3.4) {Father};
\node at (0,2.8) {XX};
\node at (6,2.8) {XY};
\node[font=\small, popmuted] at (0,2.3) {(all ova: X)};
\node[font=\small, popmuted] at (6,2.3) {(sperm: X or Y)};
% gametes arrows
\draw[->, thick, popmuted] (0,2.0) -- (1.2,1.35);
\draw[->, thick, popmuted] (6,2.0) -- (4.8,1.35);
% Punnett square
\draw[thick, popink] (1.5,1.2) rectangle (4.5,-1.2);
\draw[thick, popink] (3,1.2) -- (3,-1.2);
\draw[thick, popink] (1.5,0) -- (4.5,0);
\node[font=\bfseries, poppink] at (1.05,0.6) {X};
\node[font=\bfseries, poppink] at (1.05,-0.6) {X};
\node[font=\bfseries, popblue] at (2.25,1.65) {X};
\node[font=\bfseries, popblue] at (3.75,1.65) {Y};
\node at (2.25,0.6) {XX};
\node[font=\small, poppink] at (2.25,0.15) {girl};
\node at (3.75,0.6) {XY};
\node[font=\small, popblue] at (3.75,0.15) {boy};
\node at (2.25,-0.6) {XX};
\node[font=\small, poppink] at (2.25,-1.05) {girl};
\node at (3.75,-0.6) {XY};
\node[font=\small, popblue] at (3.75,-1.05) {boy};
% ratio
\node[font=\bfseries, popdark, align=center] at (3,-2.1) {Offspring ratio: 1 girl : 1 boy (50\% : 50\%)};
\end{tikzpicture}
\end{center}
\legende{Genetic diagram of human sex determination. The mother produces only X-bearing ova; the father produces equal numbers of X- and Y-bearing sperm. Fertilisation gives a 1:1 sex ratio.}
\end{popfigure}

\begin{aretenir}[Key points on human sex determination]
\begin{itemize}
\item The expected sex ratio at fertilisation is \textbf{1 male : 1 female}, because the two classes of sperm are produced in equal proportions and fertilisation is random.
\item \textbf{The father determines the sex of the child} in humans: the ovum always contributes an X, so it is the sperm (X or Y) that fixes the sex of the zygote. It is therefore biologically wrong to blame a woman for producing ``only girls'' --- a prejudice still encountered in some families.
\item Each child is an independent event: a couple with three daughters still has a probability of $\tfrac{1}{2}$ of having a son at the next pregnancy.
\end{itemize}
\end{aretenir}

\courssub{Sex Determination in Other Organisms}

The XX/XY system is not universal. Nature has invented several solutions, and the GCE examiner expects you to know the main ones.

\begin{center}
\begin{tabularx}{\linewidth}{|l|l|l|X|}
\hline
\rowcolor{popblueL} \textbf{System} & \textbf{Female} & \textbf{Male} & \textbf{Examples} \\
\hline
XX/XY & XX (homogametic) & XY (heterogametic) & Humans, most mammals, \emph{Drosophila} \\
\hline
ZW/ZZ & ZW (heterogametic) & ZZ (homogametic) & Birds, butterflies, moths, some reptiles \\
\hline
XX/XO & XX & XO (no second sex chromosome) & Grasshoppers, cockroaches, some bugs \\
\hline
Haplodiploidy & Diploid ($2n$, fertilised egg) & Haploid ($n$, unfertilised egg) & Bees, wasps, ants \\
\hline
\end{tabularx}
\end{center}

Note carefully the \textbf{ZW/ZZ system}: here the situation is \emph{reversed} compared with humans. The \emph{female} bird is the heterogametic sex (ZW) and produces two classes of ova (Z and W), while the male (ZZ) produces only Z sperm. It is therefore the \textbf{mother who determines the sex of the chick} --- the exact opposite of the human situation.

In the \textbf{XX/XO system} (grasshoppers), the male simply has one X chromosome and no second sex chromosome at all (the ``O'' means ``absent''). He produces two classes of sperm: with or without an X.

In \textbf{haplodiploidy} (bees), there are no sex chromosomes at all. A fertilised (diploid) egg develops into a female (queen or worker); an unfertilised (haploid) egg develops into a male (drone) by parthenogenesis. The queen controls fertilisation by releasing or withholding stored sperm as the egg passes.

\begin{savaistu}[Extension: environmental sex determination]
In some reptiles, including many turtles and crocodilians, sex is not fixed by chromosomes at all but by the \textbf{temperature of incubation} of the eggs. In several turtle species, eggs incubated at warm temperatures produce females and cooler nests produce males. This is \cle{environmental sex determination}. It is not examinable in detail at GCE, but it is a striking reminder that ``sex'' is a phenotype, and like all phenotypes it results from the interaction of genotype and environment.
\end{savaistu}

\courssub{Sex Linkage: Genes Carried on the X Chromosome}

We now return to our opening puzzle: why do some diseases run through families, striking boys far more often than girls? The answer is that the genes responsible are located \emph{on the X chromosome itself}, so their inheritance follows the inheritance of sex.

\begin{definition}[Sex linkage]
A gene is said to be \cle{sex-linked} when it is located on a sex chromosome (in practice, almost always the X chromosome), so that its pattern of inheritance is associated with the sex of the individual. The two classic X-linked recessive conditions in humans are \textbf{haemophilia} (failure of blood to clot, due to a defective clotting factor) and \textbf{red-green colour blindness} (inability to distinguish red from green).
\end{definition}

\begin{definition}[Hemizygosity]
A male is said to be \cle{hemizygous} for X-linked genes because he possesses only \textbf{one X chromosome} (and the Y carries no corresponding allele for most X-linked loci). He therefore has only \emph{one copy} of each X-linked gene, and whatever allele that copy carries --- normal or defective --- will be expressed in his phenotype.
\end{definition}

This single fact explains the whole epidemiology of X-linked recessive disorders. A female, with two X chromosomes, can be:
\begin{itemize}
\item \textbf{homozygous normal} ($X^H X^H$): normal, non-carrier;
\item \textbf{heterozygous} ($X^H X^h$): phenotypically normal but a \cle{carrier} --- the normal allele dominates;
\item \textbf{homozygous recessive} ($X^h X^h$): affected (very rare, since it requires an affected father \emph{and} at least a carrier mother).
\end{itemize}
A male, with only one X, can only be $X^H Y$ (normal) or $X^h Y$ (affected). There is no heterozygous ``carrier'' state for him.

\begin{attention}[A male cannot be a ``carrier'' of an X-linked recessive trait]
Because a male is hemizygous, the single allele on his X chromosome is always expressed. A man is therefore either \textbf{normal} ($X^H Y$) or \textbf{affected} ($X^h Y$) --- never a ``carrier''. Writing $X^H X^h$ for a male, or describing a healthy man as a carrier of haemophilia, is a serious error that costs marks. Only \emph{females} can be carriers.
\end{attention}

\begin{methode}[Writing sex-linked crosses correctly]
\begin{enumerate}
\item Identify the gene as X-linked (the question will state or imply this: ``haemophilia'', ``colour blindness'', ``sex-linked'').
\item Choose the allele symbols and write them as \textbf{superscripts on the X}: e.g.\ $X^H$ (normal clotting) dominant over $X^h$ (haemophilia); $X^B$ (normal vision) over $X^b$ (colour blindness).
\item Write genotypes \emph{with the sex chromosomes}: female $X^H X^h$, male $X^h Y$. \textbf{Never} write $Hh$ or $XhXh$ alone --- the superscript must sit on an X, and the Y must appear in male genotypes.
\item List the gametes: a carrier female gives $X^H$ and $X^h$; a normal male gives $X^H$ and $Y$.
\item Complete the Punnett square and state the phenotypic ratio \emph{separately for each sex}.
\end{enumerate}
\end{methode}

\begin{exoresolu}[Carrier mother $\times$ normal father (haemophilia)]
A woman who is a carrier of haemophilia marries a normal man. Using suitable symbols, determine the expected phenotypes of their children.
\tcblower
Let $X^H$ = X chromosome carrying the normal clotting allele (dominant); $X^h$ = X chromosome carrying the haemophilia allele (recessive).

\textbf{Parental phenotypes:} carrier mother $\times$ normal father

\textbf{Parental genotypes:} $X^H X^h \;\times\; X^H Y$

\textbf{Gametes:} mother: $X^H$, $X^h$ \quad father: $X^H$, $Y$

\begin{center}
\begin{tabular}{|c|c|c|}
\hline
\rowcolor{popblueL} & $X^H$ (father) & $Y$ (father) \\
\hline
$X^H$ (mother) & $X^H X^H$ normal girl & $X^H Y$ normal boy \\
\hline
$X^h$ (mother) & $X^H X^h$ carrier girl (normal) & $X^h Y$ \textbf{haemophiliac boy} \\
\hline
\end{tabular}
\end{center}

\textbf{Offspring:} all daughters phenotypically normal (half of them carriers); among sons, $\tfrac{1}{2}$ normal and $\tfrac{1}{2}$ haemophiliac. Overall: \textbf{no affected daughter, 50\% of sons affected}. Note that the father gives his Y to his sons, so a son's X --- and therefore his clotting allele --- comes entirely from his mother.
\end{exoresolu}

\begin{popfigure}
\begin{center}
\begin{tikzpicture}[scale=1]
% Parents
\node[font=\bfseries, poppink] at (0,3.6) {Carrier mother};
\node[font=\bfseries, popblue] at (6.4,3.6) {Normal father};
\node at (0,3.0) {$X^H X^h$};
\node at (6.4,3.0) {$X^H Y$};
\draw[->, thick, popmuted] (0,2.55) -- (1.3,1.9);
\draw[->, thick, popmuted] (6.4,2.55) -- (5.1,1.9);
% Punnett square
\draw[thick, popink] (1.6,1.6) rectangle (4.8,-1.6);
\draw[thick, popink] (3.2,1.6) -- (3.2,-1.6);
\draw[thick, popink] (1.6,0) -- (4.8,0);
\node[font=\bfseries, poppink] at (1.1,0.8) {$X^H$};
\node[font=\bfseries, poppink] at (1.1,-0.8) {$X^h$};
\node[font=\bfseries, popblue] at (2.4,2.05) {$X^H$};
\node[font=\bfseries, popblue] at (4.0,2.05) {$Y$};
\node at (2.4,0.8) {$X^H X^H$};
\node[font=\small, poppink] at (2.4,0.35) {normal girl};
\node at (4.0,0.8) {$X^H Y$};
\node[font=\small, popblue] at (4.0,0.35) {normal boy};
\node at (2.4,-0.8) {$X^H X^h$};
\node[font=\small, poppink] at (2.4,-1.25) {carrier girl};
\node[fill=popblueL, rounded corners=2pt, inner sep=2pt] at (4.0,-0.8) {$X^h Y$};
\node[font=\small, poporange] at (4.0,-1.25) {haemophiliac boy};
\node[font=\bfseries, popdark] at (3.2,-2.4) {50\% of sons affected; all daughters normal (half carriers)};
\end{tikzpicture}
\end{center}
\legende{Inheritance of haemophilia from a carrier mother ($X^H X^h$) and a normal father ($X^H Y$). Only sons can be affected, because they receive their single X from their mother.}
\end{popfigure}

\begin{exoresolu}[Colour-blind father $\times$ homozygous normal mother]
A colour-blind man marries a woman with normal vision whose family has no history of colour blindness. Predict the genotypes and phenotypes of their children, and of the children their \emph{daughters} might have with normal men.
\tcblower
Let $X^B$ = normal vision (dominant), $X^b$ = colour blindness (recessive).

\textbf{First cross:} $X^B X^B$ (normal mother) $\times$ $X^b Y$ (colour-blind father)

Gametes: mother $X^B$ only; father $X^b$ and $Y$.

Offspring: daughters all $X^B X^b$ (normal vision, \textbf{carriers}); sons all $X^B Y$ (normal). \textbf{No child is affected}, but \emph{every daughter is a carrier}, because each daughter necessarily receives her father's only X chromosome, which carries $X^b$.

\textbf{Second cross (daughter $\times$ normal man):} $X^B X^b \times X^B Y$ --- this is exactly the carrier-mother cross solved above: half of the grandsons will be colour-blind.

This is the classic pattern: the affected grandfather transmits the allele silently through his daughter to half of his grandsons.
\end{exoresolu}

\begin{propriete}[Criss-cross inheritance]
A male transmits his \textbf{X chromosome to all of his daughters} and his \textbf{Y chromosome to all of his sons}. Consequently, an X-linked trait passes from father to daughter (who is usually an unaffected carrier) and then from mother to son: the trait appears to ``criss-cross'' from one sex to the other in alternate generations. This \cle{criss-cross inheritance} is the diagnostic signature of X-linked recessive transmission in a pedigree. (Proof not examinable; it follows directly from the sex-determination diagram.)
\end{propriete}

\begin{attention}[Exam tip --- the word ``carrier'']
In GCE questions, the phrase ``a \textbf{carrier female}'' \emph{always} means a heterozygous woman, e.g.\ $X^H X^h$, who is phenotypically normal. State this explicitly in your answer (``the carrier mother has genotype $X^H X^h$'') before drawing your Punnett square. Conversely, if a question says a woman is ``normal with no family history'', assume she is homozygous $X^H X^H$ unless told otherwise.
\end{attention}

\courssub{Y-linked (Holandric) Inheritance --- Supplementary}

A very small number of genes lie on the part of the Y chromosome that has no counterpart on the X. Such genes are called \cle{holandric} (``wholly male'') genes. Because only males possess a Y chromosome, a Y-linked trait is transmitted \textbf{exclusively from father to son, to all sons, and never to daughters}. The textbook example is the \emph{hairy pinna} trait (tufts of hair on the outer ear) reported in some families. Holandric inheritance is mentioned for completeness; it is rarely examined, but you should be able to distinguish its father-to-all-sons pattern from the criss-cross pattern of X-linked genes.

\begin{aretenir}[Summary of the section]
\begin{itemize}
\item Humans: 44 autosomes + XX (female) or XY (male); the \textbf{father} determines the child's sex; expected sex ratio 1:1.
\item Other systems: ZW/ZZ (birds, butterflies --- \emph{female} heterogametic), XO (grasshoppers), haplodiploidy (bees); environmental sex determination in some reptiles (extension).
\item \textbf{Sex-linked genes} are carried on the X chromosome; males are \textbf{hemizygous}, so X-linked recessive disorders (haemophilia, colour blindness) are far commoner in males.
\item Correct notation: $X^H X^h$, $X^h Y$ --- superscripts on the X, Y written bare.
\item A male is either normal or affected, \textbf{never a carrier}; ``carrier'' in an exam question means a heterozygous female.
\item \textbf{Criss-cross inheritance}: grandfather $\rightarrow$ carrier daughter $\rightarrow$ affected grandson.
\end{itemize}
\end{aretenir}

\begin{exercice}[Practice]
\begin{enumerate}
\item A woman with normal vision, whose father was colour-blind, marries a man with normal vision. Determine the probability that (a) their first child is a colour-blind son; (b) a given son is colour-blind; (c) a given daughter is a carrier.
\item In poultry, the gene for barred plumage is carried on the Z chromosome and is dominant over non-barred. A barred cock (homozygous) is crossed with non-barred hens. Recalling that the hen is ZW, predict the phenotypes of the chicks of each sex.
\item Explain, using the word \emph{hemizygous}, why haemophilia is much commoner in males than in females, and why an affected female ($X^h X^h$) is extremely rare.
\end{enumerate}
\end{exercice}

\begin{corrige}[Exercise 1]
\textbf{1.} The woman's father was colour-blind ($X^b Y$), so he gave his only X, carrying $X^b$, to his daughter. The woman has normal vision, so her other X must carry $X^B$: she is an obligate carrier, $X^B X^b$. The husband is normal: $X^B Y$. The Punnett square is the carrier-mother cross: offspring $X^B X^B$, $X^B X^b$, $X^B Y$, $X^b Y$, each with probability $\tfrac{1}{4}$.
(a) P(colour-blind son) $= \tfrac{1}{4}$. (b) Among sons, half are $X^b Y$: probability $\tfrac{1}{2}$. (c) Among daughters, half are $X^B X^b$: probability $\tfrac{1}{2}$ (and all daughters have normal vision).

\textbf{2.} Cock: $Z^B Z^B$ (barred, homozygous); hen: $Z^b W$ (non-barred). Gametes: cock gives $Z^B$ only; hen gives $Z^b$ or $W$. Offspring: $Z^B Z^b$ barred cockerels and $Z^B W$ barred pullets --- \emph{all} chicks barred. (Note the reversal: had the cock been non-barred $Z^b Z^b$ and the hen barred $Z^B W$, the sons would all be barred carriers and the daughters all non-barred --- criss-cross inheritance in the ZW system.)

\textbf{3.} A male is hemizygous for X-linked genes: he has only one X chromosome, so a single recessive allele $X^h$ is enough to produce the disease (genotype $X^h Y$). A female has two X chromosomes and must inherit $X^h$ from \emph{both} parents to be affected ($X^h X^h$), which requires an affected father \emph{and} a carrier or affected mother --- a very unlikely combination, especially since haemophiliacs historically had reduced survival. Hence affected females are extremely rare while affected males are comparatively common.
\end{corrige}

\coursec{Dihybrid Inheritance, Linkage and Crossing Over}

In the previous section we saw how genes carried on the sex chromosomes produce characteristic criss-cross patterns of inheritance. We now widen our view and ask a new question: \emph{what happens when we follow two pairs of alleles at the same time?} Mendel himself answered this question with his pea plants, and the answer led him to his Second Law. Later, Thomas Hunt Morgan's work on the fruit fly \emph{Drosophila melanogaster} showed that Mendel's Second Law has an important exception --- \cle{linkage} --- and that the exceptions themselves, through \cle{crossing over}, allow us to map the positions of genes on chromosomes.

\courssub{The Dihybrid Cross and Mendel's Second Law}

\begin{definition}[Dihybrid cross]
A \cle{dihybrid cross} is a cross between individuals that differ in \textbf{two} contrasting characters (two pairs of alleles), for example seed shape (round/wrinkled) and seed colour (yellow/green) in peas. An individual heterozygous for both genes, e.g.\ $RrYy$, is called a \cle{dihybrid}.
\end{definition}

\textbf{Mendel's classic experiment.} Mendel crossed pure-breeding pea plants with \textbf{round, yellow} seeds ($RRYY$) with plants having \textbf{wrinkled, green} seeds ($rryy$). Round ($R$) is dominant to wrinkled ($r$); yellow ($Y$) is dominant to green ($y$).

\begin{itemize}
\item \textbf{F$_1$ generation:} all offspring are $RrYy$ --- round and yellow, as expected from complete dominance.
\item \textbf{F$_2$ generation:} when the F$_1$ dihybrids were self-pollinated ($RrYy \times RrYy$), four phenotypes appeared in the ratio \textbf{9 round yellow : 3 round green : 3 wrinkled yellow : 1 wrinkled green}.
\end{itemize}

Two of these classes (round green, wrinkled yellow) are \emph{new combinations} not seen in the original parents. This told Mendel that the allele for seed shape is inherited \emph{independently} of the allele for seed colour.

\begin{propriete}[Mendel's Second Law --- Law of Independent Assortment]
When two or more pairs of alleles are considered together, the alleles of \textbf{different genes assort independently of one another} during gamete formation, so that a gamete is equally likely to receive any combination of alleles.\\
\textbf{Cytological basis:} genes located on \emph{different} (non-homologous) chromosomes behave this way because, at \textbf{metaphase I of meiosis}, homologous chromosome pairs orient themselves at the equator \emph{independently} of one another (independent orientation of bivalents). A dihybrid $RrYy$ therefore produces four gamete types in equal proportions: $RY$, $Ry$, $rY$, $ry$ ($\tfrac{1}{4}$ each).\\
\emph{The statement and its cytological explanation are examinable at GCE Advanced Level.}
\end{propriete}

\courssub{The 16-Box Punnett Square and the Forked-Line Method}

Since each F$_1$ parent produces four gamete types, the F$_2$ Punnett square has $4 \times 4 = 16$ boxes.

\begin{popfigure}
\begin{center}
\begin{tikzpicture}[scale=0.85]
% Gamete labels
\foreach \x/\g in {0/RY, 1/Ry, 2/rY, 3/ry}{
  \node[font=\bfseries, text=popblue] at (\x+0.5, 4.3) {\g};
  \node[font=\bfseries, text=popblue] at (-0.6, 3.5-\x) {\g};
}
% Grid and genotypes
\foreach \x in {0,...,3}{
  \foreach \y in {0,...,3}{
    \draw[popline] (\x, \y) rectangle (\x+1, \y+1);
  }
}
% Genotypes (row 3 = top)
\node at (0.5,3.5) {\scriptsize RRYY}; \node at (1.5,3.5) {\scriptsize RRYy};
\node at (2.5,3.5) {\scriptsize RrYY}; \node at (3.5,3.5) {\scriptsize RrYy};
% Row 2
\node at (0.5,2.5) {\scriptsize RRYy}; \node at (1.5,2.5) {\scriptsize RRyy};
\node at (2.5,2.5) {\scriptsize RrYy}; \node at (3.5,2.5) {\scriptsize Rryy};
% Row 1
\node at (0.5,1.5) {\scriptsize RrYY}; \node at (1.5,1.5) {\scriptsize RrYy};
\node at (2.5,1.5) {\scriptsize rrYY}; \node at (3.5,1.5) {\scriptsize rrYy};
% Row 0
\node at (0.5,0.5) {\scriptsize RrYy}; \node at (1.5,0.5) {\scriptsize Rryy};
\node at (2.5,0.5) {\scriptsize rrYy}; \node at (3.5,0.5) {\scriptsize rryy};
% Highlight 9 round-yellow (R-Y-)
\foreach \coord in {(0,3),(1,3),(2,3),(3,3),(0,2),(2,2),(0,1),(1,1),(0,0)}{
  \fill[popgreenL, opacity=0.55] \coord rectangle +(1,1);
}
% Redraw grid on top
\foreach \x in {0,...,3}{
  \foreach \y in {0,...,3}{
    \draw[popline] (\x, \y) rectangle (\x+1, \y+1);
  }
}
% Redraw genotypes on top
\node at (0.5,3.5) {\scriptsize RRYY}; \node at (1.5,3.5) {\scriptsize RRYy};
\node at (2.5,3.5) {\scriptsize RrYY}; \node at (3.5,3.5) {\scriptsize RrYy};
\node at (0.5,2.5) {\scriptsize RRYy}; \node at (1.5,2.5) {\scriptsize RRyy};
\node at (2.5,2.5) {\scriptsize RrYy}; \node at (3.5,2.5) {\scriptsize Rryy};
\node at (0.5,1.5) {\scriptsize RrYY}; \node at (1.5,1.5) {\scriptsize RrYy};
\node at (2.5,1.5) {\scriptsize rrYY}; \node at (3.5,1.5) {\scriptsize rrYy};
\node at (0.5,0.5) {\scriptsize RrYy}; \node at (1.5,0.5) {\scriptsize Rryy};
\node at (2.5,0.5) {\scriptsize rrYy}; \node at (3.5,0.5) {\scriptsize rryy};
\end{tikzpicture}
\end{center}
\legende{Punnett square for $RrYy \times RrYy$. Shaded boxes (9 of 16) have at least one $R$ and one $Y$: round yellow. The remaining classes give 3 round green ($R\_yy$), 3 wrinkled yellow ($rrY\_$) and 1 wrinkled green ($rryy$): the 9:3:3:1 ratio.}
\end{popfigure}

Counting the genotypes: 9 boxes are $R\_Y\_$ (round yellow), 3 are $R\_yy$ (round green), 3 are $rrY\_$ (wrinkled yellow) and 1 is $rryy$ (wrinkled green).

\begin{methode}[The forked-line (branch) method --- faster than 16 boxes]
Treat each gene separately, then multiply the probabilities:
\begin{enumerate}
\item For shape, $Rr \times Rr$ gives $\tfrac{3}{4}$ round, $\tfrac{1}{4}$ wrinkled.
\item For colour, $Yy \times Yy$ gives $\tfrac{3}{4}$ yellow, $\tfrac{1}{4}$ green.
\item Multiply along the branches:
\[ \tfrac{3}{4}\text{ round} \times \tfrac{3}{4}\text{ yellow} = \tfrac{9}{16}\ \text{round yellow};\quad \tfrac{3}{4}\times\tfrac{1}{4} = \tfrac{3}{16}\ \text{round green}; \]
\[ \tfrac{1}{4}\times\tfrac{3}{4} = \tfrac{3}{16}\ \text{wrinkled yellow};\quad \tfrac{1}{4}\times\tfrac{1}{4} = \tfrac{1}{16}\ \text{wrinkled green}. \]
\end{enumerate}
This is exactly the 9:3:3:1 ratio, obtained in seconds. Multiplication is valid \emph{because} the genes assort independently.
\end{methode}

\begin{attention}[Common mistake]
The 9:3:3:1 ratio applies to the \textbf{F$_2$ of a cross between two heterozygotes} with \textbf{complete dominance at both loci} and \textbf{unlinked genes}. Do not quote it for a test cross, for codominance, or for linked genes.
\end{attention}

\courssub{The Dihybrid Test Cross}

A \cle{test cross} crosses the dihybrid with the double recessive homozygote: $RrYy \times rryy$. The recessive parent can only give $ry$ gametes, so the offspring phenotypes \emph{reveal directly} the gametes produced by the dihybrid. For unlinked genes the four gametes are equally frequent, giving:

\begin{center}
\begin{tabularx}{\linewidth}{|l|X|X|X|X|}
\hline
\rowcolor{popblueL} \textbf{Gamete from $RrYy$} & $RY$ & $Ry$ & $rY$ & $ry$ \\
\hline
Offspring genotype & $RrYy$ & $Rryy$ & $rrYy$ & $rryy$ \\
\hline
Phenotype & round yellow & round green & wrinkled yellow & wrinkled green \\
\hline
Proportion & $\tfrac{1}{4}$ & $\tfrac{1}{4}$ & $\tfrac{1}{4}$ & $\tfrac{1}{4}$ \\
\hline
\end{tabularx}
\end{center}

Hence the expected test cross ratio for unlinked genes is \textbf{1:1:1:1}.

\courssub{Gene Linkage: Morgan's Discovery}

Mendel was lucky: the genes he studied lay on different chromosomes. But chromosomes are few and genes are many, so most genes must share a chromosome with others.

\begin{definition}[Linkage]
\cle{Linked genes} are genes located on the \textbf{same chromosome}. Because they are physically attached, they tend to be \textbf{inherited together} and do \emph{not} assort independently. A set of linked genes is called a \cle{linkage group}; the number of linkage groups equals the haploid number of chromosomes (e.g.\ 4 in \emph{Drosophila}, 23 in humans).
\end{definition}

\textbf{Morgan's experiments with \emph{Drosophila melanogaster}.} Thomas Hunt Morgan crossed grey-bodied, normal-winged flies ($b^+b^+vg^+vg^+$) with black-bodied, vestigial-winged flies ($bb\,vgvg$), then test-crossed the F$_1$ females. Instead of 1:1:1:1, he obtained a large \textbf{excess of the two parental phenotypes} (about 41.5\% each) and only small numbers of the two new combinations (about 8.5\% each). Morgan concluded that the body-colour and wing-size genes are on the same chromosome --- linked --- and that the minority of new combinations arises by \cle{crossing over}.

\begin{attention}[Warning: linked genes break Mendelian ratios]
Linked genes do \textbf{not} give 9:3:3:1 in an F$_2$ nor 1:1:1:1 in a test cross. The diagnostic sign of linkage is an \textbf{excess of parental phenotypes} and a \textbf{shortage of recombinant phenotypes}. If your test cross gives, say, 42:42:8:8 instead of 25:25:25:25, suspect linkage immediately.
\end{attention}

\courssub{Crossing Over and Recombinant Gametes}

During \textbf{prophase I} of meiosis, homologous chromosomes pair up (synapsis) to form bivalents. At points of contact called \cle{chiasmata} (singular: chiasma), \textbf{non-sister chromatids} break at corresponding points and exchange segments. A chromatid that has exchanged a segment now carries a \emph{new combination} of alleles: it is a \cle{recombinant} chromatid.

\begin{definition}[Parental and recombinant offspring]
\textbf{Parental (non-recombinant) offspring} show the same combination of characters as the original parents. \textbf{Recombinant offspring} show new combinations of characters produced by crossing over between the linked genes.
\end{definition}

\begin{popfigure}
\begin{center}
\begin{tikzpicture}[scale=1.1]
% Before crossing over
\node[font=\bfseries] at (2.0,4.2) {Before crossing over (synapsis)};
% Homolog 1 (maternal) - blue
\draw[very thick, popblue] (0.5,0.5) -- (0.5,3.5);
\draw[very thick, popblue] (1.0,0.5) -- (1.0,3.5);
% Homolog 2 (paternal) - orange
\draw[very thick, poporange] (2.5,0.5) -- (2.5,3.5);
\draw[very thick, poporange] (3.0,0.5) -- (3.0,3.5);
% Alleles
\node[popblue, left] at (0.5,3.2) {\scriptsize A};
\node[popblue, left] at (0.5,1.5) {\scriptsize B};
\node[poporange, right] at (3.0,3.2) {\scriptsize a};
\node[poporange, right] at (3.0,1.5) {\scriptsize b};
% Chiasma
\draw[popdark, thick] (1.0,2.5) -- (2.5,2.0);
\draw[popdark, thick] (1.0,2.0) -- (2.5,2.5);
\node[popdark] at (1.75,2.25) {\scriptsize chiasma};
% After
\node[font=\bfseries] at (7.5,4.2) {After exchange (recombinant chromatids)};
% Four chromatids after meiosis
% Chromatid 1: A b (recombinant)
\draw[very thick, popblue] (5.5,0.5) -- (5.5,2.2);
\draw[very thick, poporange] (5.5,2.2) -- (5.5,3.5);
% Chromatid 2: A B (parental)
\draw[very thick, popblue] (6.0,0.5) -- (6.0,3.5);
% Chromatid 3: a b (parental)
\draw[very thick, poporange] (7.5,0.5) -- (7.5,3.5);
% Chromatid 4: a B (recombinant)
\draw[very thick, poporange] (8.0,0.5) -- (8.0,2.2);
\draw[very thick, popblue] (8.0,2.2) -- (8.0,3.5);
% Alleles after
\node[left] at (5.5,3.2) {\scriptsize A}; \node[left] at (5.5,1.5) {\scriptsize b};
\node[left] at (6.0,3.2) {\scriptsize A}; \node[left] at (6.0,1.5) {\scriptsize B};
\node[right] at (7.5,3.2) {\scriptsize a}; \node[right] at (7.5,1.5) {\scriptsize b};
\node[right] at (8.0,3.2) {\scriptsize a}; \node[right] at (8.0,1.5) {\scriptsize B};
% Gametes
\node[font=\bfseries] at (11.5,4.2) {Resulting gametes};
\node at (11.5,3.5) {\scriptsize $AB$ \ (parental)};
\node at (11.5,3.0) {\scriptsize $ab$ \ (parental)};
\node[text=popgreen] at (11.5,2.5) {\scriptsize $Ab$ \ (recombinant)};
\node[text=popgreen] at (11.5,2.0) {\scriptsize $aB$ \ (recombinant)};
\draw[->, thick, popmuted] (3.5,2.25) -- (5.0,2.25);
\draw[->, thick, popmuted] (8.5,2.25) -- (10.5,2.25);
\end{tikzpicture}
\end{center}
\legende{Crossing over at a chiasma between non-sister chromatids in prophase I. Only the two chromatids involved in the exchange become recombinant ($Ab$, $aB$); the other two remain parental ($AB$, $ab$).}
\end{popfigure}

Note that each single crossover event involves only \textbf{two of the four} chromatids, so even when crossing over occurs, at most half the gametes from that meiosis are recombinant. This is why recombinant offspring are always in the minority.

\courssub{Coupling and Repulsion Phase}

The arrangement of alleles on the homologous chromosomes of a dihybrid affects the phenotypic ratios. 

\begin{definition}[Coupling (cis) and repulsion (trans) phase]
\begin{itemize}
\item \textbf{Coupling phase (cis):} the two dominant alleles are on one homologue and the two recessives on the other, e.g.\ $AB/ab$. Parental gametes are $AB$ and $ab$.
\item \textbf{Repulsion phase (trans):} each homologue carries one dominant and one recessive allele, e.g.\ $Ab/aB$. Parental gametes are $Ab$ and $aB$.
\end{itemize}
\end{definition}

In a test cross, the \textbf{two most frequent offspring classes always reveal the parental gametes}, and therefore the phase of the dihybrid.

\courssub{Crossover Value and Chromosome Mapping}

\begin{definition}[Crossover value]
The \cle{crossover value} (COV), or \cle{recombination frequency}, is the percentage of recombinant offspring in the progeny of a test cross:
\[ \text{COV} = \frac{\text{number of recombinant offspring}}{\text{total number of offspring}} \times 100\% \]
A crossover value of 1\% corresponds to a map distance of \textbf{1 map unit} (or \textbf{1 centimorgan}, cM). The further apart two genes are on a chromosome, the more often a crossover occurs between them, so COV is used to estimate the \textbf{relative positions of genes} (chromosome mapping). Genes with COV near 50\% behave as if unlinked (they are either far apart on the same chromosome or on different chromosomes).
\end{definition}

\begin{propriete}[Properties of crossover values]
\begin{itemize}
\item COV ranges from 0\% (complete linkage, no crossing over) to a maximum of 50\% (independent assortment).
\item For genes far apart, multiple crossovers can occur; the observed COV underestimates the true map distance. Mapping functions (e.g.\ Haldane, Kosambi) correct for this, but at GCE level we use COV directly as map distance for values $\le 30\%$.
\item The sum of COVs for adjacent intervals approximates the COV for the outer genes (additivity of map distances for short intervals).
\end{itemize}
\end{propriete}

\begin{methode}[Calculating the crossover value from test cross results]
\begin{enumerate}
\item Identify the \textbf{parental classes}: the two offspring classes with the \emph{largest} numbers (they match the original parents' phenotypes).
\item Identify the \textbf{recombinant classes}: the two classes with the \emph{smallest} numbers.
\item Add the recombinant numbers; add all four classes for the total.
\item Apply the formula: $\text{COV} = \dfrac{\text{recombinants}}{\text{total}} \times 100$.
\item State the map distance: COV in \% = distance in map units (cM).
\end{enumerate}
\end{methode}

\begin{exoresolu}[Worked example: crossover value in maize]
In maize, coloured full seeds ($CS$) are dominant to colourless shrunken seeds ($cs$). A dihybrid $CcSs$ (gametes in coupling: $CS/cs$) was test-crossed with $ccss$. The offspring were: coloured full = 4032; colourless shrunken = 3968; coloured shrunken = 152; colourless full = 149. Calculate the crossover value and the map distance between the genes.
\tcblower
\textbf{Step 1 --- Parental classes} (most frequent): coloured full (4032) and colourless shrunken (3968); these match the parental combinations $CS$ and $cs$.\\
\textbf{Step 2 --- Recombinant classes} (least frequent): coloured shrunken (152) and colourless full (149), from gametes $Cs$ and $cS$.\\
\textbf{Step 3 --- Totals:}
\[ \text{recombinants} = 152 + 149 = 301; \qquad \text{total} = 4032 + 3968 + 152 + 149 = 8301. \]
\textbf{Step 4 --- Crossover value:}
\[ \text{COV} = \frac{301}{8301} \times 100 \approx 3.6\%. \]
\textbf{Step 5 --- Map distance:} the two genes are approximately \textbf{3.6 map units (cM)} apart on the same chromosome. The strong deviation from 1:1:1:1 (which would predict about 2075 per class) confirms linkage.
\end{exoresolu}

\begin{popfigure}
\begin{center}
\begin{tikzpicture}
\begin{axis}[
  ybar, bar width=14pt,
  width=0.9\linewidth, height=6.5cm,
  symbolic x coords={Coloured full, Colourless shrunken, Coloured shrunken, Colourless full},
  xtick=data, x tick label style={rotate=20, anchor=east, font=\scriptsize},
  ylabel={Number of offspring}, ymin=0, ymax=4600,
  legend style={at={(0.5,1.02)}, anchor=south, legend columns=2, font=\scriptsize},
  ymajorgrids, grid style={popline!40},
]
\addplot[fill=popblue] coordinates {(Coloured full,4032) (Colourless shrunken,3968) (Coloured shrunken,152) (Colourless full,149)};
\addplot[fill=poporange] coordinates {(Coloured full,2075) (Colourless shrunken,2075) (Coloured shrunken,2075) (Colourless full,2075)};
\legend{Observed (linked), Expected 1:1:1:1 (unlinked)}
\end{axis}
\end{tikzpicture}
\end{center}
\legende{Observed offspring from the maize test cross compared with the 1:1:1:1 expectation for unlinked genes. The huge excess of the two parental classes is the signature of linkage.}
\end{popfigure}

\begin{attention}[Exam tip]
In any test cross with linked genes, \textbf{the two most frequent offspring classes are the parental (non-recombinant) types}, and they tell you the arrangement of alleles on the dihybrid's chromosomes (coupling $AB/ab$ or repulsion $Ab/aB$). The two rare classes are always the recombinants. Never average all four classes when computing a crossover value --- only recombinants go in the numerator.
\end{attention}

\begin{savaistu}[Chi-square test for linkage]
At GCE Advanced Level, you may be asked to use a $\chi^2$ test to decide whether observed ratios deviate significantly from the expected 1:1:1:1 (unlinked) or from a hypothesised linkage ratio. The null hypothesis is usually \emph{independent assortment} (1:1:1:1). Degrees of freedom = number of classes $-$ 1 = 3. If $\chi^2_{\text{calc}} > \chi^2_{\text{crit}}$ (7.82 at $p=0.05$), reject independent assortment and conclude linkage.
\end{savaistu}

\begin{exercice}[Test your understanding]
A heterozygous female \emph{Drosophila} with grey body and normal wings was test-crossed. The offspring were: grey, normal = 965; black, vestigial = 944; grey, vestigial = 206; black, normal = 185.
\begin{enumerate}
\item State the expected ratio if the genes were unlinked.
\item Identify the parental and recombinant classes.
\item Calculate the crossover value and the map distance between the genes.
\item State the phase (coupling or repulsion) of the dihybrid female.
\end{enumerate}
\end{exercice}

\begin{corrige}[Exercise 1]
\begin{enumerate}
\item Unlinked genes would give a 1:1:1:1 ratio (about 575 offspring per class).
\item Parental (most frequent): grey normal (965) and black vestigial (944). Recombinant (least frequent): grey vestigial (206) and black normal (185).
\item Recombinants $= 206 + 185 = 391$; total $= 965 + 944 + 206 + 185 = 2300$.
\[ \text{COV} = \frac{391}{2300} \times 100 = 17\%. \]
The genes are linked, approximately \textbf{17 map units (cM)} apart.
\item The parental gametes are $b^+vg^+$ (grey normal) and $b\,vg$ (black vestigial). Since both dominant alleles are on one chromosome and both recessives on the other, the dihybrid was in \textbf{coupling phase} ($b^+vg^+/b\,vg$).
\end{enumerate}
\end{corrige}

\begin{aretenir}[Key points]
\begin{itemize}
\item A dihybrid cross $RrYy \times RrYy$ between unlinked genes gives F$_2$ ratio \textbf{9:3:3:1}; the test cross $RrYy \times rryy$ gives \textbf{1:1:1:1}.
\item Mendel's Second Law: alleles of different genes assort independently, because bivalents orient independently at metaphase I.
\item \textbf{Linked genes} lie on the same chromosome and are inherited together, giving an excess of parental phenotypes.
\item \textbf{Crossing over} at chiasmata (prophase I) between non-sister chromatids produces recombinant gametes.
\item $\text{COV} = \dfrac{\text{recombinants}}{\text{total}} \times 100$; 1\% = 1 map unit (cM); COV measures the distance between genes.
\item Maximum COV = 50\% (independent assortment). COV $\le 30\%$ is used directly as map distance.
\item Parental classes = most frequent; recombinant classes = least frequent. They reveal coupling or repulsion phase.
\end{itemize}
\end{aretenir}

\coursec{Gene Interactions: Epistasis, Lethal Genes and Multiple Alleles}

In the previous section we studied dihybrid inheritance, where two genes at two different loci assort independently and produce the classical $9:3:3:1$ phenotypic ratio in the $F_2$ generation. However, genes do not always act in isolation. Often the product of one gene is required for the expression of another, so one gene can \cle{mask} or modify the effect of a different gene. Sometimes an allele disrupts a vital function so severely that individuals carrying a particular genotype do not survive. And sometimes a gene exists in a population in three, four or more allelic forms, although any single diploid individual still carries only two. In this section we examine these three phenomena — \cle{epistasis}, \cle{lethal genes} and \cle{multiple alleles} — and learn to recognise their characteristic, non-Mendelian ratios in examination data.

\courssub{Epistasis: When One Gene Masks Another}

\textbf{Intuition.} Imagine a painter who must first apply a white primer before any colour can be seen on a wall. If the primer is missing, it does not matter which colour of paint is chosen — the wall stays unpainted. In the same way, many traits depend on a \emph{chain} of gene products: a pigment can only be deposited if a precursor substance has first been produced by another gene. The gene acting "upstream" can therefore hide the effect of the gene acting "downstream". This is epistasis.

\begin{definition}[Epistasis]
\cle{Epistasis} is a form of gene interaction in which an allele (or genotype) at one locus \textbf{masks, suppresses or modifies the expression} of alleles at a \emph{different} locus. The masking gene is called the \textbf{epistatic gene}; the gene whose expression is masked is called the \textbf{hypostatic gene}. Epistasis is \emph{not} dominance: dominance is an interaction between alleles of the \emph{same} gene, whereas epistasis is an interaction between \emph{different} genes.
\end{definition}

\begin{propriete}[Modified dihybrid ratios]
When two independently assorting genes interact epistatically, the $F_2$ generation still contains 16 equally likely zygote combinations, but several genotype classes now share the \emph{same} phenotype. The classical $9:3:3:1$ ratio is therefore modified into characteristic ratios:
\begin{itemize}
\item $9:3:4$ — recessive epistasis;
\item $12:3:1$ — dominant epistasis;
\item $9:7$ — complementary gene action (both dominant alleles needed);
\item $15:1$ — duplicate gene action (either dominant allele sufficient).
\end{itemize}
You are not required to prove these ratios from first principles; you must be able to \emph{derive} them by grouping the 16 boxes of a Punnett square, and to \emph{recognise} them in exam data.
\end{propriete}

\textbf{Recessive epistasis: coat colour in mice (and Labrador retrievers).} In mice, gene $B/b$ determines pigment colour ($B$ = black, dominant; $b$ = brown, recessive), but pigment can only be deposited in the fur if at least one dominant allele $C$ of a second gene is present. Genotype $cc$ blocks pigment deposition entirely, giving an \textbf{albino} (white) mouse whatever the $B/b$ genotype is. Thus $cc$ is epistatic to $B/b$. The same logic explains coat colour in Labrador retrievers, where $ee$ gives golden dogs regardless of the black/chocolate gene.

Crossing two double heterozygotes $BbCc \times BbCc$ gives the usual 16 combinations, which group as follows:

\begin{itemize}
\item $9\ B\_C\_$ : black (pigment deposited, black allele present);
\item $3\ bbC\_$ : brown (pigment deposited, only brown alleles);
\item $3\ B\_cc + 1\ bbcc = 4$ : albino (no pigment deposited).
\end{itemize}

\begin{popfigure}
\begin{center}
\begin{tikzpicture}[scale=0.85, every node/.style={font=\small}]
% Grid
\draw[thick, popblue] (0,0) grid (4,4);
% Gamete labels
\foreach \x/\g in {0.5/BC, 1.5/Bc, 2.5/bC, 3.5/bc}
  \node[above, font=\small\bfseries, text=popdark] at (\x,4) {\g};
\foreach \y/\g in {3.5/BC, 2.5/Bc, 1.5/bC, 0.5/bc}
  \node[left, font=\small\bfseries, text=popdark] at (0,\y) {\g};
% Genotypes
\node at (0.5,3.5) {BBCC}; \node at (1.5,3.5) {BBCc}; \node at (2.5,3.5) {BbCC}; \node at (3.5,3.5) {BbCc};
\node at (0.5,2.5) {BBCc}; \node at (1.5,2.5) {BBcc}; \node at (2.5,2.5) {BbCc}; \node at (3.5,2.5) {Bbcc};
\node at (0.5,1.5) {BbCC}; \node at (1.5,1.5) {BbCc}; \node at (2.5,1.5) {bbCC}; \node at (3.5,1.5) {bbCc};
\node at (0.5,0.5) {BbCc}; \node at (1.5,0.5) {Bbcc}; \node at (2.5,0.5) {bbCc}; \node at (3.5,0.5) {bbcc};
% Phenotype fills: black (9) = popgreenL, brown (3) = popblueL, albino (4) = popgold
\fill[popgreenL, opacity=0.6] (0,3) rectangle (1,4);   % BBCC
\fill[popgreenL, opacity=0.6] (1,3) rectangle (2,4);   % BBCc
\fill[popgreenL, opacity=0.6] (2,3) rectangle (3,4);   % BbCC
\fill[popgreenL, opacity=0.6] (3,3) rectangle (4,4);   % BbCc
\fill[popgreenL, opacity=0.6] (0,2) rectangle (1,3);   % BBCc
\fill[popgreenL, opacity=0.6] (2,2) rectangle (3,3);   % BbCc
\fill[popgreenL, opacity=0.6] (0,1) rectangle (1,2);   % BbCC
\fill[popgreenL, opacity=0.6] (1,1) rectangle (2,2);   % BbCc
\fill[popgreenL, opacity=0.6] (2,1) rectangle (3,2);   % bbCC
\fill[popblueL, opacity=0.7] (3,1) rectangle (4,2);    % bbCc
\fill[popblueL, opacity=0.7] (2,0) rectangle (3,1);    % bbCc
\fill[popblueL, opacity=0.7] (3,0) rectangle (4,1);    % bbcc? wait bbcc is albino
% Correction: bbcc is albino (gold). Let's redo fills systematically.
% Actually, better to use a loop but manual is fine for 16 cells.
% Let's list all 16 with coordinates (x,y) where x=0..3, y=0..3 (bottom-left origin)
% Row y=3 (top): BC x BC, Bc x BC, bC x BC, bc x BC -> BBCC, BBCc, BbCC, BbCc -> all black
% Row y=2: BC x Bc, Bc x Bc, bC x Bc, bc x Bc -> BBCc, BBcc, BbCc, Bbcc -> black, albino, black, albino
% Row y=1: BC x bC, Bc x bC, bC x bC, bc x bC -> BbCC, BbCc, bbCC, bbCc -> black, black, brown, brown
% Row y=0: BC x bc, Bc x bc, bC x bc, bc x bc -> BbCc, Bbcc, bbCc, bbcc -> black, albino, brown, albino
% So black (9): (0,3),(1,3),(2,3),(3,3), (0,2),(2,2), (0,1),(1,1),(2,1) -> 9
% Brown (3): (3,1),(2,0),(3,0) -> 3
% Albino (4): (1,2),(3,2), (1,0),(3,0)? wait (3,0) is bbcc albino, (1,0) is Bbcc albino, (1,2) BBcc albino, (3,2) Bbcc albino -> 4
% Let's fill correctly:
\fill[popgreenL, opacity=0.6] (0,3) rectangle (1,4);
\fill[popgreenL, opacity=0.6] (1,3) rectangle (2,4);
\fill[popgreenL, opacity=0.6] (2,3) rectangle (3,4);
\fill[popgreenL, opacity=0.6] (3,3) rectangle (4,4);
\fill[popgreenL, opacity=0.6] (0,2) rectangle (1,3);
\fill[popgreenL, opacity=0.6] (2,2) rectangle (3,3);
\fill[popgreenL, opacity=0.6] (0,1) rectangle (1,2);
\fill[popgreenL, opacity=0.6] (1,1) rectangle (2,2);
\fill[popgreenL, opacity=0.6] (2,1) rectangle (3,2);
\fill[popblueL, opacity=0.7] (3,1) rectangle (4,2);
\fill[popblueL, opacity=0.7] (2,0) rectangle (3,1);
\fill[popblueL, opacity=0.7] (3,0) rectangle (4,1); % bbcc is albino! mistake. bbcc should be gold.
% Actually (3,0) is bbcc -> albino. So brown only (3,1) and (2,0). That's 2? Wait count:
% Genotypes for brown: bbC_ -> bbCC, bbCc, bbCc. That's three: (2,1)=bbCC, (3,1)=bbCc, (2,0)=bbCc. (3,0)=bbcc albino.
% So brown: (2,1), (3,1), (2,0). Albino: (1,2), (3,2), (1,0), (3,0).
\fill[popblueL, opacity=0.7] (2,1) rectangle (3,2); % bbCC
\fill[popblueL, opacity=0.7] (3,1) rectangle (4,2); % bbCc
\fill[popblueL, opacity=0.7] (2,0) rectangle (3,1); % bbCc
\fill[popgold, opacity=0.5] (1,2) rectangle (2,3); % BBcc
\fill[popgold, opacity=0.5] (3,2) rectangle (4,3); % Bbcc
\fill[popgold, opacity=0.5] (1,0) rectangle (2,1); % Bbcc
\fill[popgold, opacity=0.5] (3,0) rectangle (4,1); % bbcc
% Redraw grid on top
\draw[thick, popblue] (0,0) grid (4,4);
% Legend
\node[font=\footnotesize, text=popdark, anchor=west] at (4.2,3.5) {green = black (9)};
\node[font=\footnotesize, text=popdark, anchor=west] at (4.2,3.0) {blue = brown (3)};
\node[font=\footnotesize, text=popdark, anchor=west] at (4.2,2.5) {gold = albino (4)};
\end{tikzpicture}
\end{center}
\legende{Punnett square for $BbCc \times BbCc$ (recessive epistasis). All $cc$ genotypes are albino, giving the modified ratio $9$ black : $3$ brown : $4$ albino.}
\end{popfigure}

\textbf{Dominant epistasis: fruit colour in summer squash.} Here a \emph{dominant} allele at one locus does the masking. In summer squash, the dominant allele $W$ inhibits colour development (white fruit), while $ww$ allows colour; the second gene then decides: $Y$ = yellow (dominant), $y$ = green. A cross $WwYy \times WwYy$ gives:

\begin{itemize}
\item $9\ W\_Y\_ + 3\ W\_yy = 12$ : white ($W$ masks everything);
\item $3\ wwY\_$ : yellow;
\item $1\ wwyy$ : green.
\end{itemize}

Hence the ratio $12:3:1$. The same pattern occurs in poultry plumage, where the dominant allele $I$ (inhibitor) produces white birds regardless of the colour gene.

\textbf{Extensions: complementary and duplicate genes.} In \emph{complementary gene} action (e.g. flower colour in sweet pea), both dominant alleles $C$ \emph{and} $P$ are needed to produce pigment, so only the $9\ C\_P\_$ class is coloured and all other classes ($3+3+1=7$) are white: ratio $9:7$. In \emph{duplicate gene} action, either dominant allele alone produces the same trait, so only the double recessive ($1/16$) differs: ratio $15:1$.

\begin{methode}[Recognising and solving epistasis problems]
\begin{enumerate}
\item Count the total offspring. If the $F_2$ total fits sixteenths (e.g. 160, 320), suspect a dihybrid-type cross.
\item Group the phenotypes: divide each class by the smallest to get a simple integer ratio. Compare with $9:3:4$, $12:3:1$, $9:7$, $15:1$.
\item Identify which genotype classes share the same phenotype — this reveals which gene is epistatic and whether the epistatic allele is dominant or recessive.
\item Group the 16 Punnett boxes accordingly and state the modified ratio with the phenotype of each class.
\end{enumerate}
\end{methode}

\begin{attention}[Exam tip: group before you conclude]
In epistasis questions, \textbf{first identify which genotype classes share the same phenotype, then group the 16 boxes accordingly}. A very common error is to write "$9:3:3:1$" automatically for any dihybrid cross, or to force data into $9:3:3:1$ when they clearly fit $9:3:4$. Always let the \emph{numbers} tell you the interaction, not habit.
\end{attention}

\begin{exoresolu}[Recognising recessive epistasis]
A geneticist crosses pure-breeding black mice ($BBCC$) with albino mice ($bbcc$). The $F_1$ are all black. Interbreeding the $F_1$ gives 320 $F_2$ offspring: 180 black, 60 brown and 80 albino. Explain these results.
\tcblower
\textbf{Solution.} The $F_1$ are $BbCc$ (black). Dividing the $F_2$ classes by 20: $180:60:80 = 9:3:4$. This is the signature of \textbf{recessive epistasis}: the 16-box Punnett square gives $9\ B\_C\_$ black, $3\ bbC\_$ brown, and the $cc$ classes ($3\ B\_cc + 1\ bbcc = 4$) albino because $cc$ prevents pigment deposition. The gene $C/c$ is epistatic to $B/b$, and the epistatic allele $c$ acts recessively. Expected numbers: $\tfrac{9}{16}\times 320 = 180$, $\tfrac{3}{16}\times 320 = 60$, $\tfrac{4}{16}\times 320 = 80$ — a perfect fit.
\end{exoresolu}

\begin{savaistu}[Cameroon context: coat colour in local livestock]
In Cameroon's indigenous cattle breeds (e.g., Kapsiki, Namchi), coat colour patterns often involve epistatic interactions. For instance, the \emph{white spotting} gene can be epistatic to base colour genes, producing pied patterns. Understanding such interactions helps animal breeders select for desired coat colours which may be linked to heat tolerance or disease resistance.
\end{savaistu}

\courssub{Lethal Genes}

\textbf{Intuition.} Some alleles disrupt a function so essential that the organism cannot develop or survive. Such alleles are usually detected only because they distort the expected Mendelian ratios: a whole genotype class is missing from the offspring because those individuals die before birth (or before counting).

\begin{definition}[Lethal gene]
A \cle{lethal gene} (lethal allele) is an allele whose presence in a particular genotype causes the death of the organism, either before birth or later in life. A \textbf{recessive lethal} kills only when homozygous; a \textbf{dominant lethal} kills even in heterozygotes (and so can only persist in a population if it acts \emph{late}, after reproductive age).
\end{definition}

\textbf{Recessive lethals and the 2:1 ratio.} The classic example is yellow coat colour in mice. The allele $Y$ (yellow) is dominant to $y$ (grey/agouti) for coat colour, but homozygous $YY$ embryos die in the uterus. Every yellow mouse is therefore an obligate heterozygote $Yy$. Crossing two yellow mice:

\[ Yy \times Yy \;\longrightarrow\; 1\ \underbrace{YY}_{\text{lethal}} : 2\ Yy\ (\text{yellow}) : 1\ yy\ (\text{grey}) \]

Among the \emph{surviving} offspring the ratio is $2$ yellow : $1$ grey — not the expected $3:1$. The same pattern occurs in the \textbf{creeper fowl}, where heterozygotes have shortened legs but homozygous dominants die.

\begin{popfigure}
\begin{center}
\begin{tikzpicture}[scale=1.0, every node/.style={font=\small}]
\draw[thick, popblue] (0,0) grid (2,2);
\node[above, font=\small\bfseries, text=popdark] at (0.5,2) {$Y$};
\node[above, font=\small\bfseries, text=popdark] at (1.5,2) {$y$};
\node[left, font=\small\bfseries, text=popdark] at (0,1.5) {$Y$};
\node[left, font=\small\bfseries, text=popdark] at (0,0.5) {$y$};
\node at (0.5,1.5) {$YY$}; \node at (1.5,1.5) {$Yy$};
\node at (0.5,0.5) {$Yy$}; \node at (1.5,0.5) {$yy$};
\draw[thick, poppink, line width=1.2pt] (0.08,1.08) -- (0.92,1.92);
\draw[thick, poppink, line width=1.2pt] (0.08,1.92) -- (0.92,1.08);
\node[font=\footnotesize, text=poppink] at (0.5,1.15) {dies};
\fill[popgreenL, opacity=0.7] (1,1) rectangle (2,2); % Yy
\fill[popgreenL, opacity=0.7] (0,0) rectangle (1,1); % Yy
\fill[popblueL, opacity=0.7] (1,0) rectangle (2,1); % yy
\draw[thick, popblue] (0,0) grid (2,2);
\node[font=\footnotesize, text=popdark, anchor=west] at (2.2,1.5) {green = yellow $Yy$ (2)};
\node[font=\footnotesize, text=popdark, anchor=west] at (2.2,1.0) {blue = grey $yy$ (1)};
\node[font=\footnotesize, text=poppink, anchor=west] at (2.2,0.5) {crossed = $YY$ lethal};
\end{tikzpicture}
\end{center}
\legende{Cross $Yy \times Yy$ in yellow mice. The $YY$ class dies, so surviving offspring show a $2:1$ ratio instead of $3:1$.}
\end{popfigure}

\textbf{Dominant lethals.} A dominant lethal allele kills even in single dose, so it normally eliminates itself from the population. It can only be transmitted if its effects appear \emph{after} reproductive age. The standard example is \textbf{Huntington's disease} in humans: affected heterozygotes are healthy for decades, have children (passing the allele to about half of them), and only later develop the fatal neurological degeneration.

\begin{attention}[Do not "correct" a 2:1 ratio]
A $2:1$ ratio instead of $3:1$ in a monohybrid cross strongly suggests a \textbf{lethal homozygous genotype}. Never "correct" the data to $3:1$ and never assume the counting was wrong. Instead, state that one genotype class (usually the homozygous dominant) is missing because it is lethal, and support this by noting that all surviving individuals showing the dominant trait are heterozygous (they never breed true).
\end{attention}

\begin{exoresolu}[Yellow mice]
Two yellow mice are crossed repeatedly and produce 48 surviving offspring: 31 yellow and 17 grey. Explain the genetic basis, and predict the result of crossing a yellow mouse with a grey mouse.
\tcblower
\textbf{Solution.} The ratio $31:17 \approx 2:1$ indicates a recessive lethal. Let $Y$ = yellow (dominant for colour, lethal when homozygous) and $y$ = grey. All yellow mice are $Yy$. The cross $Yy \times Yy$ gives $1\,YY$ (dies) : $2\,Yy$ (yellow) : $1\,yy$ (grey), i.e. $2:1$ among survivors, matching the data. A test cross $Yy \times yy$ gives $1\,Yy$ (yellow) : $1\,yy$ (grey), a $1:1$ ratio with no deaths, because no $YY$ zygote can form.
\end{exoresolu}

\courssub{Multiple Alleles and the ABO Blood Group System}

\textbf{Intuition.} Mendel worked with genes having two alleles (e.g. $T$ and $t$). But a gene can mutate in many different ways, so within a whole \emph{population} a gene may exist in three, four or more allelic forms. Any single diploid individual, however, still carries only \emph{two} of them — one on each homologous chromosome.

\begin{definition}[Multiple alleles]
\cle{Multiple alleles} are three or more alternative forms (alleles) of the same gene, occupying the same locus on homologous chromosomes, found within a population. Each individual carries at most two of these alleles. Multiple alleles typically form a \textbf{dominance hierarchy} or show codominance between some members.
\end{definition}

\textbf{The ABO blood group system.} The human ABO gene (locus $I$) has three alleles: $I^A$, $I^B$ and $i$. $I^A$ and $I^B$ are \textbf{codominant} (both expressed together), and both are dominant over $i$, which is recessive. The alleles determine antigens on the red blood cell surface, and the plasma naturally contains antibodies against the antigens the person lacks:

\begin{center}
\begin{tabularx}{\linewidth}{|l|l|X|X|}
\hline
\rowcolor{popblueL} \textbf{Phenotype (group)} & \textbf{Genotypes} & \textbf{Antigens on red cells} & \textbf{Antibodies in plasma} \\
\hline
A & $I^A I^A$ or $I^A i$ & A & anti-B \\
\hline
B & $I^B I^B$ or $I^B i$ & B & anti-A \\
\hline
AB & $I^A I^B$ & A and B & neither \\
\hline
O & $ii$ & neither & anti-A and anti-B \\
\hline
\end{tabularx}
\end{center}

This has direct medical importance in Cameroonian hospitals as everywhere else: a transfusion is safe only if the donor's red-cell antigens are not attacked by the recipient's antibodies. Group O individuals (no antigens) are \textbf{universal donors}; group AB individuals (no antibodies) are \textbf{universal recipients}.

\begin{center}
\begin{tabularx}{\linewidth}{|l|X|X|X|X|}
\hline
\rowcolor{popblueL} \textbf{Recipient} & \textbf{Can receive from} & \textbf{Can donate to} & \textbf{Antigens} & \textbf{Antibodies} \\
\hline
O & O & O, A, B, AB & none & anti-A, anti-B \\
\hline
A & A, O & A, AB & A & anti-B \\
\hline
B & B, O & B, AB & B & anti-A \\
\hline
AB & A, B, AB, O & AB & A, B & none \\
\hline
\end{tabularx}
\end{center}

\begin{methode}[Solving ABO problems]
\begin{enumerate}
\item Write down \textbf{all possible genotypes} for each phenotype involved (group A = $I^A I^A$ \emph{or} $I^A i$, etc.). Never forget the heterozygous options.
\item Choose the genotype combination consistent with the information given (e.g. a group O child proves each parent carries $i$).
\item Draw the Punnett square and read off the possible offspring phenotypes.
\item For paternity/dispute questions: a man is \emph{excluded} only if the child's phenotype is impossible from his possible gametes; a non-excluded man is merely a \emph{candidate}, not proven to be the father.
\end{enumerate}
\end{methode}

\begin{exoresolu}[Paternity dispute using blood groups]
A woman of blood group A has a child of blood group O. A man of blood group B is accused of being the father. Could he be? A second man is group AB. Could he be?
\tcblower
\textbf{Solution.} The child is $ii$, so it received $i$ from each parent. The mother (group A) must therefore be $I^A i$. The first man (group B) could be $I^B I^B$ or $I^B i$; if he is $I^B i$, the cross $I^A i \times I^B i$ gives $I^A I^B$, $I^A i$, $I^B i$, $ii$ — a group O child is possible, so he \textbf{cannot be excluded}. The second man (group AB) is $I^A I^B$ and can only pass on $I^A$ or $I^B$; all his children must be A, B or AB. A group O child is impossible, so he is \textbf{excluded} from paternity.
\end{exoresolu}

\textbf{Extension: the Bombay phenotype.} Rare individuals possess the alleles $I^A$ or $I^B$ but cannot express them because they are homozygous recessive ($hh$) at a second locus needed to build the precursor substance $H$ on which A and B antigens are made. They test as blood group O despite carrying A or B alleles — a beautiful real-life case of recessive epistasis acting on a multiple-allele system. This is beyond the core syllabus but may appear in extension questions.

\textbf{Another example: coat colour in rabbits.} The coat colour gene $C$ exists as four alleles with a strict dominance hierarchy:
\[ C \;(\text{full colour}) \;>\; c^{ch} \;(\text{chinchilla}) \;>\; c^{h} \;(\text{Himalayan}) \;>\; c \;(\text{albino}) \]
Each allele is dominant to all alleles below it and recessive to all above. Thus a Himalayan rabbit can be $c^h c^h$ or $c^h c$, but never carries $C$ or $c^{ch}$. This hierarchy explains why, for example, a cross between a full-colour rabbit ($Cc$) and a Himalayan rabbit ($c^h c$) can produce full-colour, Himalayan and albino offspring, but never chinchilla.

\begin{aretenir}[Key points of this section]
\begin{itemize}
\item \textbf{Epistasis}: one gene masks another gene at a different locus. Recessive epistasis $\rightarrow 9:3:4$; dominant epistasis $\rightarrow 12:3:1$; complementary genes $\rightarrow 9:7$; duplicate genes $\rightarrow 15:1$. All are groupings of the 16-box dihybrid square.
\item \textbf{Lethal alleles}: a missing genotype class distorts ratios. A $2:1$ ratio signals a homozygous lethal (yellow mice, creeper fowl); dominant lethals survive only if late-acting (Huntington's disease).
\item \textbf{Multiple alleles}: more than two alleles per gene in the population, two per individual. ABO system: $I^A$ and $I^B$ codominant, $i$ recessive; group O = universal donor, AB = universal recipient. Rabbit coat colour shows a dominance hierarchy $C > c^{ch} > c^h > c$.
\end{itemize}
\end{aretenir}

\begin{exercice}[Practice]
\begin{enumerate}
\item In summer squash, a white-fruited plant of genotype $WwYy$ is self-pollinated. State the expected phenotypic ratio and name the type of gene interaction.
\item Two yellow mice produce 60 offspring over several litters. How many of each colour do you expect among the survivors, and why is no pure-breeding yellow strain possible?
\item A child is blood group AB. State, with reasons, which blood groups the father \emph{cannot} have if the mother is group O.
\item A full-colour rabbit of genotype $Cc^h$ is crossed with a chinchilla rabbit of genotype $c^{ch}c$. List the expected phenotypes and their proportions.
\end{enumerate}
\end{exercice}

\begin{corrige}[Exercise 5.1]
\begin{enumerate}
\item $12$ white : $3$ yellow : $1$ green; dominant epistasis ($W$ masks the $Y/y$ gene).
\item $Yy \times Yy \rightarrow 1\,YY$ (dies) : $2\,Yy$ yellow : $1\,yy$ grey. Expected survivors: $40$ yellow, $20$ grey (the $YY$ quarter, about 15 embryos, is lost). No pure-breeding yellow strain is possible because $YY$ is lethal, so every yellow mouse is heterozygous.
\item The mother is $ii$, so the child received $I^A$ or $I^B$ from the father. The father cannot be group O ($ii$, gives only $i$). He could be A, B or AB.
\item Gametes: $C, c^h$ from the first parent; $c^{ch}, c$ from the second. Offspring genotypes: $Cc^{ch}$ (full colour), $Cc$ (full colour), $c^h c^{ch}$ (chinchilla, since $c^{ch} > c^h$), $c^h c$ (Himalayan). Phenotypic ratio: $\tfrac{1}{2}$ full colour : $\tfrac{1}{4}$ chinchilla : $\tfrac{1}{4}$ Himalayan.
\end{enumerate}
\end{corrige}

\coursec{Population Genetics: Gene Frequencies and the Hardy-Weinberg Equation}

So far in this chapter we have studied inheritance at the level of individuals and families: one gene, two parents, a Punnett square. But genes do not exist only in individuals --- they exist in \emph{populations}. A population of cocoa trees on the slopes of Mount Cameroon, a herd of cattle in the Adamaoua, or the human population of Yaound\'e each carries a vast collection of alleles, and the proportions of these alleles can change over time. That change \emph{is} evolution at its most fundamental level. Population genetics gives us the mathematical tools to describe the genetic make-up of a population and to predict whether it is evolving. The cornerstone of the whole subject is the \cle{Hardy-Weinberg principle}, a deceptively simple pair of equations with enormous power --- from estimating how many people in a village carry the sickle cell allele to managing the genetic diversity of endangered species.

\courssub{Populations, Gene Pools and Frequencies}

\begin{definition}[Population and gene pool]
A \cle{population} is a group of individuals of the same species living in the same area at the same time, capable of interbreeding. The \cle{gene pool} of a population is the total collection of all the alleles of all the genes present in that population.
\end{definition}

Think of the gene pool as a common reservoir. Each diploid individual draws two alleles per locus from it, and returns its own alleles to it through reproduction. To describe a gene pool numerically we use frequencies.

\begin{definition}[Allele frequency and genotype frequency]
The \cle{allele frequency} (or gene frequency) of an allele is the proportion of all copies of that gene in the population that are that particular allele. The \cle{genotype frequency} is the proportion of individuals in the population having a particular genotype.
\end{definition}

\begin{exemplebox}[Counting alleles in a small population]
Consider a population of $500$ diploid individuals at a locus with two alleles, $A$ and $a$. Suppose we count: $320$ individuals $AA$, $160$ individuals $Aa$, and $20$ individuals $aa$. Each individual carries $2$ alleles, so the gene pool contains $500 \times 2 = 1000$ alleles.
\begin{itemize}
\item Number of $A$ alleles $= (2 \times 320) + 160 = 800$, so frequency of $A$ is $p = \dfrac{800}{1000} = 0.8$.
\item Number of $a$ alleles $= (2 \times 20) + 160 = 200$, so frequency of $a$ is $q = \dfrac{200}{1000} = 0.2$.
\end{itemize}
Check: $p + q = 0.8 + 0.2 = 1$. The genotype frequencies are $\dfrac{320}{500} = 0.64$ for $AA$, $\dfrac{160}{500} = 0.32$ for $Aa$, and $\dfrac{20}{500} = 0.04$ for $aa$, which also sum to $1$.
\end{exemplebox}

\begin{attention}[Alleles versus individuals]
Each \emph{diploid} individual contributes \textbf{two} alleles to the gene pool. When calculating allele frequencies from genotype counts, always multiply the homozygote counts by $2$ and divide by twice the number of individuals. Forgetting the factor of $2$ is a classic source of error.
\end{attention}

\courssub{The Hardy-Weinberg Principle}

Here is the puzzle that G. H. Hardy and W. Weinberg independently solved in 1908. Dominant alleles do not ``take over'' a population, and recessive alleles do not fade away. Left to itself, a population is genetically \emph{stable}. Why?

Imagine the gene pool as a vast bag of gametes: a fraction $p$ carry allele $A$ and a fraction $q$ carry allele $a$. If mating is random, fertilisation is like drawing two gametes at random from the bag. The probability of drawing two $A$ gametes is $p \times p = p^2$; of drawing $A$ then $a$ (or $a$ then $A$) is $pq + pq = 2pq$; of drawing two $a$ gametes is $q \times q = q^2$. The figure below shows this as a Punnett square with gamete \emph{frequencies} instead of parental genotypes.

\begin{popfigure}
\begin{center}
\begin{tikzpicture}[scale=1.6]
\fill[popblueL] (0,1) rectangle (1,2);
\fill[popgreenL] (1,1) rectangle (2,2);
\fill[popgreenL] (0,0) rectangle (1,1);
\fill[popgold!30] (1,0) rectangle (2,1);
\draw[thick, popink] (0,0) grid (2,2);
\node[font=\bfseries, popblue] at (-0.55,1.5) {$A\;(p)$};
\node[font=\bfseries, poppink] at (-0.55,0.5) {$a\;(q)$};
\node[font=\bfseries, popblue] at (0.5,2.45) {$A\;(p)$};
\node[font=\bfseries, poppink] at (1.5,2.45) {$a\;(q)$};
\node at (0.5,1.5) {$AA\;\; p^2$};
\node at (1.5,1.5) {$Aa\;\; pq$};
\node at (0.5,0.5) {$Aa\;\; pq$};
\node at (1.5,0.5) {$aa\;\; q^2$};
\end{tikzpicture}
\end{center}
\legende{Derivation of the Hardy-Weinberg genotype frequencies. Random union of gametes drawn from a gene pool with allele frequencies $p$ (allele $A$) and $q$ (allele $a$) gives offspring genotypes $AA$, $Aa$ and $aa$ in the proportions $p^2 : 2pq : q^2$.}
\end{popfigure}

\begin{propriete}[The Hardy-Weinberg equilibrium]
In a large, randomly mating population, with no mutation, no migration, no natural selection and no genetic drift, the allele frequencies $p$ and $q$ and the genotype frequencies remain constant from generation to generation. The population is said to be in \cle{Hardy-Weinberg equilibrium}, and the genotype frequencies are given by:
\[ p + q = 1 \qquad \text{and} \qquad p^2 + 2pq + q^2 = 1 \]
where $p^2$ = frequency of $AA$, $2pq$ = frequency of $Aa$, $q^2$ = frequency of $aa$. (The derivation from the gamete Punnett square above is examinable.)
\end{propriete}

Notice that the second equation is just the binomial expansion $(p+q)^2 = p^2 + 2pq + q^2$, and since $p + q = 1$, the genotype frequencies automatically sum to $1$.

\begin{popfigure}
\begin{center}
\begin{tikzpicture}
\begin{axis}[
width=11cm, height=7cm,
xlabel={Allele frequency $p$ (of allele $A$)},
ylabel={Genotype frequency},
xmin=0, xmax=1, ymin=0, ymax=1,
xtick={0,0.2,0.4,0.6,0.8,1},
ytick={0,0.2,0.4,0.6,0.8,1},
legend style={at={(0.5,1.02)}, anchor=south, legend columns=3, font=\small},
grid=both, grid style={popline!40},
]
\addplot[domain=0:1, samples=100, thick, popblue] {x^2};
\addlegendentry{$p^2$ ($AA$)}
\addplot[domain=0:1, samples=100, thick, popgreen] {2*x*(1-x)};
\addlegendentry{$2pq$ ($Aa$)}
\addplot[domain=0:1, samples=100, thick, poppink] {(1-x)^2};
\addlegendentry{$q^2$ ($aa$)}
\end{axis}
\end{tikzpicture}
\end{center}
\legende{Genotype frequencies at Hardy-Weinberg equilibrium as functions of the allele frequency $p$. Note that heterozygotes ($2pq$) are most frequent when $p = q = 0.5$, and that a rare recessive allele (small $q$, i.e.\ $p$ near $1$) is carried mostly hidden in heterozygotes.}
\end{popfigure}

The graph carries an important message: when a recessive allele is rare ($q$ small), $q^2$ is tiny but $2pq$ is comparatively large. Most copies of a rare recessive allele are therefore \emph{hidden in healthy carriers} --- which is why recessive disorders cannot be eliminated simply because affected individuals are rare.

\courssub{Conditions for Equilibrium --- and What Happens When They Break Down}

The Hardy-Weinberg principle rests on five assumptions. Each one, when violated, becomes a mechanism of \emph{evolution}:

\begin{itemize}
\item \textbf{Large population size.} In small populations, chance events (\cle{genetic drift}) randomly change allele frequencies from one generation to the next.
\item \textbf{Random mating.} If individuals choose mates by genotype (e.g.\ inbreeding), homozygote frequencies rise and heterozygote frequencies fall, even though allele frequencies may stay the same.
\item \textbf{No mutation.} Mutation creates new alleles and slowly changes $p$ and $q$.
\item \textbf{No migration (gene flow).} Immigrants bring in alleles at different frequencies, altering the gene pool.
\item \textbf{No natural selection.} If one genotype survives or reproduces better, its alleles increase in frequency --- this is the driving force of adaptive evolution.
\end{itemize}

\begin{aretenir}[Why Hardy-Weinberg matters]
The Hardy-Weinberg equilibrium is the \emph{null hypothesis} of evolution. If a real population's genotype frequencies match $p^2 : 2pq : q^2$, it is not evolving at that locus. If they do not match, one of the five conditions is violated --- and the population is evolving. The equations also let us estimate things we cannot see directly, such as the number of carriers of a recessive disease allele.
\end{aretenir}

\courssub{Solving Hardy-Weinberg Problems}

\begin{methode}[The golden route: always start from $q^2$]
\begin{enumerate}
\item Identify the \textbf{recessive phenotype} (e.g.\ individuals affected by a recessive disorder). Its frequency in the population is $q^2$.
\item Take the square root: $q = \sqrt{q^2}$. This is the frequency of the recessive allele.
\item Find the dominant allele frequency: $p = 1 - q$.
\item Compute whatever is asked: carrier (heterozygote) frequency $= 2pq$; homozygous dominant frequency $= p^2$.
\item \textbf{Check:} $p^2 + 2pq + q^2$ must equal $1$.
\end{enumerate}
\end{methode}

\begin{attention}[The most common GCE mistake]
$q^2$ is the frequency of \textbf{affected individuals} (recessive phenotype), \textbf{not} the frequency of the recessive allele. If $1$ in $10\,000$ babies is affected, then $q^2 = 0.0001$ and $q = \sqrt{0.0001} = 0.01$ --- not $q = 0.0001$. Confusing the two makes the carrier frequency wrong by a factor of hundreds. Always take the square root!
\end{attention}

\begin{exoresolu}[Carriers of sickle cell anaemia]
In a Cameroonian population, the incidence of sickle cell anaemia (genotype $Hb^S Hb^S$) at birth is estimated at $4\%$. Assuming Hardy-Weinberg equilibrium, calculate (a) the frequency of the $Hb^S$ allele, (b) the frequency of the normal $Hb^A$ allele, (c) the percentage of the population who are healthy carriers (heterozygotes, $Hb^A Hb^S$).
\tcblower
\textbf{Step 1 --- Start from the recessive phenotype.} Affected individuals are $Hb^S Hb^S$, so
\[ q^2 = 0.04. \]
\textbf{Step 2 --- Find $q$.}
\[ q = \sqrt{0.04} = 0.2. \]
The frequency of the $Hb^S$ allele is $0.2$.
\textbf{Step 3 --- Find $p$.}
\[ p = 1 - q = 1 - 0.2 = 0.8. \]
The frequency of the $Hb^A$ allele is $0.8$.
\textbf{Step 4 --- Carrier frequency.}
\[ 2pq = 2 \times 0.8 \times 0.2 = 0.32. \]
So $32\%$ of the population are heterozygous carriers.
\textbf{Step 5 --- Check.} $p^2 + 2pq + q^2 = 0.64 + 0.32 + 0.04 = 1$. The answer is consistent.
\end{exoresolu}

\begin{exoresolu}[A slightly harder case]
In another population, $1$ child in $2\,500$ is born with cystic fibrosis, a recessive disorder. Estimate the proportion of carriers in the population.
\tcblower
$q^2 = \dfrac{1}{2500} = 0.0004$, so $q = \sqrt{0.0004} = 0.02$ and $p = 1 - 0.02 = 0.98$.
Carrier frequency: $2pq = 2 \times 0.98 \times 0.02 = 0.0392 \approx 3.9\%$, i.e.\ roughly $1$ person in $25$ is a carrier.
Check: $p^2 + 2pq + q^2 = 0.9604 + 0.0392 + 0.0004 = 1$. The answer is consistent.
Notice how carriers vastly outnumber affected individuals ($3.9\%$ versus $0.04\%$) --- the allele persists mainly in heterozygotes.
\end{exoresolu}

\begin{exoresolu}[Testing whether a population is in equilibrium]
In a population of $1\,000$ maize plants, the observed genotype counts are: $AA = 490$, $Aa = 420$, $aa = 90$. (a) Calculate the allele frequencies. (b) Calculate the genotype frequencies expected at Hardy-Weinberg equilibrium. (c) State whether the population appears to be in equilibrium.
\tcblower
\textbf{(a) Allele frequencies.} Total alleles $= 2 \times 1000 = 2000$.
\[ p = \frac{2(490) + 420}{2000} = \frac{1400}{2000} = 0.7, \qquad q = \frac{2(90) + 420}{2000} = \frac{600}{2000} = 0.3. \]
\textbf{(b) Expected genotype frequencies.}
\[ p^2 = 0.49, \qquad 2pq = 2 \times 0.7 \times 0.3 = 0.42, \qquad q^2 = 0.09. \]
Expected numbers out of $1\,000$: $AA = 490$, $Aa = 420$, $aa = 90$.
\textbf{(c) Conclusion.} The observed numbers match the Hardy-Weinberg expectations exactly, so the population appears to be in equilibrium at this locus: there is no evidence of selection, non-random mating, migration, mutation or drift acting on it.
\end{exoresolu}

\courssub{Applications of the Hardy-Weinberg Principle}

\begin{itemize}
\item \textbf{Medical genetics.} As shown above, the equations allow genetic counsellors to estimate carrier frequencies and the risk that two healthy carriers have an affected child --- essential for disorders such as sickle cell anaemia, which is common in Cameroon.
\item \textbf{Conservation genetics.} In small, isolated populations (for example the Cross River gorilla of the Cameroon--Nigeria border forests), genetic drift violates the large-population assumption, heterozygosity falls, and harmful recessive alleles are exposed. Conservationists use Hardy-Weinberg expectations to measure genetic diversity and design breeding programmes.
\item \textbf{Detecting evolution.} Comparing observed genotype frequencies with $p^2 : 2pq : q^2$ tells biologists whether selection, migration or non-random mating is acting on a locus.
\end{itemize}

\begin{savaistu}[Extension: heterozygote advantage and sickle cell in Cameroon]
Why is the $Hb^S$ allele so frequent in Central Africa when $Hb^S Hb^S$ individuals suffer severe anaemia? The answer is \emph{selection maintaining polymorphism}. In malaria-endemic regions of Cameroon, heterozygotes ($Hb^A Hb^S$) are significantly more resistant to severe malaria than $Hb^A Hb^A$ individuals, while not suffering sickle cell disease. The heterozygote is fitter than \emph{either} homozygote --- a situation called \cle{heterozygote advantage} (or balanced polymorphism). Selection therefore \emph{maintains} both alleles in the population, and the population is \textbf{not} in Hardy-Weinberg equilibrium at this locus: the ``no selection'' assumption is violated. This is a beautiful example of how breaking a Hardy-Weinberg condition drives and maintains evolutionary change.
\end{savaistu}

\begin{exercice}[Practice]
\begin{enumerate}
\item In a population of $1\,000$ plants, $480$ are $RR$ (red), $400$ are $Rr$ (pink) and $120$ are $rr$ (white). Calculate the allele frequencies $p$ and $q$.
\item In a human population in Hardy-Weinberg equilibrium, $9\%$ of individuals show a recessive phenotype. Calculate (a) the frequency of the recessive allele, (b) the frequency of carriers, (c) the frequency of homozygous dominant individuals.
\item Albinism (recessive) affects about $1$ in $20\,000$ newborns in a certain population. Estimate the percentage of carriers.
\item A population of $2\,000$ beetles contains $1\,280$ $BB$, $640$ $Bb$ and $80$ $bb$ individuals. Calculate $p$ and $q$, then state whether the population is in Hardy-Weinberg equilibrium, justifying your answer with the expected genotype numbers.
\end{enumerate}
\end{exercice}

\begin{corrige}[Exercise 1]
Total alleles $= 2000$. $R$ alleles $= 2(480) + 400 = 1360$, so $p = 0.68$. $r$ alleles $= 2(120) + 400 = 640$, so $q = 0.32$. Check: $0.68 + 0.32 = 1$.
\end{corrige}

\begin{corrige}[Exercise 2]
(a) $q^2 = 0.09 \Rightarrow q = 0.3$. (b) $p = 0.7$, so $2pq = 2 \times 0.7 \times 0.3 = 0.42$ ($42\%$ carriers). (c) $p^2 = 0.49$ ($49\%$). Check: $0.49 + 0.42 + 0.09 = 1$.
\end{corrige}

\begin{corrige}[Exercise 3]
$q^2 = \dfrac{1}{20000} = 0.00005$, so $q = \sqrt{0.00005} \approx 0.00707$ and $p \approx 0.99293$. Carriers: $2pq \approx 2 \times 0.99293 \times 0.00707 \approx 0.014$, i.e.\ about $1.4\%$ of the population.
\end{corrige}

\begin{corrige}[Exercise 4]
Total alleles $= 4000$. $p = \dfrac{2(1280) + 640}{4000} = \dfrac{3200}{4000} = 0.8$ and $q = \dfrac{2(80) + 640}{4000} = \dfrac{800}{4000} = 0.2$. Expected genotype frequencies: $p^2 = 0.64$, $2pq = 0.32$, $q^2 = 0.04$, giving expected numbers $BB = 1280$, $Bb = 640$, $bb = 80$. The observed numbers match the expectations exactly, so the population is in Hardy-Weinberg equilibrium at this locus.
\end{corrige}

\begin{attention}[Exam tip --- always verify]
Before leaving a Hardy-Weinberg question, check that $p + q = 1$ and that $p^2 + 2pq + q^2 = 1$. If your numbers fail either test, you have almost certainly forgotten a square root or mixed up $q$ with $q^2$. This ten-second check can save you several marks.
\end{attention}

\coursec{Variation, Polymorphism and Pedigree Analysis}

In the previous section we studied genes at the level of the \emph{population}: allele frequencies, genotype frequencies and the Hardy--Weinberg equilibrium. But why do populations contain different alleles at all? Because individuals of the same species are never identical. Look around your classroom in Yaound\'e or Bamenda: students differ in height, skin tone, blood group and ability to roll the tongue. These differences are the raw material of population genetics and of evolution. In this final section we study \cle{variation} itself, its types and causes, the special case of \cle{polymorphism}, and finally \cle{pedigree analysis}, the tool geneticists use to track traits through families.

\courssub{Variation and Its Types}

\begin{definition}[Variation]
\cle{Variation} is the occurrence of differences in characteristics (morphological, physiological or behavioural) between individuals of the same species. No two individuals are exactly alike, except identical twins, because they differ in genotype, in the environment they experienced, or in both.
\end{definition}

Variation is of two fundamental types, easily distinguished by asking one question: \emph{can the character take any value within a range, or only a few distinct classes?}

\begin{definition}[Continuous and discontinuous variation]
\begin{itemize}
\item \cle{Continuous variation} (quantitative variation): the character shows a complete range of intermediate values between two extremes, with no distinct classes. Examples: height, weight, skin colour, milk yield in cattle, grain yield in maize. When plotted, the population gives a \textbf{normal distribution} (bell-shaped curve).
\item \cle{Discontinuous variation} (qualitative variation): the character falls into a few clearly distinct classes with no intermediates. Examples: ABO blood groups (A, B, AB or O), ability/inability to roll the tongue, albinism, sex.
\end{itemize}
\end{definition}

\begin{propriete}[Genetic basis of the two types of variation]
\begin{itemize}
\item Continuous variation is \textbf{polygenic}: it is controlled by many genes (polygenes), each with a small additive effect, and it is strongly influenced by the \textbf{environment} (nutrition, climate, disease). The many possible genotypes, blurred by environmental effects, produce the smooth bell curve.
\item Discontinuous variation is controlled by \textbf{one gene or very few genes}, usually with two or a few alleles, and is little or not affected by the environment. Your blood group is fixed at fertilisation: no diet can change group A into group O.
\end{itemize}
This property is directly examinable: you must be able to classify a given character and justify your answer.
\end{propriete}

\begin{popfigure}
\begin{center}
\begin{tikzpicture}
\begin{axis}[
  width=0.52\linewidth, height=5.2cm,
  title={Continuous: height},
  title style={font=\small},
  xlabel={Height (cm)}, ylabel={Number of students},
  xlabel style={font=\small}, ylabel style={font=\small},
  tick label style={font=\scriptsize},
  ymin=0, ymax=0.06, xmin=140, xmax=200,
  xtick={150,160,170,180,190},
  axis lines=left]
\addplot[domain=145:195, samples=80, thick, popblue] {0.05*exp(-((x-170)^2)/(2*6^2))};
\addplot[domain=145:195, samples=80, smooth, fill=popblueL, draw=none] {0.05*exp(-((x-170)^2)/(2*6^2))} \closedcycle;
\end{axis}
\end{tikzpicture}\hfill
\begin{tikzpicture}
\begin{axis}[
  width=0.45\linewidth, height=5.2cm,
  ybar, bar width=16pt,
  title={Discontinuous: ABO groups},
  title style={font=\small},
  ylabel={Number of students},
  ylabel style={font=\small},
  tick label style={font=\scriptsize},
  symbolic x coords={O, A, B, AB},
  xtick=data, ymin=0, ymax=60,
  axis lines=left, nodes near coords, nodes near coords style={font=\scriptsize}]
\addplot[fill=poporange, draw=popdark] coordinates {(O,48) (A,32) (B,15) (AB,5)};
\end{axis}
\end{tikzpicture}
\end{center}
\legende{Left: continuous variation in height among 100 Upper Sixth students gives a bell-shaped (normal) distribution. Right: discontinuous variation in ABO blood group gives separate classes with no intermediates.}
\end{popfigure}

\begin{attention}[Do not confuse the two types]
Height is continuous even though we sometimes group people as ``short'' or ``tall'' --- those classes are artificial. Blood group is discontinuous even if we measure it on thousands of people. The test is: \emph{are there real intermediates in nature?} Also note that tongue rolling is discontinuous, \textbf{not} because few people can do it, but because there are only two classes.
\end{attention}

\courssub{Causes of Variation}

Variation arises from two great sources, which usually act together.

\textbf{1. Genetic (inherited) causes.}
\begin{itemize}
\item \cle{Mutation}: the ultimate source of \emph{new} alleles (see below).
\item \cle{Independent assortment} of chromosomes at meiosis I: with $n = 23$ pairs, humans can produce about $2^{23}$ different gamete types from this alone.
\item \cle{Crossing over} at prophase I: exchanges between non-sister chromatids create new combinations of alleles on the same chromosome.
\item \cle{Random fertilisation}: any sperm may fuse with any ovum, multiplying combinations further.
\end{itemize}

\textbf{2. Environmental causes.} Nutrition, temperature, light, disease and lifestyle modify the \emph{phenotype} without changing the genotype. A well-fed child in Douala may grow taller than his genetic potential would allow in conditions of malnutrition; skin darkens under the tropical sun.

\begin{definition}[Genotype--environment interaction]
The \cle{phenotype} is the product of the \textbf{genotype} and the \textbf{environment}: phenotype $=$ genotype $+$ environmental effect. Only the genetic component of variation is heritable; acquired characteristics (a scar, a trained muscle) are \emph{not} transmitted to offspring.
\end{definition}

\textbf{Mutations revisited.} A \cle{gene mutation} (point mutation: substitution, insertion or deletion of bases) creates a new allele of a gene --- this is how alleles such as $H\!b^S$ (sickle cell) arose. A \cle{chromosome mutation} (deletion, duplication, inversion, translocation of chromosome segments, or change in chromosome number as in Down syndrome, trisomy 21) alters many genes at once and is usually more severe. Mutations occurring in \textbf{germ cells} (gamete-forming cells) are heritable; those in somatic cells are not.

\courssub{Polymorphism}

\begin{definition}[Polymorphism]
\cle{Polymorphism} is the occurrence, within a population of the same species living together, of two or more sharply distinct forms (morphs) at frequencies too high to be maintained by recurrent mutation alone. When the forms are based on different alleles of a gene, we speak of \textbf{genetic polymorphism} (e.g. ABO blood groups, sickle cell trait).
\end{definition}

We distinguish:
\begin{itemize}
\item \textbf{Transient polymorphism}: one form is progressively replacing another, so the polymorphism is temporary. Classic example: the peppered moth (\emph{Biston betularia}), where the dark (melanic) form spread during industrial pollution and declined afterwards.
\item \textbf{Balanced (stable) polymorphism}: the different forms are maintained at stable frequencies by natural selection, typically because the \textbf{heterozygote is advantaged} (heterozygote advantage).
\end{itemize}

\begin{savaistu}[Sickle cell and malaria: balanced polymorphism in Cameroon]
The gene for haemoglobin has two common alleles: $H\!b^A$ (normal) and $H\!b^S$ (sickle). Genotype $H\!b^S H\!b^S$ causes sickle cell anaemia, a serious disease; $H\!b^A H\!b^A$ individuals are healthy but fully susceptible to severe malaria. Heterozygotes $H\!b^A H\!b^S$ (sickle cell \emph{trait}) are healthy \emph{and} more resistant to malaria. In malaria-endemic regions such as most of Cameroon, selection favours the heterozygote, so the $H\!b^S$ allele is maintained in the population at a stable frequency despite being lethal in double dose. This is the best-known example of \cle{balanced polymorphism} in humans --- and a frequent GCE question.
\end{savaistu}

\courssub{Pedigree Analysis}

When a trait runs in a human family, we cannot perform controlled crosses as Mendel did with peas. Instead we use \cle{pedigree analysis}: a diagram of the family's history showing who is affected, from which we deduce the mode of inheritance.

\begin{definition}[Pedigree]
A \cle{pedigree} is a standardised chart showing the occurrence of a trait through several generations of a family, using conventional symbols. It allows geneticists to deduce genotypes, identify carriers and calculate the risk of the trait appearing in future children.
\end{definition}

\begin{popfigure}
\begin{center}
\begin{tikzpicture}[scale=0.9, every node/.style={font=\small}]
% symbols legend
\draw (0,0) rectangle (0.5,0.5); \node[right] at (0.7,0.25) {Unaffected male};
\draw[fill=popink] (3.6,0) rectangle (4.1,0.5); \node[right] at (4.3,0.25) {Affected male};
\draw (0,-1) circle (0.25); \node[right] at (0.7,-1) {Unaffected female};
\draw[fill=popink] (3.85,-1) circle (0.25); \node[right] at (4.3,-1) {Affected female};
\draw (0,-2) circle (0.25); \fill (0,-2) circle (0.05); \node[right] at (0.7,-2) {Carrier female (X-linked)};
\draw (3.7,-2) -- (4.3,-2); \draw (3.85,-2) circle (0.25); \draw (3.6,-2) rectangle (4.1,-1.5); \node[right] at (4.5,-1.9) {Mating line};
\draw (0,-3) circle (0.25); \draw (-0.2,-3.25) -- (0.2,-2.75); \node[right] at (0.7,-3) {Deceased individual};
\draw (3.85,-3) circle (0.25); \node at (3.85,-3) {?}; \node[right] at (4.3,-3) {Sex unspecified / unknown};
\draw (0,-4) -- (0,-3.6); \draw (-0.5,-3.6) -- (0.5,-3.6); \draw (-0.5,-3.6) -- (-0.5,-4.2); \draw (0.5,-3.6) -- (0.5,-4.2);
\draw (-0.5,-4.45) rectangle (-0.25,-4.2); \draw (0.25,-4.2) circle (0.25);
\node[right] at (0.9,-4.1) {Sibship line (children)};
\end{tikzpicture}
\end{center}
\legende{Standard pedigree symbols. Generations are numbered I, II, III\ldots\ from top to bottom; individuals are numbered within each generation from left to right.}
\end{popfigure}

\begin{popfigure}
\begin{center}
\begin{tikzpicture}[scale=1.0, every node/.style={font=\small}]
% Generation I
\draw (1,4) rectangle (1.5,4.5); \node[above] at (1.25,4.5) {I-1};
\draw (2.5,4.25) circle (0.25); \node[above] at (2.5,4.5) {I-2};
\draw (1.5,4.25) -- (2.25,4.25); % mating line
% vertical down from mating
\draw (1.875,4.25) -- (1.875,3.4);
\draw (0.5,3.4) -- (3.25,3.4); % sibship line
% children gen II
\draw (0.5,3.4) -- (0.5,2.9); \draw (0.25,2.4) rectangle (0.75,2.9); \node[below] at (0.5,2.4) {II-1};
\draw (1.875,3.4) -- (1.875,2.9); \draw (1.875,2.65) circle (0.25); \node[below] at (1.875,2.4) {II-2};
\draw (3.25,3.4) -- (3.25,2.9); \draw (3,2.4) rectangle (3.5,2.9); \node[below] at (3.25,2.4) {II-3};
% II-2 marries a spouse (unrelated male)
\draw (2.125,2.65) -- (2.75,2.65); % mating line to spouse
\draw (3,2.4) rectangle (3.5,2.9); \node[below] at (3.25,2.4) {Spouse};
% II-2 x Spouse children (Gen III)
\draw (2.94,2.65) -- (2.94,1.8);
\draw (2.2,1.8) -- (3.7,1.8); % sibship line
\draw (2.2,1.8) -- (2.2,1.35); \draw[fill=popink] (1.95,0.85) rectangle (2.45,1.35); \node[below] at (2.2,0.85) {III-1};
\draw (2.94,1.8) -- (2.94,1.35); \draw (2.94,1.1) circle (0.25); \node[below] at (2.94,0.85) {III-2};
\draw (3.7,1.8) -- (3.7,1.35); \draw (3.7,1.1) circle (0.25); \node[below] at (3.7,0.85) {III-3};
\end{tikzpicture}
\end{center}
\legende{Pedigree of an autosomal recessive trait (e.g. albinism). Unaffected parents I-1 and I-2 have a carrier daughter II-2; she marries an unrelated carrier male (Spouse). Their son III-1 is affected. II-2 and her spouse must both be carriers.}
\end{popfigure}

\begin{methode}[Four-step strategy to determine the mode of inheritance from a pedigree]
\begin{enumerate}
\item \textbf{Are both sexes affected in roughly equal numbers?} If yes, suspect an \textbf{autosomal} trait; if mostly males are affected, suspect \textbf{X-linked recessive}.
\item \textbf{Do unaffected parents produce an affected child?} If yes, the trait must be \textbf{recessive} and both parents are heterozygous carriers. If every affected child has at least one affected parent, suspect \textbf{dominance}.
\item \textbf{Does the trait appear in every generation?} Dominant traits typically appear in every generation; recessive traits often \emph{skip} generations.
\item \textbf{Check the transmission through males.} In X-linked recessive inheritance, an affected father transmits the allele to \emph{all} his daughters (carriers) but to \emph{none} of his sons; affected grandfathers reappear as affected grandsons through carrier daughters.
\end{enumerate}
\end{methode}

\begin{attention}[The golden rule of recessive traits]
``Unaffected parents producing an affected child'' means the trait is \textbf{recessive} and \textbf{both parents are carriers} ($Aa \times Aa$). It can \emph{never} be dominant: a dominant allele cannot hide in an unaffected person. Conversely, two affected parents with an unaffected child implies a \textbf{dominant} trait with both parents heterozygous.
\end{attention}

\begin{exoresolu}[Identifying genotypes and calculating a risk]
In the pedigree above (albinism, autosomal recessive; let $A$ = normal pigment, $a$ = albinism): (a) Give the genotypes of I-1, I-2, II-2, Spouse and III-1. (b) If III-2 marries a heterozygous man, what is the probability that their first child is albino, given that III-2 has a $2/3$ chance of being a carrier?
\tcblower
(a) III-1 is affected, so his genotype is $aa$. His parents II-2 and Spouse are unaffected but produced an affected child, so both must be carriers: $Aa \times Aa$. 
II-2 is a carrier ($Aa$). She inherited the $a$ allele from one of her parents (I-1 or I-2). Since both I-1 and I-2 are unaffected, at least one of them is a carrier ($Aa$). The other parent could be $AA$ or $Aa$; without evidence of consanguinity or multiple affected offspring in generation II, we cannot assume both are carriers. Standard notation: the parent who transmitted $a$ is $Aa$, the other is $A\_$ (phenotypically normal, genotype uncertain). If the trait is rare, the other is usually $AA$.

(b) III-2 is unaffected from the cross $Aa \times Aa$, so $P(\text{carrier}) = \dfrac{2}{3}$. Her husband is $Aa$. A child is albino ($aa$) only if III-2 is a carrier \emph{and} transmits $a$ \emph{and} the husband transmits $a$:
\[
P(aa) = \frac{2}{3} \times \frac{1}{2} \times \frac{1}{2} = \frac{1}{6}.
\]
The probability that their first child is albino is $\boxed{\dfrac{1}{6}}$.
\end{exoresolu}

\begin{exemplebox}[Recognising X-linked recessive inheritance]
In a pedigree of red--green colour blindness: the grandfather (generation I) is affected; his daughter (generation II) is unaffected but is an \emph{obligatory carrier} (she received his only X chromosome); her son (generation III) is affected. No male-to-male transmission occurs, and far more males than females are affected. This pattern --- \textbf{affected grandfather $\rightarrow$ carrier daughter $\rightarrow$ affected grandson} --- is the signature of X-linked recessive inheritance. Always look for it first in GCE pedigree questions.
\end{exemplebox}

\begin{exercice}[Pedigree practice]
A couple, both with normal skin pigmentation, have four children; one daughter is albino. (1) State the mode of inheritance of albinism and justify. (2) Give the genotypes of the parents and of the albino daughter. (3) What is the probability that their next child is an albino son? (4) The albino daughter later marries an albino man; what proportion of their children will be albino?
\end{exercice}

\begin{corrige}[Exercise 1]
(1) Autosomal \textbf{recessive}: unaffected parents produced an affected child, so the allele was hidden in both parents; both sexes can be affected. (2) Parents: $Aa \times Aa$; albino daughter: $aa$. (3) $P(\text{albino}) = \tfrac{1}{4}$ and $P(\text{son}) = \tfrac{1}{2}$, independent events: $P = \tfrac{1}{4} \times \tfrac{1}{2} = \tfrac{1}{8}$. (4) Cross $aa \times aa$: all children ($100\%$) are albino.
\end{corrige}

\begin{aretenir}[Key points of the section]
\begin{itemize}
\item Variation = differences between individuals of the same species; it is \textbf{continuous} (polygenic, environment-sensitive, bell curve) or \textbf{discontinuous} (one/few genes, distinct classes).
\item Variation comes from mutation, independent assortment, crossing over, random fertilisation, and the environment; only genetic variation is heritable.
\item Polymorphism = stable coexistence of distinct forms in a population; sickle cell/malaria is the classic case of \textbf{balanced polymorphism} by heterozygote advantage.
\item Pedigrees use standard symbols; unaffected parents with an affected child $\Rightarrow$ recessive trait, both parents carriers; affected grandfather $\rightarrow$ carrier daughter $\rightarrow$ affected grandson $\Rightarrow$ X-linked recessive.
\item Risk calculations combine Mendelian ratios with the rules of probability.
\end{itemize}
\end{aretenir}

\newpage

\coursec{Integration activity}

\coursec{Integration Activity: Genetic Counselling in Nkambe}

\courssub{Aims of the activity}
\begin{itemize}
    \item Consolidate understanding of Mendelian inheritance patterns: complete dominance, codominance, incomplete dominance, sex linkage, multiple alleles, epistasis and lethal genes.
    \item Apply the Hardy–Weinberg principle to estimate allele and genotype frequencies in a real population.
    \item Construct and interpret a three‑generation pedigree, deducing genotypes and calculating probabilities for future offspring.
    \item Use ABO blood group data to evaluate a disputed paternity case.
    \item Communicate genetic risk information clearly and compassionately in a written advisory note.
    \item Develop teamwork, critical reasoning and oral presentation skills through group work and plenary defence.
\end{itemize}

\courssub{Situation}
You are a team of genetic counsellors at the Nkambe District Hospital in the Donga‑Mantung Division, North West Region of Cameroon. A family has been referred to you with three interconnected concerns. The case file contains the following information.

\textbf{1. Family pedigree and haemophilia history}\\
Haemophilia A is an X‑linked recessive disorder (\cle{X\textsuperscript{h}} = haemophilia allele, \cle{X\textsuperscript{H}} = normal allele). The family spans three generations. The known phenotypes are:
\begin{itemize}
    \item Generation I: Grandfather (I‑1) – normal; Grandmother (I‑2) – phenotypically normal but known to be a carrier (one affected brother in her family).
    \item Generation II: They have two children: a son (II‑1) who has severe haemophilia, and a daughter (II‑2) who is phenotypically normal. II‑1 marries a normal woman (II‑3). II‑2 marries a normal man (II‑4).
    \item Generation III: From II‑1 × II‑3: a son (III‑1) normal, a daughter (III‑2) phenotypically normal. From II‑2 × II‑4: four children – a son (III‑3) with haemophilia, a daughter (III‑4) phenotypically normal, a son (III‑5) normal, a daughter (III‑6) normal.
    \item III‑2 (now pregnant) is married to a phenotypically normal man. The couple requests the probability that their unborn child will have haemophilia.
\end{itemize}

\textbf{2. Disputed paternity – ABO blood groups}\\
The mother (III‑2) has blood group A. Her child (already born, III‑2’s first child) has blood group O. The alleged father (III‑2’s husband) has blood group B. The ABO system is controlled by three alleles: \cle{I\textsuperscript{A}}, \cle{I\textsuperscript{B}} (codominant) and \cle{i} (recessive). Determine whether the alleged father can be excluded as the biological father.

\textbf{3. Population data – Sickle cell disease}\\
Health statistics for the Donga‑Mantung Division show that 4\% of newborns are homozygous for sickle cell disease (genotype \cle{Hb\textsuperscript{S}Hb\textsuperscript{S}}). The sickle cell allele (\cle{Hb\textsuperscript{S}}) is recessive to the normal allele (\cle{Hb\textsuperscript{A}}). Assuming the population is in Hardy–Weinberg equilibrium for this locus, estimate the percentage of carriers (heterozygotes \cle{Hb\textsuperscript{A}Hb\textsuperscript{S}}) in the population.

\textbf{4. Advisory note}\\
The family needs a clear, non‑technical summary of the risks and recommendations.

\begin{popfigure}
\begin{tikzpicture}[scale=0.85,
    male/.style={draw, rectangle, minimum size=0.9cm},
    female/.style={draw, circle, minimum size=0.9cm},
    affected/.style={fill=black, text=white},
    carrier/.style={fill=gray!50},
    normal/.style={fill=white}
]
% Generation I
\node[male, normal] (I1) at (0,5) {I-1};
\node[female, carrier] (I2) at (3,5) {I-2};
\draw (I1) -- (I2);
\draw (1.5,5) -- (1.5,4);
% Generation II
\node[male, affected] (II1) at (0,3) {II-1};
\node[female, carrier] (II2) at (3,3) {II-2};
\node[female, normal] (II3) at (-2,3) {II-3};
\node[male, normal] (II4) at (5,3) {II-4};
\draw (II1) -- (II3);
\draw (II2) -- (II4);
% Vertical lines to Gen III
\draw (-1,3) -- (-1,2);
\draw (1.5,3) -- (1.5,2);
\draw (4,3) -- (4,2);
% Generation III - children of II-1 x II-3
\node[male, normal] (III1) at (-2,1) {III-1};
\node[female, carrier] (III2) at (0,1) {III-2};
% Generation III - children of II-2 x II-4
\node[male, affected] (III3) at (3,1) {III-3};
\node[female, carrier] (III4) at (4.5,1) {III-4};
\node[male, normal] (III5) at (6,1) {III-5};
\node[female, normal] (III6) at (7.5,1) {III-6};
% Connections
\draw (II1) |- (-1,2) -| (III1);
\draw (II1) |- (-1,2) -| (III2);
\draw (II2) |- (1.5,2) -| (III3);
\draw (II2) |- (1.5,2) -| (III4);
\draw (II4) |- (4,2) -| (III5);
\draw (II4) |- (4,2) -| (III6);
% Legend
\node[anchor=west] at (8.5,5) {\textbf{Legend:}};
\node[male, affected, anchor=west] at (8.5,4.5) {};
\node[anchor=west] at (9.5,4.5) {Affected male};
\node[female, carrier, anchor=west] at (8.5,4) {};
\node[anchor=west] at (9.5,4) {Carrier female};
\node[male, normal, anchor=west] at (8.5,3.5) {};
\node[anchor=west] at (9.5,3.5) {Normal male};
\node[female, normal, anchor=west] at (8.5,3) {};
\node[anchor=west] at (9.5,3) {Normal female};
\end{tikzpicture}
\legende{Three‑generation pedigree of the Nkambe family. Black squares = males with haemophilia; grey circles = carrier females; white symbols = unaffected individuals.}
\end{popfigure}

\courssub{Tasks}
\begin{enumerate}
    \item \textbf{Pedigree analysis and risk calculation}
    \begin{enumerate}
        \item Using the pedigree, deduce the most likely genotypes of all individuals for the haemophilia locus. Write each genotype in the form \cle{X\textsuperscript{H}X\textsuperscript{h}} etc.
        \item Calculate the probability that the unborn child of III‑2 (carrier female) and her normal husband will be affected by haemophilia. Show all steps.
    \end{enumerate}
    \item \textbf{ABO paternity exclusion}
    \begin{enumerate}
        \item List all possible genotypes for the mother (group A), the child (group O) and the alleged father (group B).
        \item Using a Punnett square or logical deduction, decide whether the alleged father can be excluded as the biological father. Justify your answer.
    \end{enumerate}
    \item \textbf{Hardy–Weinberg estimation for sickle cell}
    \begin{enumerate}
        \item Write the Hardy–Weinberg equations for allele frequencies (\cle{p}, \cle{q}) and genotype frequencies.
        \item Given that \cle{q\textsuperscript{2} = 0.04}, calculate \cle{p}, \cle{q} and the carrier frequency \cle{2pq}. Express the carrier frequency as a percentage.
    \end{enumerate}
    \item \textbf{Advisory note}
    Write a concise advisory note (150–200 words) addressed to the family. It must include:
    \begin{itemize}
        \item The probability that the unborn child will have haemophilia.
        \item The conclusion of the paternity test.
        \item The estimated carrier frequency of sickle cell trait in the community.
        \item Practical recommendations (e.g., prenatal diagnosis, genetic counselling for relatives, newborn screening).
    \end{itemize}
    The note should be in plain language, avoiding jargon.
\end{enumerate}

\courssub{Model answer}

\begin{methode}[Step‑by‑step pedigree analysis for X‑linked recessive traits]
\begin{enumerate}
\item Identify the mode of inheritance (here: X‑linked recessive).
\item Assign alleles: \cle{X\textsuperscript{H}} (normal), \cle{X\textsuperscript{h}} (haemophilia).
\item Start with known genotypes: affected males are \cle{X\textsuperscript{h}Y}; known carrier females are \cle{X\textsuperscript{H}X\textsuperscript{h}}.
\item Use the fact that a male inherits his X chromosome from his mother and his Y from his father.
\item A female inherits one X from each parent.
\item If a phenotypically normal female has an affected son, she must be a carrier (\cle{X\textsuperscript{H}X\textsuperscript{h}}).
\item Deduce genotypes generation by generation, using the phenotypes of offspring to resolve uncertainties.
\end{enumerate}
\end{methode}

\textbf{Task 1 – Pedigree analysis and risk calculation}

\begin{enumerate}
    \item \textbf{Genotype deduction}\\
    Haemophilia A is X‑linked recessive. Let \cle{X\textsuperscript{H}} = normal allele, \cle{X\textsuperscript{h}} = haemophilia allele.
    \begin{itemize}
        \item \textbf{I‑1 (normal male):} \cle{X\textsuperscript{H}Y} — only one X, must be normal.
        \item \textbf{I‑2 (carrier female):} \cle{X\textsuperscript{H}X\textsuperscript{h}} — stated as known carrier (affected brother).
        \item \textbf{II‑1 (affected male):} \cle{X\textsuperscript{h}Y} — inherited \cle{X\textsuperscript{h}} from mother I‑2.
        \item \textbf{II‑2 (normal female):} Initially 50\% chance carrier. Because she has an affected son (III‑3), she \emph{must} have passed an \cle{X\textsuperscript{h}} allele; therefore her genotype is \cle{X\textsuperscript{H}X\textsuperscript{h}}.
        \item \textbf{II‑3 (normal female, wife of II‑1):} \cle{X\textsuperscript{H}X\textsuperscript{H}} — no family history, and all her children are phenotypically normal (though daughters are obligate carriers).
        \item \textbf{II‑4 (normal male, husband of II‑2):} \cle{X\textsuperscript{H}Y}.
        \item \textbf{III‑1 (son of II‑1 × II‑3):} Father gives Y, mother gives \cle{X\textsuperscript{H}} → \cle{X\textsuperscript{H}Y} (normal).
        \item \textbf{III‑2 (daughter of II‑1 × II‑3):} Father gives \cle{X\textsuperscript{h}}, mother gives \cle{X\textsuperscript{H}} → \cle{X\textsuperscript{H}X\textsuperscript{h}} (obligate carrier).
        \item \textbf{III‑3 (son of II‑2 × II‑4):} Mother gives \cle{X\textsuperscript{h}} (50\% chance), father gives Y → \cle{X\textsuperscript{h}Y} (affected).
        \item \textbf{III‑4 (daughter of II‑2 × II‑4):} Mother gives \cle{X\textsuperscript{h}} (50\%), father gives \cle{X\textsuperscript{H}} → \cle{X\textsuperscript{H}X\textsuperscript{h}} (carrier).
        \item \textbf{III‑5 (son):} Mother gives \cle{X\textsuperscript{H}}, father gives Y → \cle{X\textsuperscript{H}Y} (normal).
        \item \textbf{III‑6 (daughter):} Mother gives \cle{X\textsuperscript{H}}, father gives \cle{X\textsuperscript{H}} → \cle{X\textsuperscript{H}X\textsuperscript{H}} (normal).
    \end{itemize}
    \item \textbf{Probability for unborn child of III‑2}\\
    III‑2 genotype: \cle{X\textsuperscript{H}X\textsuperscript{h}} (carrier). Husband genotype: \cle{X\textsuperscript{H}Y} (normal).
    Possible gametes:
    \begin{itemize}
        \item Mother: \cle{X\textsuperscript{H}} or \cle{X\textsuperscript{h}} (each 1/2).
        \item Father: \cle{X\textsuperscript{H}} or \cle{Y} (each 1/2).
    \end{itemize}
    Punnett square:
    \[
    \begin{array}{c|cc}
    & X^H & Y \\ \hline
    X^H & X^HX^H\ (\text{normal female}) & X^HY\ (\text{normal male}) \\
    X^h & X^HX^h\ (\text{carrier female}) & X^hY\ (\text{affected male})
    \end{array}
    \]
    Each cell probability = 1/4. Only the \cle{X\textsuperscript{h}Y} combination yields an affected child (a son with haemophilia). Therefore:
    \[
    P(\text{affected child}) = \frac{1}{4} = 0.25 = 25\%.
    \]
    If the sex of the child is known (e.g., by ultrasound), the conditional probabilities are:
    \begin{itemize}
        \item If male: 1/2 (50\%) affected.
        \item If female: 0\% affected (but 50\% carrier).
    \end{itemize}
\end{enumerate}

\begin{attention}[Common mistake in X‑linked risk calculation]
Do not simply say “50\% chance” because the mother is a carrier. The 50\% refers to the chance of passing the allele, but the child must also be male to be affected. The overall probability is 1/2 (pass allele) × 1/2 (male) = 1/4.
\end{attention}

\textbf{Task 2 – ABO paternity exclusion}

\begin{enumerate}
    \item \textbf{Possible genotypes}\\
    \begin{itemize}
        \item Mother (group A): \cle{I$^{A}$ I$^{A}$} or \cle{I$^{A}$ i}.
        \item Child (group O): \cle{ii} (only genotype for group O).
        \item Alleged father (group B): \cle{I$^{B}$ I$^{B}$} or \cle{I$^{B}$ i}.
    \end{itemize}
    \item \textbf{Exclusion test}\\
    The child \cle{ii} must receive an \cle{i} allele from each parent.
    \begin{itemize}
        \item Mother can provide \cle{i} only if she is \cle{I$^{A}$ i}. Since she has a group O child, she \emph{must} be \cle{I$^{A}$ i}.
        \item Alleged father can provide \cle{i} only if he is \cle{I$^{B}$ i}. If he were \cle{I$^{B}$ I$^{B}$}, he could not give an \cle{i} allele and would be excluded.
    \end{itemize}
    The case file states the alleged father is group B but does not specify his genotype. In the absence of further information, a group B individual \emph{can} be \cle{I$^{B}$ i} (heterozygous). Therefore, \textbf{the alleged father cannot be excluded} on the basis of ABO blood groups alone. (If subsequent testing showed he was \cle{I$^{B}$ I$^{B}$}, he would be excluded.)
\end{enumerate}

\begin{definition}[Hardy–Weinberg equilibrium]
For a diallelic locus in a large, randomly mating population with no selection, migration, mutation or genetic drift, allele frequencies remain constant from generation to generation, and genotype frequencies are given by the binomial expansion \cle{(p+q)$^{2}$ = p$^{2}$ + 2pq + q$^{2}$ = 1}, where \cle{p} is the frequency of the dominant allele and \cle{q} the frequency of the recessive allele.
\end{definition}

\textbf{Task 3 – Hardy–Weinberg estimation for sickle cell}

\begin{enumerate}
    \item \textbf{Hardy–Weinberg equations}\\
    For a diallelic locus with alleles \cle{A} (normal, frequency \cle{p}) and \cle{S} (sickle, frequency \cle{q}):
    \[
    p + q = 1 \qquad \text{(allele frequencies)}
    \]
    \[
    p^2 + 2pq + q^2 = 1 \qquad \text{(genotype frequencies: } AA, AS, SS\text{)}
    \]
    \item \textbf{Calculations}\\
    Given: \cle{q$^{2}$ = 0.04} (frequency of \cle{SS} newborns).
    \[
    q = \sqrt{0.04} = 0.20
    \]
    \[
    p = 1 - q = 1 - 0.20 = 0.80
    \]
    Carrier frequency (heterozygotes \cle{AS}) = \cle{2pq}:
    \[
    2pq = 2 \times 0.80 \times 0.20 = 0.32 = 32\%.
    \]
    Thus, approximately \textbf{32\% of the population are carriers of the sickle cell trait}.
\end{enumerate}

\textbf{Task 4 – Advisory note (sample)}

\begin{quote}
\textbf{Advisory Note for the Nkambe Family}\\
\textbf{Date:} [Insert date] \quad \textbf{From:} Genetic Counselling Team, Nkambe District Hospital\\
\textbf{To:} The family of [Family name]\\

Dear Family Members,

After careful study of your family history and the laboratory data, we wish to share the following information in simple terms.

\begin{enumerate}
    \item \textbf{Haemophilia risk:} The mother (III‑2) carries the haemophilia gene. For each pregnancy, there is a 1 in 4 chance (25\%) that the child will be a boy with haemophilia. If the baby is a girl, she will not have the disease but has a 50\% chance of being a carrier like her mother. We recommend prenatal diagnosis (e.g., chorionic villus sampling or amniocentesis) to know the baby’s status early, and delivery in a centre equipped to manage haemophilia.
    \item \textbf{Paternity question:} The ABO blood groups of the mother (group A), the child (group O) and the alleged father (group B) are compatible. The alleged father \textbf{cannot be excluded} as the biological father based on this test alone. If absolute certainty is required, DNA‑based paternity testing is available.
    \item \textbf{Sickle cell in the community:} Health records show that 4\% of newborns in our division have sickle cell disease. This means about 32\% of people in the community carry the sickle cell trait (AS). Carriers are usually healthy but can pass the gene to their children. We advise that all family members, especially those planning marriage, know their genotype (AA, AS or SS) through a simple blood test (haemoglobin electrophoresis).
    \item \textbf{Recommendations:}
    \begin{itemize}
        \item Schedule prenatal diagnosis for the current pregnancy.
        \item Offer haemoglobin electrophoresis to all relatives of reproductive age.
        \item Ensure that any male child born with haemophilia is referred promptly to the haematology clinic for factor replacement therapy.
        \item Consider genetic counselling for future marriages within the extended family.
    \end{itemize}
\end{enumerate}

We are here to support you. Please do not hesitate to ask questions or request further tests.

Yours sincerely,\\
The Genetic Counselling Team\\
Nkambe District Hospital
\end{quote}

\begin{aretenir}[Key integration points]
\begin{itemize}
\item Pedigree analysis combines Mendelian rules with sex‑linked inheritance; always track the X chromosome.
\item ABO paternity exclusion relies on the fact that a group O child (\cle{ii}) must receive an \cle{i} from each parent.
\item Hardy–Weinberg equilibrium assumes random mating, no selection, migration, mutation or drift; it provides a good first estimate for carrier frequencies.
\item Genetic counselling must translate probabilities into clear, actionable advice while respecting cultural sensitivities.
\end{itemize}
\end{aretenir}

\newpage

\coursec{Methods and worked examples}

\courssub{Skill 1: Solving a Monohybrid Cross with Complete Dominance}

\begin{methode}[Predicting offspring from a monohybrid cross]
\begin{enumerate}
\item \textbf{Read the statement carefully} and identify the \cle{character}, the two contrasting \cle{traits}, and which trait is \cle{dominant} (often revealed by the F$_1$ result or stated explicitly).
\item \textbf{Choose symbols}: one letter only; capital for the dominant allele, small letter for the recessive allele (e.g. $T$ = tall, $t$ = dwarf). State your symbols clearly.
\item \textbf{Write the genotypes} of the parents from their phenotypes: a dominant phenotype may be $TT$ or $Tt$ (use the cross details to decide); a recessive phenotype must be $tt$.
\item \textbf{Determine the gametes}: each parent passes only \emph{one} allele of each gene to each gamete (Mendel's law of segregation).
\item \textbf{Draw a Punnett square} (checkerboard) and fill in the zygote genotypes.
\item \textbf{State the genotypic ratio} (e.g. $1\,TT : 2\,Tt : 1\,tt$) and the \textbf{phenotypic ratio} (e.g. $3$ tall : $1$ dwarf).
\item \textbf{Check}: the fractions must add up to 1 (or the numbers to the total offspring).
\end{enumerate}
\textbf{Reflexes:} a $3:1$ F$_2$ ratio signals complete dominance with heterozygous parents; a $1:1$ ratio signals a test cross ($Aa \times aa$).
\end{methode}

\begin{exoresolu}[Monohybrid cross in maize]
In maize, the allele for purple grains ($P$) is dominant over the allele for yellow grains ($p$). A pure-breeding purple maize plant is crossed with a yellow-grained plant. \textbf{(a)} Give the genotypes of the parents and of the F$_1$. \textbf{(b)} The F$_1$ plants are self-pollinated. Using a Punnett square, predict the genotypic and phenotypic ratios of the F$_2$. \textbf{(c)} 800 F$_2$ grains are harvested; how many would you expect to be yellow?
\tcblower
\textbf{(a) Parental genotypes.} ``Pure-breeding'' means \cle{homozygous}. The purple parent is $PP$ and the yellow parent, showing the recessive trait, is $pp$.\par
Gametes: the purple parent produces only $P$ gametes; the yellow parent only $p$ gametes.\par
F$_1$ zygotes: all $Pp$ (heterozygous), phenotype \textbf{purple} because $P$ is dominant.\par
\textbf{(b) F$_1$ self-pollination: $Pp \times Pp$.} Each F$_1$ plant produces two gamete types, $P$ and $p$, in equal proportions.
\begin{center}
\begin{tabular}{|c|c|c|}
\hline
\rowcolor{popblueL} Gametes & $P$ & $p$ \\
\hline
$P$ & $PP$ (purple) & $Pp$ (purple) \\
\hline
$p$ & $Pp$ (purple) & $pp$ (yellow) \\
\hline
\end{tabular}
\end{center}
Genotypic ratio: $1\,PP : 2\,Pp : 1\,pp$, i.e. $1:2:1$.\par
Phenotypic ratio: $3$ purple : $1$ yellow, i.e. $3:1$.\par
\textbf{(c) Expected yellow grains.} Yellow grains are $pp$, with probability $\tfrac{1}{4}$:
\[ \frac{1}{4} \times 800 = 200 \text{ yellow grains.} \]
\textbf{Examiner's remark:} always distinguish \emph{genotypic} ratio ($1:2:1$) from \emph{phenotypic} ratio ($3:1$); confusing them is the commonest error in this type of question.
\end{exoresolu}

\courssub{Skill 2: The Test Cross — Revealing an Unknown Genotype}

\begin{methode}[Identifying the genotype of a dominant-phenotype individual]
\begin{enumerate}
\item An organism showing the \cle{dominant} trait has genotype $AA$ or $Aa$. To decide, cross it with a \textbf{homozygous recessive} individual ($aa$): this is a \cle{test cross}.
\item Write the two possibilities separately: $AA \times aa$ and $Aa \times aa$.
\item Work out the offspring of each cross.
\item \textbf{Interpret}: if \emph{all} offspring show the dominant trait, the tested parent was $AA$; if the offspring are $1$ dominant : $1$ recessive, the tested parent was $Aa$.
\end{enumerate}
\textbf{Key idea:} the recessive parent contributes only recessive alleles, so the offspring phenotypes directly reveal the gametes produced by the tested parent.
\end{methode}

\begin{exoresolu}[Test cross in goats]
In goats, black coat ($B$) is dominant to white coat ($b$). A farmer in the North West Region owns a black male goat and wants to know whether it is pure-breeding before using it for reproduction. He crosses it with a white female. \textbf{(a)} Explain why a white female is used. \textbf{(b)} The cross produces 6 black kids and 5 white kids. Deduce the genotype of the black male, showing your working.
\tcblower
\textbf{(a)} A white goat shows the recessive phenotype, so its genotype is certainly $bb$. It can only donate a $b$ allele to every offspring. Any white offspring must therefore have received a $b$ allele \emph{from the black male}, revealing that the male carries $b$. This is the principle of the test cross.\par
\textbf{(b)} Hypothesis 1: male is $BB$. Then $BB \times bb$ gives all offspring $Bb$ (all black) — this does not match the results, since white kids appeared.\par
Hypothesis 2: male is $Bb$. Then:
\begin{center}
\begin{tabular}{|c|c|c|}
\hline
\rowcolor{popblueL} Gametes & $B$ & $b$ \\
\hline
$b$ & $Bb$ (black) & $bb$ (white) \\
\hline
$b$ & $Bb$ (black) & $bb$ (white) \\
\hline
\end{tabular}
\end{center}
Expected ratio: $1\,Bb$ (black) : $1\,bb$ (white), i.e. $\tfrac{1}{2}$ black, $\tfrac{1}{2}$ white. The observed 6 black : 5 white is close to the expected $1:1$ (small samples rarely give exactly equal numbers).\par
\textbf{Conclusion:} the black male is \textbf{heterozygous $Bb$} and is not pure-breeding.
\end{exoresolu}

\courssub{Skill 3: Codominance and Incomplete Dominance}

\begin{methode}[Recognising and solving non-complete dominance]
\begin{enumerate}
\item \textbf{Diagnose the pattern} from the statement:
\begin{itemize}
\item \cle{Incomplete dominance}: the heterozygote has an \emph{intermediate} phenotype (e.g. red $\times$ white $\to$ pink).
\item \cle{Codominance}: the heterozygote expresses \emph{both} traits fully and separately (e.g. roan cattle with red and white hairs; blood group AB).
\end{itemize}
\item Use two different capital letters (e.g. $C^R$, $C^W$) since neither allele is recessive.
\item Work the Punnett square exactly as for a monohybrid cross.
\item \textbf{Key signature}: the phenotypic ratio equals the genotypic ratio, $1:2:1$, with \emph{three} distinct phenotypes.
\end{enumerate}
\end{methode}

\begin{exoresolu}[Incomplete dominance in ``four o'clock'' flowers]
In \emph{Mirabilis jalapa} (the four o'clock plant, common in Cameroonian gardens), a plant with red flowers crossed with a plant with white flowers produces only pink-flowered offspring. \textbf{(a)} Explain this result genetically. \textbf{(b)} Two pink-flowered plants are crossed. Predict the phenotypic ratio of the offspring and the probability of obtaining a white-flowered plant among 120 seedlings.
\tcblower
\textbf{(a)} The F$_1$ phenotype (pink) is \emph{intermediate} between the two parents: this is \cle{incomplete dominance}. Neither allele is fully dominant. Let $C^R$ = allele for red, $C^W$ = allele for white.\par
Parents: $C^R C^R$ (red) $\times$ $C^W C^W$ (white). Gametes: $C^R$ and $C^W$. F$_1$: all $C^R C^W$, phenotype pink.\par
\textbf{(b)} Cross: $C^R C^W \times C^R C^W$.
\begin{center}
\begin{tabular}{|c|c|c|}
\hline
\rowcolor{popblueL} Gametes & $C^R$ & $C^W$ \\
\hline
$C^R$ & $C^R C^R$ (red) & $C^R C^W$ (pink) \\
\hline
$C^W$ & $C^R C^W$ (pink) & $C^W C^W$ (white) \\
\hline
\end{tabular}
\end{center}
Phenotypic ratio: $1$ red : $2$ pink : $1$ white (identical to the genotypic ratio $1:2:1$).\par
Probability of white $= \tfrac{1}{4}$; expected number $= \tfrac{1}{4} \times 120 = 30$ white-flowered plants.\par
\textbf{Warning:} do not write $Rr$ here — using a capital/small letter pair would wrongly imply that one allele is dominant.
\end{exoresolu}

\courssub{Skill 4: Sex-Linked Inheritance}

\begin{methode}[Solving a sex-linkage problem]
\begin{enumerate}
\item Check that the gene is carried on the \cle{X chromosome} (stated in the question, or deduced because males and females are affected differently).
\item Write genotypes with the sex chromosomes: female $X^A X^A$, $X^A X^a$ or $X^a X^a$; male $X^A Y$ or $X^a Y$. The Y carries \emph{no allele} of the gene.
\item Remember: a male receives his X \emph{from his mother} and gives his X \emph{to all his daughters}; a male with the recessive allele on his X \emph{expresses} the trait (he is \cle{hemizygous}).
\item Draw the Punnett square with gametes ($X^A$ or $X^a$ from the mother; $X$ or $Y$ from the father).
\item Quote phenotypic ratios \emph{separately for each sex} if the question asks.
\end{enumerate}
\end{methode}

\begin{exoresolu}[Haemophilia in a royal-style pedigree]
Haemophilia is caused by a recessive allele $h$ carried on the X chromosome; $H$ is the normal allele. A carrier woman marries a normal man. \textbf{(a)} Give the genotypes of the couple. \textbf{(b)} Determine, for their children, the probability of a haemophiliac son, and the probability of a carrier daughter. \textbf{(c)} Explain why haemophiliac daughters are extremely rare.
\tcblower
\textbf{(a)} Carrier woman: $X^H X^h$ (she is phenotypically normal but carries $h$). Normal man: $X^H Y$.\par
\textbf{(b)} Gametes: mother $\to X^H$ or $X^h$; father $\to X^H$ or $Y$.
\begin{center}
\begin{tabular}{|c|c|c|}
\hline
\rowcolor{popblueL} & $X^H$ (father) & $Y$ (father) \\
\hline
$X^H$ (mother) & $X^H X^H$ normal daughter & $X^H Y$ normal son \\
\hline
$X^h$ (mother) & $X^H X^h$ carrier daughter & $X^h Y$ haemophiliac son \\
\hline
\end{tabular}
\end{center}
Each box has probability $\tfrac{1}{4}$:\par
P(haemophiliac son, $X^h Y$) $= \tfrac{1}{4}$ of all children (i.e. $\tfrac{1}{2}$ of the sons).\par
P(carrier daughter, $X^H X^h$) $= \tfrac{1}{4}$ of all children (i.e. $\tfrac{1}{2}$ of the daughters).\par
\textbf{(c)} A haemophiliac daughter would need genotype $X^h X^h$: she must receive $X^h$ from \emph{both} parents, i.e. a carrier (or affected) mother \emph{and} a haemophiliac father ($X^h Y$). Since affected males are uncommon and such a couple is rare, $X^h X^h$ daughters are extremely rare. Males need only \emph{one} copy of $h$ to be affected, which is why X-linked recessive diseases are far commoner in males.
\end{exoresolu}

\courssub{Skill 5: Dihybrid Inheritance and Independent Assortment}

\begin{methode}[Solving a dihybrid cross]
\begin{enumerate}
\item Identify the \textbf{two genes} and their dominant/recessive alleles; assign symbols (e.g. $R/r$, $Y/y$).
\item Write parental genotypes (double heterozygote: $RrYy$).
\item \textbf{List the gametes} using the FOIL rule: $RrYy$ gives $RY$, $Ry$, $rY$, $ry$ in equal proportions (law of independent assortment — the genes are on different chromosomes).
\item Draw the $4 \times 4$ Punnett square (16 boxes) \emph{or} use the faster probability method: treat each gene separately ($\tfrac{3}{4}$ dominant, $\tfrac{1}{4}$ recessive) and multiply.
\item Collect phenotypes into the classic $9:3:3:1$ ratio.
\item \textbf{Check}: $9+3+3+1 = 16$ boxes.
\end{enumerate}
\textbf{Shortcut formula:} P(dominant for both) $= \tfrac{3}{4}\times\tfrac{3}{4} = \tfrac{9}{16}$; P(recessive for both) $= \tfrac{1}{4}\times\tfrac{1}{4} = \tfrac{1}{16}$.
\end{methode}

\begin{exoresolu}[Dihybrid cross in the garden pea]
In peas, round seeds ($R$) are dominant to wrinkled ($r$), and yellow cotyledons ($Y$) are dominant to green ($y$). Two plants heterozygous for both genes are crossed. \textbf{(a)} State the gametes produced by each parent. \textbf{(b)} Calculate the proportion of offspring that are (i) round and yellow, (ii) wrinkled and green, (iii) round and green. \textbf{(c)} Out of 1600 seeds, how many are expected to be wrinkled and yellow?
\tcblower
\textbf{(a)} Each parent is $RrYy$. By independent assortment the gametes are: $RY$, $Ry$, $rY$, $ry$, each with probability $\tfrac{1}{4}$.\par
\textbf{(b)} Using the probability method (each gene segregates as $\tfrac{3}{4}$ dominant : $\tfrac{1}{4}$ recessive):
\begin{align*}
\text{(i)}\quad &P(\text{round, yellow}) = \frac{3}{4}\times\frac{3}{4} = \frac{9}{16}\\
\text{(ii)}\quad &P(\text{wrinkled, green}) = \frac{1}{4}\times\frac{1}{4} = \frac{1}{16}\\
\text{(iii)}\quad &P(\text{round, green}) = \frac{3}{4}\times\frac{1}{4} = \frac{3}{16}
\end{align*}
(The fourth class, wrinkled yellow, is also $\tfrac{3}{16}$; the full ratio is $9:3:3:1$.)\par
\textbf{(c)} Expected wrinkled yellow seeds:
\[ \frac{3}{16}\times 1600 = 300 \text{ seeds.} \]
\textbf{Examiner's remark:} if the question demands the full Punnett square, draw all 16 boxes and group them: 9 round-yellow, 3 round-green, 3 wrinkled-yellow, 1 wrinkled-green. Never write gametes like $RR$ or $Rr$ — a gamete carries only \emph{one} allele per gene.
\end{exoresolu}

\courssub{Skill 6: Linkage and Crossover Value}

\begin{methode}[Detecting linkage and computing the crossover value]
\begin{enumerate}
\item Compare the observed offspring numbers with the $1:1:1:1$ (test cross) or $9:3:3:1$ (F$_2$) expectation for \emph{independent} genes.
\item If two classes are \textbf{much more frequent} (the \cle{parental} types) and two are \textbf{rare} (the \cle{recombinant} types), the genes are \cle{linked} on the same chromosome.
\item Identify recombinants: the classes with the \emph{smallest} numbers.
\item Apply the formula:
\[ \text{Crossover value (COV)} = \frac{\text{number of recombinant offspring}}{\text{total offspring}} \times 100\ \% \]
\item The COV estimates the \textbf{distance between the genes} on the chromosome: $1\ \%$ recombination $= 1$ map unit (centimorgan, cM).
\end{enumerate}
\textbf{Reflex:} in a test cross of linked genes, parental classes are roughly equal and large; recombinant classes are roughly equal and small.
\end{methode}

\begin{exoresolu}[Linkage in maize]
In maize, coloured grains ($C$) are dominant to colourless ($c$), and full grains ($S$) are dominant to shrunken ($s$). A plant heterozygous for both genes is test-crossed with a colourless, shrunken plant. The offspring are: coloured full 4032, colourless shrunken 3968, coloured shrunken 149, colourless full 152. \textbf{(a)} Show that the genes are linked. \textbf{(b)} Identify the parental and recombinant classes. \textbf{(c)} Calculate the crossover value and state the distance between the two genes.
\tcblower
\textbf{(a)} If the genes assorted independently, a test cross ($CcSs \times ccss$) would give four classes in a $1:1:1:1$ ratio, i.e. about $2075$ in each class out of $8301$. Instead, two classes are very large ($\approx 4000$) and two very small ($\approx 150$). This strong departure from $1:1:1:1$ shows the genes are \textbf{linked} on the same chromosome.\par
\textbf{(b)} Parental classes (most frequent, produced without crossing over): \textbf{coloured full} (4032) and \textbf{colourless shrunken} (3968). The heterozygous parent therefore had the alleles in \emph{coupling}: $CS/cs$.\par
Recombinant classes (produced by crossing over between the two loci): \textbf{coloured shrunken} (149) and \textbf{colourless full} (152).\par
\textbf{(c)} Total offspring $= 4032 + 3968 + 149 + 152 = 8301$. Recombinants $= 149 + 152 = 301$.
\[ \text{COV} = \frac{301}{8301}\times 100 \approx 3.6\ \% \]
The two genes are therefore about \textbf{3.6 map units (cM) apart} on the chromosome.\par
\textbf{Examiner's remark:} always divide by the \emph{total} number of offspring, not by the parental total; and quote the answer as a percentage.
\end{exoresolu}

\courssub{Skill 7: Applying the Hardy-Weinberg Equation}

\begin{methode}[Calculating allele and genotype frequencies]
\begin{enumerate}
\item Check the conditions: a \textbf{large, randomly mating} population with no selection, no mutation and no migration — then $p + q = 1$ and $p^2 + 2pq + q^2 = 1$.
\item \textbf{Anchor on the recessive phenotype}: its frequency is $q^2$ (only $aa$ individuals show it). This is usually the only genotype frequency directly observable.
\item Take the square root to get $q$: $q = \sqrt{q^2}$; then $p = 1 - q$.
\item Compute the required genotype frequencies: $p^2$ (homozygous dominant), $2pq$ (heterozygous carriers), $q^2$ (homozygous recessive).
\item Multiply by the population size if numbers of individuals are requested.
\end{enumerate}
\textbf{Key formulas:} $p + q = 1$; \quad $p^2 + 2pq + q^2 = 1$.
\end{methode}

\begin{exoresolu}[Sickle cell trait in a Cameroonian population]
Sickle cell anaemia is caused by a recessive allele $H^S$ of the haemoglobin gene; $H^A$ is the normal allele. In a large, randomly mating population of 10 000 newborns in a Cameroonian town, 160 babies are born with sickle cell anaemia (genotype $H^S H^S$). Assume Hardy-Weinberg equilibrium. \textbf{(a)} Calculate the frequency of the $H^S$ allele. \textbf{(b)} Calculate the frequency of carriers ($H^A H^S$) and the expected number of carriers among the 10 000 newborns. \textbf{(c)} State one biological reason why the $H^S$ allele remains frequent in malaria-endemic regions.
\tcblower
\textbf{(a)} Affected babies are $H^S H^S$, so
\[ q^2 = \frac{160}{10000} = 0.016 \quad\Rightarrow\quad q = \sqrt{0.016} \approx 0.126. \]
The frequency of the $H^S$ allele is about $0.126$ (i.e. $12.6\ \%$).\par
\textbf{(b)} $p = 1 - q = 1 - 0.126 = 0.874$.\par
Carrier frequency:
\[ 2pq = 2 \times 0.874 \times 0.126 \approx 0.220. \]
Expected number of carriers:
\[ 0.220 \times 10000 \approx 2203 \text{ carrier babies.} \]
(So roughly 22\ \% of newborns carry the sickle cell trait without being ill.)\par
\textbf{(c)} In regions where malaria is endemic, heterozygotes ($H^A H^S$) are \textbf{more resistant to severe malaria} than $H^A H^A$ individuals. This \emph{heterozygote advantage} (balanced polymorphism) means natural selection maintains the $H^S$ allele at a relatively high frequency, despite the elimination of many $H^S H^S$ individuals.\par
\textbf{Examiner's remark:} never equate the frequency of the \emph{recessive phenotype} with the frequency of the \emph{recessive allele} — you must take the square root first.
\end{exoresolu}

\courssub{Skill 8: Reading a Pedigree}

\begin{methode}[Deducing the mode of inheritance from a pedigree]
\begin{enumerate}
\item Draw/inspect the symbols: squares = males, circles = females; shaded = affected.
\item \textbf{Dominant or recessive?} If two \emph{unaffected} parents produce an \emph{affected} child, the trait must be \textbf{recessive} (the parents are carriers). If every affected child has at least one affected parent, suspect dominance.
\item \textbf{Autosomal or sex-linked?} If mainly males are affected, and affected fathers never transmit to sons but carrier mothers transmit to half their sons, suspect \textbf{X-linked recessive}.
\item Assign genotypes starting from \emph{certain} cases: affected recessive individuals ($aa$ or $X^a Y$), then deduce carriers among their relatives.
\item Compute risks (e.g. probability that a future child is affected) using the genotypes deduced.
\end{enumerate}
\end{methode}

\begin{exoresolu}[Pedigree of albinism]
Albinism in humans is controlled by a single autosomal gene. In a family, two parents with normal skin pigmentation have four children: three normal and one albino girl. \textbf{(a)} Determine the mode of inheritance of albinism and the genotypes of the parents (use $A$ and $a$). \textbf{(b)} What is the probability that their next child will be albino? \textbf{(c)} One normal daughter later marries a normal man whose father was albino. What is the probability that their first child is albino?
\tcblower
\textbf{(a)} Unaffected parents producing an affected child proves the allele is \textbf{recessive} and that both parents are \textbf{heterozygous carriers} $Aa \times Aa$. The albino girl is $aa$. (The fact that a girl is affected, with an unaffected father, also excludes X-linked recessive inheritance — consistent with an autosomal gene.)\par
\textbf{(b)} From $Aa \times Aa$: P($aa$) $= \tfrac{1}{4}$ for \emph{each} pregnancy, whatever the previous children were. Probability $= \tfrac{1}{4}$ ($25\ \%$).\par
\textbf{(c)} The normal daughter of an $Aa \times Aa$ cross is not $aa$; among the \emph{normal} offspring the genotypes are $\tfrac{1}{3}\,AA$ and $\tfrac{2}{3}\,Aa$. So P(she is a carrier) $= \tfrac{2}{3}$.\par
Her husband is normal but his father was albino ($aa$): he must have received $a$ from his father and $A$ from his mother, so he is \emph{certainly} $Aa$.\par
A child can be albino only if the mother is a carrier:
\[ P(\text{albino child}) = \underbrace{\frac{2}{3}}_{\text{mother } Aa} \times \underbrace{\frac{1}{4}}_{Aa \times Aa \to aa} = \frac{2}{12} = \frac{1}{6} \approx 16.7\ \%. \]
\textbf{Examiner's remark:} the $\tfrac{2}{3}$ (not $\tfrac{1}{2}$) carrier probability for a \emph{known-normal} child of two carriers is a classic GCE trap — the $aa$ possibility is excluded by the normal phenotype, so only three genotypes remain.
\end{exoresolu}

\newpage

\coursec{Exercises}

\courssub{I check my knowledge}

\begin{exercice}[Multiple choice questions]
For each question, choose the single correct answer.
\begin{enumerate}
\item The physical appearance of an organism resulting from its genotype and the environment is called its:
\begin{enumerate}
\item genotype \quad b) phenotype \quad c) allele \quad d) locus
\end{enumerate}
\item In a monohybrid cross between two heterozygotes showing complete dominance, the expected F$_2$ phenotypic ratio is:
\begin{enumerate}
\item 1 : 1 \quad b) 1 : 2 : 1 \quad c) 3 : 1 \quad d) 9 : 3 : 3 : 1
\end{enumerate}
\item A woman is a carrier of the allele for red-green colour blindness. Her genotype is:
\begin{enumerate}
\item X$^N$X$^N$ \quad b) X$^n$X$^n$ \quad c) X$^N$X$^n$ \quad d) X$^n$Y
\end{enumerate}
\item The ABO blood group system in humans is an example of:
\begin{enumerate}
\item complete dominance only \quad b) multiple alleles and codominance \quad c) sex linkage \quad d) epistasis
\end{enumerate}
\end{enumerate}
\end{exercice}

\begin{exercice}[True or false]
State whether each statement is true or false, and justify your answer in one or two sentences.
\begin{enumerate}
\item Mendel's law of independent assortment applies to genes located on the same chromosome.
\item In incomplete dominance, the heterozygote shows a phenotype intermediate between the two homozygotes.
\item A father passes his X chromosome to all of his sons.
\item In a population in Hardy-Weinberg equilibrium for a gene with two alleles, $p + q = 1$.
\end{enumerate}
\end{exercice}

\begin{exercice}[Key definitions]
Define each of the following terms in one or two clear sentences, giving one example where possible:
\begin{enumerate}
\item allele;
\item homozygous and heterozygous;
\item test cross;
\item linkage;
\item epistasis;
\item polymorphism.
\end{enumerate}
\end{exercice}

\begin{exercice}[Direct calculations]
\begin{enumerate}
\item In a population in Hardy-Weinberg equilibrium, the frequency of a recessive allele $a$ is $q = 0.3$. Calculate $p$, the frequency of heterozygotes, and the frequency of the homozygous recessive genotype.
\item A cross between two heterozygous tall pea plants (Tt $\times$ Tt) produces 200 offspring. How many would you expect to be dwarf?
\item A couple who are both carriers of sickle cell anaemia (Hb$^A$Hb$^S$) have four children. What is the probability that any one child has sickle cell anaemia?
\end{enumerate}
\end{exercice}

\courssub{I practise}

\begin{exercice}[Incomplete dominance in flowers]
A cross between a red-flowered plant and a white-flowered plant of the same species produces an F$_1$ generation in which all the flowers are pink. When the F$_1$ plants are self-pollinated, the F$_2$ consists of 62 red, 131 pink and 59 white-flowered plants.
\begin{enumerate}
\item Identify the pattern of inheritance and explain your reasoning.
\item Using suitable symbols, draw a genetic diagram to explain the F$_2$ result.
\item State the expected F$_2$ phenotypic ratio and compare it with the observed numbers.
\end{enumerate}
\end{exercice}

\begin{exercice}[Sex-linked inheritance of colour blindness]
Red-green colour blindness is caused by a recessive allele carried on the X chromosome. A woman with normal vision, whose father was colour-blind, marries a colour-blind man.
\begin{enumerate}
\item State the genotypes of the woman, her father and her husband.
\item Using a full genetic diagram, determine the probability that:
\begin{enumerate}
\item a son will be colour-blind;
\item a daughter will be colour-blind.
\end{enumerate}
\item Explain why colour blindness is more common in males than in females.
\end{enumerate}
\end{exercice}

\begin{exercice}[An ABO paternity case]
A mother of blood group A has a child of blood group O. A man of blood group AB is alleged to be the father.
\begin{enumerate}
\item State all possible genotypes of the mother and the genotype of the child.
\item Using a genetic diagram, show whether the alleged father could be the biological father of this child. Justify your conclusion.
\item State the possible blood group(s) and genotype(s) of the real biological father.
\end{enumerate}
\end{exercice}

\begin{exercice}[Hardy-Weinberg in a Cameroonian town]
In a town of 10\,000 people, 36 individuals suffer from a recessive genetic disorder. The population is assumed to be in Hardy-Weinberg equilibrium for this gene.
\begin{enumerate}
\item State three assumptions of the Hardy-Weinberg equilibrium.
\item Calculate the frequency of the recessive allele ($q$) and of the dominant allele ($p$).
\item Calculate the number of carriers (heterozygotes) in the town.
\item Explain why the majority of recessive alleles in the population are found in phenotypically normal individuals.
\end{enumerate}
\end{exercice}

\courssub{I prepare for the GCE}

\begin{exercice}[Linkage and crossover value in guinea pigs (20 marks)]
In guinea pigs, black coat (B) is dominant to brown coat (b), and rough fur (R) is dominant to smooth fur (r). A guinea pig heterozygous for both genes was test-crossed with a double recessive individual. The offspring obtained were:
\begin{center}
\begin{tabular}{|l|c|}
\hline
\rowcolor{popblueL} \textbf{Phenotype of offspring} & \textbf{Number} \\
\hline
Black, rough & 410 \\
\hline
Brown, smooth & 398 \\
\hline
Black, smooth & 102 \\
\hline
Brown, rough & 90 \\
\hline
\end{tabular}
\end{center}
\begin{enumerate}
\item Define the terms \emph{test cross} and \emph{linkage}. \textbf{(4 marks)}
\item State the results that would have been expected if the two genes assorted independently. \textbf{(3 marks)}
\item Using the data, explain why the two genes must be linked, and state which offspring classes are the recombinants. \textbf{(5 marks)}
\item Calculate the crossover value (recombination frequency) between the two genes, showing your working. \textbf{(4 marks)}
\item Draw a diagram to show how crossing over between the two loci during meiosis produces the recombinant classes. \textbf{(4 marks)}
\end{enumerate}
\end{exercice}

\begin{exercice}[Epistasis in summer squash (15 marks)]
In summer squash, fruit colour is controlled by two genes. The dominant allele W produces white fruit and is epistatic to the second gene; in plants with the genotype ww, the allele Y produces yellow fruit while yy produces green fruit.
\begin{enumerate}
\item Define the term \emph{epistasis} and distinguish it from dominance. \textbf{(4 marks)}
\item State the phenotypes of plants with the following genotypes: WWYy; wwYy; wwyy. \textbf{(3 marks)}
\item Using a genetic diagram (Punnett square), predict the F$_2$ phenotypic ratio obtained from the cross WwYy $\times$ WwYy. \textbf{(6 marks)}
\item Name the type of epistasis illustrated by this cross. \textbf{(2 marks)}
\end{enumerate}
\end{exercice}

\begin{exercice}[Essay: variation and polymorphism (20 marks, GCE Paper 2 style)]
``Describe the causes and types of variation in a named population, and explain how polymorphism can be maintained by natural selection, using the sickle cell trait in Cameroon as an example.''

Your essay should include:
\begin{enumerate}
\item a clear distinction between continuous and discontinuous variation, with one example of each;
\item an account of the genetic and environmental causes of variation (mutation, meiosis, random fertilisation, environmental effects);
\item an explanation of balanced polymorphism, referring to the advantage of the heterozygote Hb$^A$Hb$^S$ in regions where malaria is endemic;
\item a brief, correctly labelled genetic diagram showing the expected offspring of two carrier parents.
\end{enumerate}
\end{exercice}

\newpage

\coursec{Solutions to the exercises}

\courssub{I check my knowledge}

\begin{corrige}[Exercise 1 — Multiple choice questions]
\begin{enumerate}
\item \textbf{Answer: b) phenotype.} The \cle{phenotype} is the observable characteristic of an organism (e.g.\ tall stem, blood group A), resulting from the interaction of the genotype with the environment. The genotype is the genetic constitution itself, an allele is one version of a gene, and a locus is the position of a gene on a chromosome.
\item \textbf{Answer: c) 3 : 1.} For Tt $\times$ Tt with complete dominance, the genotypes are 1 TT : 2 Tt : 1 tt, and since TT and Tt have the same phenotype, the phenotypic ratio is 3 dominant : 1 recessive.
\item \textbf{Answer: c) X$^N$X$^n$.} A \cle{carrier} female is heterozygous: she carries one normal allele (X$^N$) and one recessive colour-blindness allele (X$^n$) but has normal vision because the normal allele is dominant. X$^n$X$^n$ would be a colour-blind woman and X$^n$Y a colour-blind man.
\item \textbf{Answer: b) multiple alleles and codominance.} The ABO system involves three alleles (I$^A$, I$^B$, I$^o$) — hence \cle{multiple alleles} — and I$^A$ and I$^B$ are \cle{codominant} (both expressed in group AB), while both are dominant to I$^o$.
\end{enumerate}
\end{corrige}

\begin{corrige}[Exercise 2 — True or false]
\begin{enumerate}
\item \textbf{False.} The law of independent assortment applies to genes located on \emph{different} chromosomes (or far apart on the same chromosome). Genes on the same chromosome are \cle{linked} and tend to be inherited together, so they do not assort independently.
\item \textbf{True.} In \cle{incomplete dominance}, neither allele is fully dominant, so the heterozygote (e.g.\ Rr) shows an intermediate phenotype (e.g.\ pink flowers from red $\times$ white parents).
\item \textbf{False.} A father passes his \textbf{Y} chromosome to his sons (this is what makes them male); he passes his X chromosome to all of his \emph{daughters}.
\item \textbf{True.} For a gene with only two alleles, the sum of their frequencies must equal 1: $p + q = 1$, since every allele in the population is either one or the other.
\end{enumerate}
\end{corrige}

\begin{corrige}[Exercise 3 — Key definitions]
\begin{enumerate}
\item \textbf{Allele:} one of the alternative forms of a gene, occupying the same locus on homologous chromosomes. Example: the alleles T (tall) and t (dwarf) of the height gene in peas.
\item \textbf{Homozygous / heterozygous:} an organism is \cle{homozygous} for a gene when it carries two identical alleles (e.g.\ TT or tt); it is \cle{heterozygous} when it carries two different alleles (e.g.\ Tt).
\item \textbf{Test cross:} a cross between an individual of unknown genotype showing the dominant phenotype and a homozygous recessive individual, used to reveal whether the unknown individual is homozygous or heterozygous.
\item \textbf{Linkage:} the tendency of genes located on the same chromosome to be inherited together because they are physically connected; linkage reduces the number of gamete types produced.
\item \textbf{Epistasis:} an interaction between genes at \emph{different} loci in which one gene masks or suppresses the expression of another gene. Example: the W allele in summer squash masks the Y/y colour gene.
\item \textbf{Polymorphism:} the occurrence, within a population, of two or more distinct forms (morphs) at frequencies too high to be maintained by mutation alone. Example: the ABO blood groups, or the Hb$^A$/Hb$^S$ haemoglobin polymorphism in Cameroon.
\end{enumerate}
\end{corrige}

\begin{corrige}[Exercise 4 — Direct calculations]
\begin{enumerate}
\item Given $q = 0.3$:
\begin{align*}
p &= 1 - q = 1 - 0.3 = 0.7 \\
\text{Frequency of heterozygotes} &= 2pq = 2 \times 0.7 \times 0.3 = 0.42 \\
\text{Frequency of } aa &= q^2 = 0.3^2 = 0.09
\end{align*}
So $p = 0.7$; 42\,\% of the population are heterozygous carriers and 9\,\% are homozygous recessive.
\item Tt $\times$ Tt gives $\tfrac{1}{4}$ tt (dwarf). Expected number of dwarf plants:
\[ \frac{1}{4} \times 200 = 50 \text{ dwarf plants.} \]
\item Each child of Hb$^A$Hb$^S$ $\times$ Hb$^A$Hb$^S$ has a probability $\tfrac{1}{4}$ (25\,\%) of being Hb$^S$Hb$^S$ (sickle cell anaemia). \textbf{Note:} this probability is the same for each of the four children, because each fertilisation is an independent event — the fact that the couple has four children does not change the probability for any one child.
\end{enumerate}
\end{corrige}

\courssub{I practise}

\begin{corrige}[Exercise 5 — Incomplete dominance in flowers]
\begin{enumerate}
\item \textbf{Pattern of inheritance: incomplete dominance.} Reasoning: the F$_1$ are all \emph{pink}, a phenotype \textbf{intermediate} between the two parents (red and white). Neither allele is completely dominant over the other, so the heterozygote has its own distinct phenotype. This is confirmed by the F$_2$, where three phenotypes appear.
\item Let R = allele for red, W = allele for white (no dominance, so both are written as capital letters or as R/R$'$).

\begin{center}
\begin{tabular}{|c|c|c|}
\hline
\rowcolor{popblueL} Gametes & R & W \\
\hline
R & RR (red) & RW (pink) \\
\hline
W & RW (pink) & WW (white) \\
\hline
\end{tabular}
\end{center}

F$_1$ self-cross: RW $\times$ RW $\rightarrow$ 1 RR (red) : 2 RW (pink) : 1 WW (white).
\item \textbf{Expected F$_2$ ratio: 1 red : 2 pink : 1 white.} With $62 + 131 + 59 = 252$ plants, the expected numbers are:
\[ \text{red} = \tfrac{1}{4}\times 252 = 63; \quad \text{pink} = \tfrac{1}{2}\times 252 = 126; \quad \text{white} = \tfrac{1}{4}\times 252 = 63. \]
The observed numbers (62, 131, 59) are very close to the expected values (63, 126, 63); the small differences are due to chance. The data therefore fit the 1 : 2 : 1 ratio characteristic of incomplete dominance.
\end{enumerate}
\end{corrige}

\begin{corrige}[Exercise 6 — Sex-linked inheritance of colour blindness]
Let X$^N$ = X chromosome carrying the normal allele; X$^n$ = X chromosome carrying the colour-blindness allele.
\begin{enumerate}
\item \textbf{Genotypes:}
\begin{itemize}
\item The woman's father was colour-blind: X$^n$Y. He must give his X$^n$ to his daughter.
\item The woman has normal vision but received X$^n$ from her father: she is a \textbf{carrier}, X$^N$X$^n$.
\item The husband is colour-blind: X$^n$Y.
\end{itemize}
\item Genetic diagram: X$^N$X$^n$ $\times$ X$^n$Y.

\begin{center}
\begin{tabular}{|c|c|c|}
\hline
\rowcolor{popblueL} Gametes & X$^n$ (father) & Y (father) \\
\hline
X$^N$ (mother) & X$^N$X$^n$ carrier daughter, normal & X$^N$Y normal son \\
\hline
X$^n$ (mother) & X$^n$X$^n$ colour-blind daughter & X$^n$Y colour-blind son \\
\hline
\end{tabular}
\end{center}

\begin{enumerate}
\item Probability that a son is colour-blind: among sons (X$^N$Y, X$^n$Y), half are colour-blind $\Rightarrow$ $\tfrac{1}{2}$ (50\,\%).
\item Probability that a daughter is colour-blind: among daughters (X$^N$X$^n$, X$^n$X$^n$), half are colour-blind $\Rightarrow$ $\tfrac{1}{2}$ (50\,\%).
\end{enumerate}
\item \textbf{Why males are more often affected:} males have only \textbf{one X chromosome} (they are \cle{hemizygous}). A single recessive allele on that X is therefore always expressed, since there is no second allele to mask it. Females have two X chromosomes, so they are affected only if they inherit the recessive allele from \emph{both} parents (X$^n$X$^n$), which is much rarer; a heterozygous woman is a normal-vision carrier.
\end{enumerate}
\end{corrige}

\begin{corrige}[Exercise 7 — An ABO paternity case]
\begin{enumerate}
\item \textbf{Mother (group A):} possible genotypes I$^A$I$^A$ or I$^A$I$^o$. \textbf{Child (group O):} genotype must be I$^o$I$^o$ (the only genotype giving group O). Since the child received one I$^o$ from the mother, the mother must be I$^A$I$^o$.
\item \textbf{Alleged father (group AB):} genotype I$^A$I$^B$. Cross: I$^A$I$^o$ (mother) $\times$ I$^A$I$^B$ (alleged father).

\begin{center}
\begin{tabular}{|c|c|c|}
\hline
\rowcolor{popblueL} Gametes & I$^A$ & I$^B$ \\
\hline
I$^A$ & I$^A$I$^A$ (A) & I$^A$I$^B$ (AB) \\
\hline
I$^o$ & I$^A$I$^o$ (A) & I$^B$I$^o$ (B) \\
\hline
\end{tabular}
\end{center}

Possible children: groups A, AB or B — but \textbf{never group O}, because the AB man has no I$^o$ allele to transmit. \textbf{Conclusion:} the man of blood group AB \textbf{cannot} be the biological father of this child.
\item The child is I$^o$I$^o$, so the real father must carry at least one I$^o$ allele. Possible genotypes and groups:
\begin{itemize}
\item I$^o$I$^o$ — blood group O;
\item I$^A$I$^o$ — blood group A;
\item I$^B$I$^o$ — blood group B.
\end{itemize}
The real father can therefore be of group O, A or B (but not AB).
\end{enumerate}
\end{corrige}

\begin{corrige}[Exercise 8 — Hardy-Weinberg in a Cameroonian town]
\begin{enumerate}
\item \textbf{Assumptions of the Hardy-Weinberg equilibrium} (any three):
\begin{itemize}
\item the population is very large (no genetic drift);
\item mating is random with respect to the gene;
\item there is no mutation, no migration (no gene flow) and no natural selection acting on the gene.
\end{itemize}
\item The affected individuals are homozygous recessive:
\[ q^2 = \frac{36}{10\,000} = 0.0036 \quad\Rightarrow\quad q = \sqrt{0.0036} = 0.06 \]
\[ p = 1 - q = 1 - 0.06 = 0.94 \]
\item Frequency of carriers: $2pq = 2 \times 0.94 \times 0.06 = 0.1128$.
\[ \text{Number of carriers} = 0.1128 \times 10\,000 = 1\,128 \text{ people.} \]
\item The recessive allele is \emph{hidden} in heterozygotes: carriers (Aa) are phenotypically normal because the dominant allele masks the recessive one. Here the carrier frequency ($2pq = 0.1128$) is about 31 times greater than the frequency of affected individuals ($q^2 = 0.0036$). Recessive alleles are therefore far more numerous in healthy heterozygotes than in affected homozygotes, which is why recessive disorders can persist in a population even when selection acts against affected individuals.
\end{enumerate}
\end{corrige}

\courssub{I prepare for the GCE}

\begin{corrige}[Exercise 9 — Linkage and crossover value in guinea pigs (20 marks)]
\begin{enumerate}
\item \textbf{Test cross (2 marks):} a cross between an individual showing the dominant phenotype and a homozygous recessive individual, used to determine the genotype of the dominant individual (or the gametes it produces). \textbf{Linkage (2 marks):} the tendency of genes situated on the same chromosome to be inherited together, because they are physically attached on the same DNA molecule.
\item \textbf{Independent assortment (3 marks):} a test cross BbRr $\times$ bbrr would give four phenotypic classes in \textbf{equal proportions}: 1 black rough : 1 black smooth : 1 brown rough : 1 brown smooth, i.e.\ 25\,\% each (about 250 of each class out of 1\,000).
\item \textbf{Evidence of linkage (5 marks):} total offspring $= 410 + 398 + 102 + 90 = 1\,000$. The four classes are \emph{not} equal: two classes are very frequent (black rough 410; brown smooth 398) and two are rare (black smooth 102; brown rough 90). The frequent classes have the \textbf{same allele combinations as the parents} (parental classes), showing that B is linked with R and b with r on the same chromosome. The rare classes are the \textbf{recombinants}: \emph{black, smooth} (102) and \emph{brown, rough} (90), produced by crossing over.
\item \textbf{Crossover value (4 marks):}
\[ \text{COV} = \frac{\text{number of recombinant offspring}}{\text{total offspring}} \times 100 = \frac{102 + 90}{1\,000} \times 100 = 19.2\,\% \]
The two loci are about \textbf{19.2 map units (cM)} apart.
\item \textbf{Diagram of crossing over (4 marks):} the heterozygote has the genes in \emph{coupling} (cis): BR / br. During prophase I, non-sister chromatids of the homologous chromosomes exchange segments between the two loci:

\begin{popfigure}
\begin{center}
\begin{tikzpicture}[scale=1.0, every node/.style={font=\small}]
% before
\draw[very thick, popblue] (0,2.4) -- (4,2.4);
\draw[very thick, poporange] (0,1.6) -- (4,1.6);
\fill[popblue] (1,2.4) circle (2pt) node[above, popink]{B};
\fill[popblue] (3,2.4) circle (2pt) node[above, popink]{R};
\fill[poporange] (1,1.6) circle (2pt) node[below, popink]{b};
\fill[poporange] (3,1.6) circle (2pt) node[below, popink]{r};
\node[popink] at (-0.6,2) {before};
% crossover mark
\draw[popdark, dashed, thick] (2,2.7) -- (2,1.3) node[midway, right, popdark]{chiasma};
% after
\draw[very thick, popblue] (6,2.4) -- (8,2.4);
\draw[very thick, poporange] (8,2.4) -- (10,2.4);
\draw[very thick, poporange] (6,1.6) -- (8,1.6);
\draw[very thick, popblue] (8,1.6) -- (10,1.6);
\fill[popblue] (7,2.4) circle (2pt) node[above, popink]{B};
\fill[poporange] (9,2.4) circle (2pt) node[above, popink]{r};
\fill[poporange] (7,1.6) circle (2pt) node[below, popink]{b};
\fill[popblue] (9,1.6) circle (2pt) node[below, popink]{R};
\node[popink] at (5.4,2) {after};
\end{tikzpicture}
\end{center}
\legende{Crossing over between the B and R loci produces recombinant chromatids Br and bR; the parental chromatids BR and br remain unchanged.}
\end{popfigure}

Four gamete types result: BR and br (parental, frequent) and Br and bR (recombinant, rare). Fertilised by the gamete br from the double recessive parent, these give respectively black rough, brown smooth (parental classes) and black smooth, brown rough (recombinant classes).
\end{enumerate}
\textbf{Examiner expectations:} clear definitions; correct 1:1:1:1 expectation; explicit comparison of observed vs expected; correct recombinant classes; COV formula with working and \% unit; diagram showing homologous chromosomes, chiasma and the four chromatids labelled.
\end{corrige}

\begin{corrige}[Exercise 10 — Epistasis in summer squash (15 marks)]
\begin{enumerate}
\item \textbf{Epistasis (2 marks):} an interaction between genes at different loci in which one gene (the epistatic gene) masks the expression of another gene (the hypostatic gene). \textbf{Distinction from dominance (2 marks):} dominance is an interaction between \emph{alleles of the same gene} at the same locus, whereas epistasis is an interaction between \emph{different genes} at different loci.
\item \textbf{Phenotypes (3 marks):} WWYy = \textbf{white} (W is epistatic); wwYy = \textbf{yellow} (no W, Y present); wwyy = \textbf{green}.
\item \textbf{F$_2$ from WwYy $\times$ WwYy (6 marks):} each parent produces gametes WY, Wy, wY, wy.

\begin{center}
\begin{tabular}{|c|c|c|c|c|}
\hline
\rowcolor{popblueL} & WY & Wy & wY & wy \\
\hline
WY & WWYY white & WWYy white & WwYY white & WwYy white \\
\hline
Wy & WWYy white & WWyy white & WwYy white & Wwyy white \\
\hline
wY & WwYY white & WwYy white & wwYY yellow & wwYy yellow \\
\hline
wy & WwYy white & Wwyy white & wwYy yellow & wwyy green \\
\hline
\end{tabular}
\end{center}

Counting: any genotype containing at least one W (12/16) is white; wwY\_ (3/16) is yellow; wwyy (1/16) is green.
\[ \boxed{12 \text{ white} : 3 \text{ yellow} : 1 \text{ green}} \]
\item This is \textbf{dominant epistasis} (2 marks): a single \emph{dominant} allele (W) at the first locus is sufficient to mask the second gene.
\end{enumerate}
\textbf{Examiner expectations:} precise definition contrasting gene--gene vs allele--allele interaction; correct phenotypes with justification; a complete 16-box Punnett square with phenotypes written in each box; the modified ratio 12:3:1 explicitly derived from the 9:3:3:1 ratio by grouping classes.
\end{corrige}

\begin{corrige}[Exercise 11 — Essay: variation and polymorphism (20 marks, indicative mark scheme)]
\textbf{Introduction (2 marks):} variation is the existence of differences between individuals of the same species; it is the raw material of evolution by natural selection.

\textbf{1. Types of variation (4 marks):}
\begin{itemize}
\item \textbf{Discontinuous variation:} individuals fall into distinct classes with no intermediates; controlled by one or few genes, little affected by environment. Example: ABO blood groups in humans (or tongue rolling).
\item \textbf{Continuous variation:} characters show a complete range between two extremes, giving a normal distribution; controlled by many genes (polygenic) and strongly influenced by the environment. Example: height or mass in a school population in Yaoundé.
\end{itemize}

\textbf{2. Causes of variation (6 marks):}
\begin{itemize}
\item \textbf{Mutation (2):} random changes in DNA create new alleles — the ultimate source of all genetic variation (e.g.\ the Hb$^S$ allele arose by a point mutation in the $\beta$-globin gene).
\item \textbf{Meiosis (2):} crossing over in prophase I and independent assortment of chromosomes reshuffle alleles, so each gamete is genetically unique.
\item \textbf{Random fertilisation (1):} any sperm may fuse with any ovum, multiplying combinations.
\item \textbf{Environment (1):} diet, climate, disease and lifestyle modify the phenotype without changing the genotype (e.g.\ body mass, skin tanning).
\end{itemize}

\textbf{3. Balanced polymorphism and sickle cell trait (6 marks):} \cle{polymorphism} is the stable coexistence of two or more forms in a population. In Cameroon, where malaria (caused by \emph{Plasmodium}, transmitted by \emph{Anopheles} mosquitoes) is endemic, three genotypes occur: Hb$^A$Hb$^A$ (normal but fully susceptible to severe malaria), Hb$^S$Hb$^S$ (sickle cell anaemia, often lethal without treatment) and Hb$^A$Hb$^S$ (healthy carriers with \textbf{increased resistance to malaria}). Selection \emph{against} Hb$^S$Hb$^S$ removes Hb$^S$ alleles, but selection \emph{for} the heterozygote (\cle{heterozygote advantage}) maintains Hb$^S$ in the population. The two opposing selective forces hold both alleles at stable frequencies — this is \textbf{balanced polymorphism}. In regions without malaria, the Hb$^S$ allele confers no advantage and is progressively eliminated.

\textbf{4. Genetic diagram (2 marks):} Hb$^A$Hb$^S$ $\times$ Hb$^A$Hb$^S$:

\begin{center}
\begin{tabular}{|c|c|c|}
\hline
\rowcolor{popblueL} Gametes & Hb$^A$ & Hb$^S$ \\
\hline
Hb$^A$ & Hb$^A$Hb$^A$ normal & Hb$^A$Hb$^S$ carrier (resistant) \\
\hline
Hb$^S$ & Hb$^A$Hb$^S$ carrier (resistant) & Hb$^S$Hb$^S$ anaemia \\
\hline
\end{tabular}
\end{center}

Offspring: 1 normal : 2 carriers : 1 with sickle cell anaemia.

\textbf{Conclusion (within marks above):} variation, generated by mutation and sexual reproduction and modulated by the environment, allows natural selection to maintain advantageous polymorphisms such as the sickle cell trait in malarial regions of Cameroon.

\textbf{Examiner expectations:} structured essay with introduction and conclusion; at least one correct example per type of variation; all four causes addressed; the terms \emph{heterozygote advantage} and \emph{balanced polymorphism} used correctly; a labelled genetic diagram with gametes, genotypes and phenotypes.
\end{corrige}

\newpage

\coursec{Summary sheet}

\begin{popfigure}
\begin{tikzpicture}[mindmap, grow cyclic, every node/.style={concept}, root concept/.append style={concept color=popblue, font=\large\bfseries, text=white}, level 1/.append style={level distance=3.4cm, sibling angle=51}, level 2/.append style={level distance=2.1cm, font=\small}]
\node[root concept]{Genetics}
  child[concept color=poporange]{ node {Mendel's Laws}
    child { node {Segregation} }
    child { node {Independent assortment} }
  }
  child[concept color=popgreen]{ node {Dominance}
    child { node {Complete} }
    child { node {Codominance} }
    child { node {Incomplete} }
  }
  child[concept color=poppurple]{ node {Sex \& Linkage}
    child { node {Sex determination} }
    child { node {Sex linkage} }
    child { node {Crossing over} }
  }
  child[concept color=popgold]{ node {Gene Interactions}
    child { node {Epistasis} }
    child { node {Lethal genes} }
    child { node {Multiple alleles} }
  }
  child[concept color=poppink]{ node {Population Genetics}
    child { node {Allele frequencies} }
    child { node {Hardy--Weinberg} }
  }
  child[concept color=popdark]{ node {Variation}
    child { node {Polymorphism} }
    child { node {Pedigrees} }
  };
\end{tikzpicture}
\legende{Mind map of the chapter: Genetics.}
\end{popfigure}

\begin{bilan}
\textbf{Key definitions}
\begin{itemize}
\item \cle{Gene}: a segment of DNA coding for a character; \cle{allele}: an alternative form of a gene; \cle{locus}: position of a gene on a chromosome.
\item \cle{Genotype}: genetic constitution; \cle{phenotype}: observable expression of the genotype (influenced by environment).
\item \cle{Homozygous} (two identical alleles) vs \cle{heterozygous} (two different alleles); \cle{dominant} allele expressed in heterozygote; \cle{recessive} allele expressed only when homozygous.
\item $F_1$, $F_2$: first and second filial generations; \cle{test cross}: cross with homozygous recessive to reveal an unknown genotype.
\item \cle{Codominance}: both alleles fully expressed (e.g.\ AB blood group); \cle{incomplete dominance}: intermediate phenotype (e.g.\ pink snapdragon).
\item \cle{Sex linkage}: gene carried on a sex chromosome (usually X); \cle{linkage}: genes on the same chromosome inherited together; \cle{crossing over}: exchange between non-sister chromatids at prophase I producing recombinants.
\item \cle{Epistasis}: one gene masks the expression of another; \cle{lethal gene}: allele causing death in a particular genotype; \cle{multiple alleles}: more than two alleles of one gene in a population (e.g.\ ABO system).
\item \cle{Polymorphism}: coexistence of two or more distinct phenotypes in a population (e.g.\ sickle cell trait in malaria zones of Cameroon).
\end{itemize}

\textbf{Laws and formulas to know}
\begin{itemize}
\item \cle{Mendel's first law} (segregation): allele pairs separate during gamete formation; monohybrid $F_2$ ratio $3:1$ (phenotype), $1:2:1$ (genotype).
\item \cle{Mendel's second law} (independent assortment): genes on different chromosomes assort independently; dihybrid $F_2$ ratio $9:3:3:1$.
\item Crossover value $= \dfrac{\text{number of recombinant offspring}}{\text{total offspring}} \times 100\ \%$; it estimates map distance between linked genes.
\item \cle{Hardy--Weinberg equation}: $p^2 + 2pq + q^2 = 1$ with $p + q = 1$; valid for a large, randomly mating population with no mutation, migration or selection.
\end{itemize}

\textbf{Methods}
\begin{itemize}
\item Always state the cross in symbols: parents $\rightarrow$ gametes $\rightarrow$ Punnett square $\rightarrow$ offspring genotypes and phenotypes with ratios.
\item For dihybrid problems, check whether genes assort independently ($9:3:3:1$) or are linked (excess of parental phenotypes).
\item For pedigrees: recessive traits skip generations; X-linked recessive traits affect mainly males, transmitted via carrier mothers.
\end{itemize}

\textbf{Common mistakes}
\begin{itemize}
\item Confusing genotype and phenotype ratios; forgetting that a $3:1$ ratio requires complete dominance.
\item Writing both alleles of a pair in one gamete (gametes are haploid!).
\item Assuming all genes assort independently; linked genes deviate from $9:3:3:1$.
\item In Hardy--Weinberg, confusing $2pq$ (heterozygotes) with $q^2$ (recessive homozygotes).
\end{itemize}

\textbf{GCE tips}
\begin{itemize}
\item Draw a clear genetic diagram for every cross: examiners award marks for gametes and Punnett squares.
\item Quote expected ratios precisely and compare them with observed data using words such as ``consistent with''.
\item Use Cameroonian contexts (sickle cell and malaria resistance, ABO blood groups) to illustrate polymorphism and selection.
\end{itemize}
\end{bilan}

\findecours

\end{document}
